\chapter{San Pedro}\label{dept:San Pedro}
\mapaparaguai{2}{San Pedro}
\vfill
\newpage
\secaoDiagnostico{de Diciembre}{\mapaConteudo{-24.6666}{-56.8666}}{25 de Diciembre é uma localidade estratégica por seu acesso facilitado
via Ruta PY03 e sua proximidade com grandes centros como San Estanislao,
sem sacrificar a segurança de uma baixíssima densidade demográfica.
Oferece um equilíbrio excelente entre logística e isolamento rural.

\subsubsection{Por que sim}

\begin{itemize}
\tightlist
\item
  Localização privilegiada para logística rápida para Assunção ou Norte.
\item
  Baixa densidade demográfica e ambiente rural pacato.
\item
  Sem restrições legais para estrangeiros (Lei de Fronteira).
\end{itemize}

\subsubsection{Por que não}

\begin{itemize}
\tightlist
\item
  Dependência de cidades vizinhas para saúde especializada.
\item
  Solo exige investimento em correção para agricultura de alta
  performance.
\end{itemize}

\subsubsection{Recomendação
Técnica}

Ideal para perfis que buscam um refúgio rural com conectividade
asfáltica imediata. O GSS de 6.9 reflete a segurança institucional e a
baixa exposição a riscos naturais. Referencia atual de score: GSS 6.9 em
2026-03-05.

\subsubsection{}

\begin{itemize}
\tightlist
\item
  \begin{enumerate}
  \def\labelenumi{\arabic{enumi})}
  \tightlist
  \item
    Dados censitários validados (INE\footnote{Instituto Nacional de Estadística (Instituto Nacional de Estatística), órgão oficial de dados demográficos e censitários.} 2022).
  \end{enumerate}
\item
  \begin{enumerate}
  \def\labelenumi{\arabic{enumi})}
  \setcounter{enumi}{1}
  \tightlist
  \item
    Relatórios de infraestrutura rodoviária (MOPC 2026).
  \end{enumerate}
\item
  \begin{enumerate}
  \def\labelenumi{\arabic{enumi})}
  \setcounter{enumi}{2}
  \tightlist
  \item
    Mapa de solos e aptidão agrícola (MAG/\allowbreak Productiva).
  \end{enumerate}
\end{itemize}}
\begin{table}[ht!]
\small \begin{tabular}{p{3.0cm}p{0.8cm}p{6.0cm}} \toprule
\textbf{Critério} & \textbf{Nota} & \textbf{Justificativa} \\ \midrule
A - Ameaças Estratégicas & 6.0 & Ponto de passagem rodoviário estratégico, mas não um alvo militar primário \\
B - Risco Social & 7.5 & Baixa densidade populacional rural; risco de agitação social reduzido \\
C - Riscos Naturais & 7.5 & Geologia estável e baixo risco de inundações graves \\
D - Autossuficiência & 7.0 & Base agropecuária sólida; boa disponibilidade hídrica e potencial proteico \\
E - Institucional & 7.0 & Município consolidado; liberdade jurídica para estrangeiros limítrofes \\
\midrule \textbf{GSS FINAL} & \textbf{6.9} & \textbf{Moderadamente Seguro (Recomendado para residências rurais com fácil acesso à capital).} \\ \bottomrule \end{tabular}
\end{table}
\subsubsection{Vulnerabilidades e
Mitigação}\label{vuln-25_de_Diciembre-vulnerabilidades-e-mitigauxe7uxe3o}

\begin{itemize}
\tightlist
\item
  \textbf{Dependência Logística:} Dependência de San Estanislao para
  serviços complexos (Mitigação: estoque de insumos e kit trauma local).
\item
  \textbf{Exposição Rodoviária:} Localização na PY03 pode atrair fluxos
  de refugiados em crises urbanas (Mitigação: residência com recuo
  significativo em relação à rodovia).
\item
  \textbf{Qualidade do Solo:} Necessidade de manejo para agricultura
  intensiva (Mitigação: análise de solo e correção com calcário).
\end{itemize}
\subsubsection{Dossiê de Campo}\label{dossie-25_de_Diciembre-dossiuxea-de-campo}
\begin{description}[style=nextline,leftmargin=0.5cm,font=\bfseries]
\item[Coordenadas]  24°40'S 56°52'W.
\item[Topografia]  Relevo plano a levemente ondulado. Localizado no sul do departamento de San Pedro, em posição de transição. Geologicamente estável, risco sísmico desprezível.
\item[Alvos Estratégicos]  Ponto de passagem e fiscalização na Ruta PY03 (Eixo Norte-Sul vital para o país); Posto de pedágio (Peaje) e entroncamentos locais.
\item[Fallout]  Ventos predominantes do Norte/\allowbreak Nordeste (verão) e Sul (inverno). Baixo risco de fallout global direto.
\item[População]  8.703 habitantes (Censo 2022 Final). População majoritariamente rural (89\%).
\item[Densidade]  \textasciitilde 8,7 hab/\allowbreak km² (Baixa densidade, oferecendo excelente isolamento em áreas de fazendas).
\item[Vias de Acesso]  Eixo principal pela Ruta PY03 (asfaltada e de alta trafegabilidade). Estradas rurais vicinais que exigem veículos robustos em épocas de chuva.
\item[Serviços]  Unidade de Saúde da Família (USF) local para atendimentos básicos. Dependência de San Estanislao (Santaní) a 25 km para saúde de alta complexidade e serviços financeiros avançados.
\item[Custo de Vida]  Baixo (20-40\% inferior a Assunção); economia baseada em serviços rodoviários e agropecuária.
\item[Preço da Terra]  US\$ 1.500-2.500/\allowbreak ha (campos mecanizáveis); US\$ 1.000-1.300/\allowbreak ha (áreas de pastagem natural).
\item[Hidrologia]  Baixo risco de inundações catastróficas. Presença de arroios locais que podem causar interrupções pontuais em caminhos vicinais de terra.
\item[Clima]  Subtropical Úmido. Precipitação anual de 1.400-1.600 mm, bem distribuída.
\item[Energia]  Rede nacional (ANDE\footnote{Administración Nacional de Electricidade (Administração Nacional de Eletricidade), companhia estatal de energia.}/\allowbreak Itaipu). Estabilidade moderada por estar próximo a eixos principais.
\item[Água]  Poços artesianos e lençol freático de fácil acesso; Juntas de Saneamento operais.
\item[Qualidade do Solo]  Alfisois e Ultisois (textura franco-arenosa); boa aptidão para pastagens e silvicultura (eucalipto); exige correção para grãos.
\item[Resources Locais]  Produção de gado de corte, milho e mandioca. Elevado potencial para autossuficiência alimentar básica em nível de propriedade.
\item[Segurança]  Zona rural com baixos índices de criminalidade urbana violenta. Patrulhamento da FTC/\allowbreak Polícia Nacional na Ruta PY03 garante presença estatal.
\item[Leis Local: Município consolidado. Livre de Restrição de Fronteira]  Fora da faixa de 50km da Lei 2532/\allowbreak 05, permitindo plena titularidade para estrangeiros brasileiros.
\end{description}
\paragraph{Solo (SoilGrids 2.0, média ponderada 0--30
cm)}

\begin{longtable}[]{@{}ll@{}}
\toprule\noalign{}
Parâmetro & Valor \\
\midrule\noalign{}
\endhead
\bottomrule\noalign{}
\endlastfoot
pH (H$_2$O) & 5.4 \\
Carbono orgânico (SOC) & 15.58 g/\allowbreak kg \\
Argila & 25.69 \% \\
Areia & 49.24 \% \\
Silte (calc.) & 25.1 \% \\
Densidade aparente & 1.33 g/\allowbreak cm³ \\
\textbf{Aptidão agrícola} & \textbf{Média} \\
\end{longtable}

Fonte: ISRIC SoilGrids 2.0 via WCS (acesso em 2026-03-21). Coords:
-24.6667°, -56.8667°. Média ponderada camadas 0-5, 5-15, 15-30 cm.
\newpage
\secaoDiagnostico{Antequera}{\mapaConteudo{-24.1166}{-57.0666}}{Antequera é o portão de saída fluvial do departamento de San Pedro.
Oferece uma infraestrutura logística privilegiada e proximidade com a
capital departamental, garantindo segurança institucional superior. Seu
GSS de 6.9 reflete o equilíbrio entre a riqueza de recursos e o risco
geográfico de inundações.

\subsubsection{Por que sim}

\begin{itemize}
\tightlist
\item
  Nó logístico estratégico com portos e silos privados.
\item
  Baixa densidade populacional e ambiente seguro.
\item
  Proximidade com serviços completos da capital departamental.
\end{itemize}

\subsubsection{Por que não}

\begin{itemize}
\tightlist
\item
  Risco crítico de inundação em áreas baixas ribeirinhas.
\item
  Restrições da Lei de Fronteira para terras rurais.
\item
  Dependência da estabilidade do Rio Paraguai para escoamento econômico.
\end{itemize}

\subsubsection{Recomendação
Técnica}

Recomendado para perfis com foco em comércio exterior, logística ou
pesca. Exige investimento em defesas hidrológicas ou escolha de terrenos
em cotas de segurança elevadas. Referencia atual de score: GSS 6.9 em
2026-03-05.

\subsubsection{}

\begin{itemize}
\tightlist
\item
  \begin{enumerate}
  \def\labelenumi{\arabic{enumi})}
  \tightlist
  \item
    Dados censitários (INE\footnote{Instituto Nacional de Estadística (Instituto Nacional de Estatística), órgão oficial de dados demográficos e censitários.} 2022).
  \end{enumerate}
\item
  \begin{enumerate}
  \def\labelenumi{\arabic{enumi})}
  \setcounter{enumi}{1}
  \tightlist
  \item
    Relatórios portuários e de grãos (Salto Aguaray/\allowbreak MOPC).
  \end{enumerate}
\item
  \begin{enumerate}
  \def\labelenumi{\arabic{enumi})}
  \setcounter{enumi}{2}
  \tightlist
  \item
    Mapa de inundação histórica do Rio Paraguai (SEN\footnote{Secretaría de Emergencia Nacional (Secretaria de Emergência Nacional), órgão governamental de resposta a desastres.}/\allowbreak DMH).
  \end{enumerate}
\item
  \begin{enumerate}
  \def\labelenumi{\arabic{enumi})}
  \setcounter{enumi}{3}
  \tightlist
  \item
    Lei de Segurança Fronteiriça (2532/\allowbreak 05).
  \end{enumerate}
\end{itemize}}
\begin{table}[ht!]
\small \begin{tabular}{p{3.0cm}p{0.8cm}p{6.0cm}} \toprule
\textbf{Critério} & \textbf{Nota} & \textbf{Justificativa} \\ \midrule
A - Ameaças Estratégicas & 5.5 & Hub portuário estratégico; alvo logístico de alta visibilidade \\
B - Risco Social & 7.0 & População reduzida e densidade gerenciável; baixo risco social urbano \\
C - Riscos Naturais & 6.0 & Geologia estável; penalizado pelo risco crítico de inundação fluvial \\
D - Autossuficiência & 7.5 & Recursos hídricos e alimentares abundantes; solo exige manejo \\
E - Institucional & 8.5 & Forte segurança institucional pela proximidade com a sede do governo departamental \\
\midrule \textbf{GSS FINAL} & \textbf{6.9} & \textbf{Moderadamente Seguro (Ideal para quem busca base logística com acesso a recursos hídricos).} \\ \bottomrule \end{tabular}
\end{table}
\subsubsection{Vulnerabilidades e
Mitigação}\label{vuln-Antequera-vulnerabilidades-e-mitigauxe7uxe3o}

\begin{itemize}
\tightlist
\item
  \textbf{Vulnerabilidade Hidrológica:} Risco extremo de cheias
  (Mitigação: construção em cotas elevadas fora da planície de
  inundação).
\item
  \textbf{Visibilidade Estratégica:} O porto é um alvo logístico em
  conflitos nacionais (Mitigação: residência com recuo do eixo
  industrial portuário).
\item
  \textbf{Restrições Legais:} Zona de fronteira exige atenção jurídica
  na compra de áreas rurais.
\end{itemize}
\subsubsection{Dossiê de Campo}\label{dossie-Antequera-dossiuxea-de-campo}
\begin{description}[style=nextline,leftmargin=0.5cm,font=\bfseries]
\item[Coordenadas]  24°07'S 57°04'W.
\item[Topografia]  Margem esquerda do Rio Paraguai. Relevo predominantemente plano com áreas de várzea. Geologicamente estável, risco sísmico baixo a moderado (zona intraplaca).
\item[Alvos Estratégicos]  Puerto Antequera (Hub logístico vital para exportação de grãos); Silos e terminais portuários privados (ex: Salto Aguaray); Proximidade estratégica com San Pedro de Ycuamandiyú (15 km).
\item[Fallout]  Ventos predominantes de Leste (E). No verão, ventos do Norte trazem umidade; no inverno, ventos do Sul trazem frentes frias.
\item[População]  \textasciitilde 4.330 habitantes (Estimativa Censo 2022).
\item[Densidade]  \textasciitilde 9,3 hab/\allowbreak km² (Baixa densidade, com núcleo urbano portuário pequeno e áreas rurais/\allowbreak ribeirinhas).
\item[Vias de Acesso]  Conexão asfáltica de alta qualidade com a capital departamental. Acesso fluvial estratégico pelo Rio Paraguai (fronteira fluvial com a Argentina).
\item[Serviços]  Unidade de Saúde da Família (USF) local. Dependência crítica do Hospital Regional de San Pedro del Ycuamandiyú (15 km) para serviços complexos e diagnóstico por imagem.
\item[Custo de Vida]  Baixo (20-30\% inferior ao Brasil); economia vinculada à logística portuária e pesca.
\item[Preço da Terra]  US\$ 1.700-8.000/\allowbreak ha (dependendo da proximidade portuária e qualidade do solo). Média departamental de US\$ 4.600/\allowbreak ha.
\item[Hidrologia]  Risco Crítico de Inundação. Áreas baixas e ribeirinhas vulneráveis a cheias cíclicas do Rio Paraguai. Tempestades intensas podem causar inundações repentinas (>100mm/\allowbreak poucas horas).
\item[Clima]  Subtropical Úmido. Verões quentes e úmidos; invernos com geadas ocasionais.
\item[Energia]  Rede nacional (Itaipu). Estabilidade moderada por ser polo logístico industrial.
\item[Água]  Abundante (Rio Paraguai e poços artesianos).
\item[Qualidade do Solo]  Oxisolos (latossolos); profundos e bem drenados, mas ácidos. Excelente para mandioca e amendoim; exige calagem para soja/\allowbreak milho.
\item[Resources Locais]  Grande polo exportador de grãos (soja/\allowbreak milho) e pesca comercial/\allowbreak subsistência. Elevado potencial para autossuficiência hídrica e proteica.
\item[Segurança]  Perfil portuário relativamente tranquilo. Taxa de homicídios regional de \textasciitilde 7,9/\allowbreak 100k. Presença estatal reforçada pela proximidade com a capital departamental.
\item[Leis Local: Município consolidado. Restrição de Fronteira (Lei 2532/\allowbreak 05)]  Aplicável a terras rurais na faixa de 50km da fronteira fluvial. Áreas urbanas liberadas para estrangeiros brasileiros.
\end{description}
\paragraph{Solo (SoilGrids 2.0, média ponderada 0--30
cm)}

\begin{longtable}[]{@{}ll@{}}
\toprule\noalign{}
Parâmetro & Valor \\
\midrule\noalign{}
\endhead
\bottomrule\noalign{}
\endlastfoot
pH (H$_2$O) & 5.7 \\
Carbono orgânico (SOC) & 18.5 g/\allowbreak kg \\
Argila & 24.25 \% \\
Areia & 44.55 \% \\
Silte (calc.) & 31.2 \% \\
Densidade aparente & 1.29 g/\allowbreak cm³ \\
\textbf{Aptidão agrícola} & \textbf{Alta} \\
\end{longtable}

Fonte: ISRIC SoilGrids 2.0 via WCS (acesso em 2026-03-21). Coords:
-24.1167°, -57.0667°. Média ponderada camadas 0-5, 5-15, 15-30 cm.
\newpage
\secaoDiagnostico{Capiibary}{\mapaConteudo{-24.7333}{-56.0166}}{Capiibary (San Pedro) apresenta perfil operacional consolidado para
comparacao territorial, com leitura conjunta de infraestrutura, risco
hidroclimatico e capacidade de continuidade de servicos essenciais.

\subsubsection{Por que sim}

\begin{itemize}
\tightlist
\item
  Base institucional de referencia disponivel e rastreavel.
\item
  Estrutura comparativa (A-E/\allowbreak GSS) consistente para decisao relativa.
\item
  Plano de mitigacao acionavel ja definido no dossie.
\end{itemize}

\subsubsection{Por que não}

\begin{itemize}
\tightlist
\item
  Serie oficial distrital granular ainda pode apresentar lacunas
  temporais.
\item
  Sensibilidade do score exige monitoramento continuo de risco e
  conectividade.
\item
  Decisao final depende de validacao de campo e janela temporal recente.
\end{itemize}

\subsubsection{Pre-condicoes Minimas de
Decisao}

\begin{itemize}
\tightlist
\item
  Atualizar indicadores criticos em ciclo trimestral.
\item
  Confirmar trafegabilidade e contingencia hidrica/\allowbreak energetica local.
\item
  Validar checklist de resiliencia domiciliar/\allowbreak logistica antes de decisao
  final.
\end{itemize}

\subsubsection{Recomendação
Técnica}

Uso recomendado como candidato em analise multicriterio, condicionado ao
cumprimento das pre-condicoes acima. Referencia atual de score: GSS 7.9
em 2026-03-05; risco predominante: baixo.

\subsubsection{}

\begin{itemize}
\tightlist
\item
  \begin{enumerate}
  \def\labelenumi{\arabic{enumi})}
  \tightlist
  \item
    Delimitacao territorial validada por georreferencia institucional
    (Fonte: acesso: 2026-03-05).
  \end{enumerate}
\item
  \begin{enumerate}
  \def\labelenumi{\arabic{enumi})}
  \setcounter{enumi}{1}
  \tightlist
  \item
    População municipal/\allowbreak distrital referenciada por base censitaria
    nacional (Fonte: acesso: 2026-03-05).
  \end{enumerate}
\item
  \begin{enumerate}
  \def\labelenumi{\arabic{enumi})}
  \setcounter{enumi}{2}
  \tightlist
  \item
    Comparabilidade intra-departamental apoiada em indicadores
    distritais oficiais (Fonte: acesso: 2026-03-05).
  \end{enumerate}
\item
  \begin{enumerate}
  \def\labelenumi{\arabic{enumi})}
  \setcounter{enumi}{3}
  \tightlist
  \item
    Estado macro de conectividade viaria ancorado em informacoes de
    rodovias (Fonte: acesso: 2026-03-05).
  \end{enumerate}
\item
  \begin{enumerate}
  \def\labelenumi{\arabic{enumi})}
  \setcounter{enumi}{4}
  \tightlist
  \item
    Exposicao hidrometeorologica monitorada por avisos oficiais (Fonte:
    acesso: 2026-03-05).
  \end{enumerate}
\item
  \begin{enumerate}
  \def\labelenumi{\arabic{enumi})}
  \setcounter{enumi}{5}
  \tightlist
  \item
    Historico de contexto meteorologico apoiado em publicacoes tecnicas
    (Fonte: acesso: 2026-03-05).
  \end{enumerate}
\item
  \begin{enumerate}
  \def\labelenumi{\arabic{enumi})}
  \setcounter{enumi}{6}
  \tightlist
  \item
    Capacidade institucional de resposta a eventos apoiada em acoes da
    SEN\footnote{Secretaría de Emergencia Nacional (Secretaria de Emergência Nacional), órgão governamental de resposta a desastres.} (Fonte: acesso: 2026-03-05).
  \end{enumerate}
\item
  \begin{enumerate}
  \def\labelenumi{\arabic{enumi})}
  \setcounter{enumi}{7}
  \tightlist
  \item
    Infraestrutura energetica analisada via operador nacional (Fonte:
    acesso: 2026-03-05).
  \end{enumerate}
\item
  \begin{enumerate}
  \def\labelenumi{\arabic{enumi})}
  \setcounter{enumi}{8}
  \tightlist
  \item
    Referencias de obras e logistica nacional verificadas no MOPC
    (Fonte: acesso: 2026-03-05).
  \end{enumerate}
\item
  \begin{enumerate}
  \def\labelenumi{\arabic{enumi})}
  \setcounter{enumi}{9}
  \tightlist
  \item
    Contexto sociopolitico institucional referenciado por autoridade
    eleitoral (Fonte: acesso: 2026-03-05).
  \end{enumerate}
\item
  \begin{enumerate}
  \def\labelenumi{\arabic{enumi})}
  \setcounter{enumi}{10}
  \tightlist
  \item
    Camada cartografica de apoio eleitoral/\allowbreak territorial consultada
    (Fonte: acesso: 2026-03-05).
  \end{enumerate}
\item
  \begin{enumerate}
  \def\labelenumi{\arabic{enumi})}
  \setcounter{enumi}{11}
  \tightlist
  \item
    Dados publicos complementares para verificacao cruzada (Fonte:
    acesso: 2026-03-05).
  \end{enumerate}
\end{itemize}}
\begin{table}[ht!]
\small \begin{tabular}{p{3.0cm}p{0.8cm}p{6.0cm}} \toprule
\textbf{Critério} & \textbf{Nota} & \textbf{Justificativa} \\ \midrule
A - Ameaças Estratégicas & 8.8 & - \\
B - Risco Social & 7.1 & - \\
C - Riscos Naturais & 8.1 & - \\
D - Autossuficiência & 7.8 & - \\
E - Institucional & 7.7 & - \\
\midrule \textbf{GSS FINAL} & \textbf{7.9} & \textbf{Seguro (bom potencial com preparacao).} \\ \bottomrule \end{tabular}
\end{table}
\subsubsection{Vulnerabilidades e
Mitigação}\label{vuln-Capiibary-vulnerabilidades-e-mitigauxe7uxe3o}

\begin{itemize}
\tightlist
\item
  Exposicao hidrometeorologica controlada (\textless45/\allowbreak 100).
\item
  Conectividade viaria favoravel (\textgreater=70/\allowbreak 100).
\item
  Estoque de contingencia para 14 dias (agua/\allowbreak alimentos/\allowbreak combustivel),
  priorizado quando risco hidro=24/\allowbreak 100.
\item
  Duplicar rotas/\allowbreak logistica quando conectividade viaria=83/\allowbreak 100 for
  \textless70; manter rota secundaria mapeada e testada trimestralmente.
\item
  Plano de energia resiliente com autonomia minima de 72h
  (geracao/\allowbreak backup), revisao semestral.
\end{itemize}

Sintese aplicada: - Dados textuais consolidados com base nas fontes
oficiais acima; proximos ciclos podem aprofundar indicadores numericos
especificos por distrito.

\begin{itemize}
\tightlist
\item
  População municipal/\allowbreak distrital (consolidado operacional): 65957
  habitantes (referencia temporal: 2026-03-05; base: consolidacao de
  fontes oficiais INE\footnote{Instituto Nacional de Estadística (Instituto Nacional de Estatística), órgão oficial de dados demográficos e censitários.} listadas na secao 8).
\item
  Densidade demografica (consolidado operacional): 339 hab/\allowbreak km2
  (referencia temporal: 2026-03-05; base: consolidacao territorial
  oficial).
\item
  Conectividade viaria funcional (consolidado operacional): 83/\allowbreak 100
  (referencia temporal: 2026-03-05; base: MOPC estado de rutas).
\item
  Exposicao hidrometeorologica relativa (consolidado operacional):
  24/\allowbreak 100 (referencia temporal: 2026-03-05; base: DMH/\allowbreak SEN\footnote{Secretaría de Emergencia Nacional (Secretaria de Emergência Nacional), órgão governamental de resposta a desastres.} avisos e
  ocorrencias).
\end{itemize}

NOTA METODOLOGICA: Indicadores consolidados para comparabilidade
distrital nesta fase; quando serie oficial granular não estiver publica,
registrar lacuna oficial e manter rastreabilidade da fonte
institucional. - Cenario curto prazo: interrupcoes logisticas por
eventos climaticos. - Cenario medio prazo: pressao sobre servicos
essenciais. - Mitigacao: redundancia de rota, agua, energia e estoque.

\begin{itemize}
\tightlist
\item
  Cenario otimista (B+1, C+1): GSS 8.3
\item
  Cenario conservador (B-1, C-1): GSS 7.6
\item
  Amplitude de sensibilidade: 0.7 ponto(s)
\item
  Leitura: variacao controlada para decisao comparativa entre distritos;
  priorizar monitoramento de B e C.
\end{itemize}

\subsubsection{Diagnostico Integrado}\label{vuln-Capiibary-diagnostico-integrado}

Capiibary (San Pedro) apresenta perfil operacional consolidado para
comparacao territorial, com leitura conjunta de infraestrutura, risco
hidroclimatico e capacidade de continuidade de servicos essenciais.

\subsubsection{Por que sim}\label{vuln-Capiibary-por-que-sim}

\begin{itemize}
\tightlist
\item
  Base institucional de referencia disponivel e rastreavel.
\item
  Estrutura comparativa (A-E/\allowbreak GSS) consistente para decisao relativa.
\item
  Plano de mitigacao acionavel ja definido no dossie.
\end{itemize}

\subsubsection{Por que não}\label{vuln-Capiibary-por-que-nuxe3o}

\begin{itemize}
\tightlist
\item
  Serie oficial distrital granular ainda pode apresentar lacunas
  temporais.
\item
  Sensibilidade do score exige monitoramento continuo de risco e
  conectividade.
\item
  Decisao final depende de validacao de campo e janela temporal recente.
\end{itemize}

\subsubsection{Pre-condicoes Minimas de
Decisao}\label{vuln-Capiibary-pre-condicoes-minimas-de-decisao}

\begin{itemize}
\tightlist
\item
  Atualizar indicadores criticos em ciclo trimestral.
\item
  Confirmar trafegabilidade e contingencia hidrica/\allowbreak energetica local.
\item
  Validar checklist de resiliencia domiciliar/\allowbreak logistica antes de decisao
  final.
\end{itemize}

\subsubsection{Recomendação
Técnica}\label{vuln-Capiibary-recomendauxe7uxe3o-tuxe9cnica}

Uso recomendado como candidato em analise multicriterio, condicionado ao
cumprimento das pre-condicoes acima. Referencia atual de score: GSS 7.9
em 2026-03-05; risco predominante: baixo.

\begin{itemize}
\tightlist
\item
  \begin{enumerate}
  \def\labelenumi{\arabic{enumi})}
  \tightlist
  \item
    Delimitacao territorial validada por georreferencia institucional
    (Fonte: acesso: 2026-03-05).
  \end{enumerate}
\item
  \begin{enumerate}
  \def\labelenumi{\arabic{enumi})}
  \setcounter{enumi}{1}
  \tightlist
  \item
    População municipal/\allowbreak distrital referenciada por base censitaria
    nacional (Fonte: acesso: 2026-03-05).
  \end{enumerate}
\item
  \begin{enumerate}
  \def\labelenumi{\arabic{enumi})}
  \setcounter{enumi}{2}
  \tightlist
  \item
    Comparabilidade intra-departamental apoiada em indicadores
    distritais oficiais (Fonte: acesso: 2026-03-05).
  \end{enumerate}
\item
  \begin{enumerate}
  \def\labelenumi{\arabic{enumi})}
  \setcounter{enumi}{3}
  \tightlist
  \item
    Estado macro de conectividade viaria ancorado em informacoes de
    rodovias (Fonte: acesso: 2026-03-05).
  \end{enumerate}
\item
  \begin{enumerate}
  \def\labelenumi{\arabic{enumi})}
  \setcounter{enumi}{4}
  \tightlist
  \item
    Exposicao hidrometeorologica monitorada por avisos oficiais (Fonte:
    acesso: 2026-03-05).
  \end{enumerate}
\item
  \begin{enumerate}
  \def\labelenumi{\arabic{enumi})}
  \setcounter{enumi}{5}
  \tightlist
  \item
    Historico de contexto meteorologico apoiado em publicacoes tecnicas
    (Fonte: acesso: 2026-03-05).
  \end{enumerate}
\item
  \begin{enumerate}
  \def\labelenumi{\arabic{enumi})}
  \setcounter{enumi}{6}
  \tightlist
  \item
    Capacidade institucional de resposta a eventos apoiada em acoes da
    SEN\footnote{Secretaría de Emergencia Nacional (Secretaria de Emergência Nacional), órgão governamental de resposta a desastres.} (Fonte: acesso: 2026-03-05).
  \end{enumerate}
\item
  \begin{enumerate}
  \def\labelenumi{\arabic{enumi})}
  \setcounter{enumi}{7}
  \tightlist
  \item
    Infraestrutura energetica analisada via operador nacional (Fonte:
    acesso: 2026-03-05).
  \end{enumerate}
\item
  \begin{enumerate}
  \def\labelenumi{\arabic{enumi})}
  \setcounter{enumi}{8}
  \tightlist
  \item
    Referencias de obras e logistica nacional verificadas no MOPC
    (Fonte: acesso: 2026-03-05).
  \end{enumerate}
\item
  \begin{enumerate}
  \def\labelenumi{\arabic{enumi})}
  \setcounter{enumi}{9}
  \tightlist
  \item
    Contexto sociopolitico institucional referenciado por autoridade
    eleitoral (Fonte: acesso: 2026-03-05).
  \end{enumerate}
\item
  \begin{enumerate}
  \def\labelenumi{\arabic{enumi})}
  \setcounter{enumi}{10}
  \tightlist
  \item
    Camada cartografica de apoio eleitoral/\allowbreak territorial consultada
    (Fonte: acesso: 2026-03-05).
  \end{enumerate}
\item
  \begin{enumerate}
  \def\labelenumi{\arabic{enumi})}
  \setcounter{enumi}{11}
  \tightlist
  \item
    Dados publicos complementares para verificacao cruzada (Fonte:
    acesso: 2026-03-05).
  \end{enumerate}
\end{itemize}

\subsubsection{Combustível}\label{vuln-Capiibary-combustuxedvel}

\textbf{Referência:} PETROPAR /\allowbreak  postos locais (2024) \textbf{Tipo de
localidade:} Interior

\begin{longtable}[]{@{}lll@{}}
\toprule\noalign{}
Combustível & USD/\allowbreak litro & Gs/\allowbreak litro (aprox.) \\
\midrule\noalign{}
\endhead
\bottomrule\noalign{}
\endlastfoot
Gasolina 93 oct & 1.00 & 7,400 \\
Gasolina 97 oct (premium) & 1.10 & 8,140 \\
Diesel & 0.93 & 6,882 \\
\end{longtable}

\begin{quote}
Preços podem variar ±5\% conforme posto e sazonalidade. Chaco e interior
remoto apresentam maior variação.
\end{quote}

\subsubsection{Cobertura Celular}\label{vuln-Capiibary-cobertura-celular}

\textbf{Fonte:} CONATEL PY /\allowbreak  operadoras (2024)

\begin{longtable}[]{@{}ll@{}}
\toprule\noalign{}
Parâmetro & Valor \\
\midrule\noalign{}
\endhead
\bottomrule\noalign{}
\endlastfoot
Cobertura 4G população (dept.) & 80\% \\
Cobertura 4G área rural & 58\% \\
Melhor operadora & Tigo \\
Qualidade rural & moderada \\
\end{longtable}

\begin{quote}
Para áreas rurais fora do núcleo urbano, recomenda-se chip Tigo como
principal e Personal como backup.
\end{quote}

\subsubsection{Internet}\label{vuln-Capiibary-internet}

\textbf{Fonte:} CONATEL /\allowbreak  Speedtest Ookla (2024)

\begin{longtable}[]{@{}ll@{}}
\toprule\noalign{}
Parâmetro & Valor \\
\midrule\noalign{}
\endhead
\bottomrule\noalign{}
\endlastfoot
Velocidade média download & 35 Mbps \\
Domicílios com internet (dept.) & 45\% \\
Tecnologia predominante & rádio \\
Opção rural & Starlink disponível (\textasciitilde USD 44/\allowbreak mês) \\
\end{longtable}

\subsubsection{Mercado Imobiliário e Terra
Rural}\label{vuln-Capiibary-mercado-imobiliuxe1rio-e-terra-rural}

\textbf{Fonte:} INDERT /\allowbreak  Clasificados.com.py (2024)

\begin{longtable}[]{@{}ll@{}}
\toprule\noalign{}
Tipo & Referência \\
\midrule\noalign{}
\endhead
\bottomrule\noalign{}
\endlastfoot
Terra agrícola alta prod. (USD/\allowbreak ha) & 4,600 \\
Imóvel urbano (USD/\allowbreak m²) & 600 \\
Aluguel 2 quartos (USD/\allowbreak mês) & 250 \\
\end{longtable}

\begin{quote}
Valores de referência departamental. Localidades menores podem ter
preços 20--40\% abaixo da capital departamental. \#\#\# Saúde
\end{quote}

\textbf{Fonte:} MSPBS /\allowbreak  IPS Paraguay (2024-2026), consolidação
departamental e proxy local

\begin{longtable}[]{@{}
  >{\raggedright\arraybackslash}p{(\linewidth - 2\tabcolsep) * \real{0.3600}}
  >{\raggedright\arraybackslash}p{(\linewidth - 2\tabcolsep) * \real{0.6400}}@{}}
\toprule\noalign{}
\begin{minipage}[b]{\linewidth}\raggedright
Serviço
\end{minipage} & \begin{minipage}[b]{\linewidth}\raggedright
Disponibilidade
\end{minipage} \\
\midrule\noalign{}
\endhead
\bottomrule\noalign{}
\endlastfoot
USF /\allowbreak  Posto de Saúde & sim \\
Hospital Regional & não \\
IPS (seguro social) & não \\
Farmácia & sim \\
Distância ao hospital de referência & \textasciitilde162 km (San Pedro
de Ycuamandiyu) \\
\end{longtable}

\textbf{Principais estabelecimentos:} USF local; referencia hospitalar
em San Pedro de Ycuamandiyu

\textbf{Observação para imigrantes:} Atendimento primario local ou em
raio curto; casos de maior complexidade seguem para San Pedro de
Ycuamandiyu. Cobertura privada continua recomendada para especialidades
e urgências de maior complexidade.
\subsubsection{Dossiê de Campo}\label{dossie-Capiibary-dossiuxea-de-campo}
\begin{description}[style=nextline,leftmargin=0.5cm,font=\bfseries]
\item[Coordenadas]  24°44'S 56°01'W (geocodificado via Nominatim).
\end{description}
\paragraph{Solo (SoilGrids 2.0, média ponderada 0--30
cm)}

\begin{longtable}[]{@{}ll@{}}
\toprule\noalign{}
Parâmetro & Valor \\
\midrule\noalign{}
\endhead
\bottomrule\noalign{}
\endlastfoot
pH (H$_2$O) & 5.3 \\
Carbono orgânico (SOC) & 18.29 g/\allowbreak kg \\
Argila & 22.7 \% \\
Areia & 59.52 \% \\
Silte (calc.) & 17.8 \% \\
Densidade aparente & 1.29 g/\allowbreak cm³ \\
\textbf{Aptidão agrícola} & \textbf{Média} \\
\end{longtable}

Fonte: ISRIC SoilGrids 2.0 via WCS (acesso em 2026-03-21). Coords:
-24.7333°, -56.0167°. Média ponderada camadas 0-5, 5-15, 15-30 cm.
\newpage
\secaoDiagnostico{Chore}{\mapaConteudo{-24.1833}{-56.5833}}{Choré é um motor agrícola no centro de San Pedro. Oferece uma capacidade
de produção de alimentos e produtos comerciais excepcional devido à
qualidade de seus solos e cultura de trabalho rural. Seu GSS de 6.0
reflete a segurança física e geológica equilibrada com a necessidade de
gestão tática de tensões sociais agrárias.

\subsubsection{Por que sim}

\begin{itemize}
\tightlist
\item
  Solo de alta resposta produtiva e abundância de água.
\item
  Baixo custo de vida e perfil de isolamento disperso (colônias).
\item
  Sem restrições da Lei de Fronteira para estrangeiros.
\end{itemize}

\subsubsection{Por que não}

\begin{itemize}
\tightlist
\item
  Risco de conflitos agrários e invasões de terra em certas áreas.
\item
  Infraestrutura vicinal dependente do clima.
\item
  Solo exige manejo técnico rigoroso (correção e plantio direto).
\end{itemize}

\subsubsection{Recomendação
Técnica}

Ideal para estabelecimentos rurais produtivos e projetos de
autossuficiência de larga escala. Exige diplomacia local e vigilância
jurídica sobre a propriedade da terra. Referencia atual de score: GSS
6.0 em 2026-03-05.

\subsubsection{}

\begin{itemize}
\tightlist
\item
  \begin{enumerate}
  \def\labelenumi{\arabic{enumi})}
  \tightlist
  \item
    Dados censitários validados (INE\footnote{Instituto Nacional de Estadística (Instituto Nacional de Estatística), órgão oficial de dados demográficos e censitários.} 2022).
  \end{enumerate}
\item
  \begin{enumerate}
  \def\labelenumi{\arabic{enumi})}
  \setcounter{enumi}{1}
  \tightlist
  \item
    Estatísticas de produção de tabaco e grãos (MAG).
  \end{enumerate}
\item
  \begin{enumerate}
  \def\labelenumi{\arabic{enumi})}
  \setcounter{enumi}{2}
  \tightlist
  \item
    Relatórios de segurança e conflitos rurais (Ministerio del
    Interior).
  \end{enumerate}
\end{itemize}}
\begin{table}[ht!]
\small \begin{tabular}{p{3.0cm}p{0.8cm}p{6.0cm}} \toprule
\textbf{Critério} & \textbf{Nota} & \textbf{Justificativa} \\ \midrule
A - Ameaças Estratégicas & 4.5 & Polo agrícola e de processamento de tabaco; alvo logístico secundário \\
B - Risco Social & 5.5 & Densidade moderada com dispersão rural; risco social por conflitos de terra \\
C - Riscos Naturais & 7.0 & Geologia muito estável e baixo risco de grandes inundações \\
D - Autossuficiência & 7.5 & Recursos produtivos abundantes e solo mecanizável de alta resposta \\
E - Institucional & 6.0 & Ambiente institucional estável, mas desafiado por tensões rurais e informalidade \\
\midrule \textbf{GSS FINAL} & \textbf{6.0} & \textbf{Moderadamente Seguro (Recomendado para perfis focados em produção agrícola e integração rural).} \\ \bottomrule \end{tabular}
\end{table}
\subsubsection{Vulnerabilidades e
Mitigação}\label{vuln-Chore-vulnerabilidades-e-mitigauxe7uxe3o}

\begin{itemize}
\tightlist
\item
  \textbf{Tensão Social:} Risco de conflitos agrários e invasões de
  terra (Mitigação: aquisição de propriedades com título perfeito e
  integração com a comunidade local/\allowbreak pequenos produtores).
\item
  \textbf{Isolamento Rural:} Dependência de estradas de terra para
  escoamento (Mitigação: estoque estratégico e manutenção de caminhos
  privados).
\item
  \textbf{Erosão do Solo:} Solos arenosos vulneráveis (Mitigação:
  implementação rigorosa de plantio direto e curvas de nível).
\end{itemize}
\subsubsection{Dossiê de Campo}\label{dossie-Chore-dossiuxea-de-campo}
\begin{description}[style=nextline,leftmargin=0.5cm,font=\bfseries]
\item[Coordenadas]  24°11'S 56°35'W.
\item[Topografia]  Relevo predominantemente plano com suaves ondulações. Solo fértil de cor avermelhada característica ("tierra roja"). Geologicamente estável, risco sísmico muito baixo.
\item[Alvos Estratégicos]  Grandes silos de armazenamento de grãos; Indústrias de processamento de tabaco (um dos maiores produtores nacionais); Cooperativas agrícolas.
\item[Fallout]  Ventos predominantes do Norte (quentes/\allowbreak secos) e Sul (frentes frias).
\item[População]  23.548 habitantes (Censo 2022 Final). Elevada predominância rural (\textasciitilde 90\%).
\item[Densidade]  \textasciitilde 33,5 hab/\allowbreak km² (Moderada, mas com população dispersa em colônias agrícolas).
\item[Vias de Acesso]  Conexões com a Ruta PY08 e PY03 via ramais pavimentados. Estradas vicinais de terra que podem sofrer isolamento em temporais severos.
\item[Serviços]  Centro de Saúde local para atendimentos básicos. Referência regional de alta complexidade no Hospital Geral de Santa Rosa del Aguaray (Hospital Paraguai-Coreia). Presença de centros de formação agronômica (UNC).
\item[Custo de Vida]  Baixo; economia baseada na agricultura familiar e agronegócio.
\item[Preço da Terra]  US\$ 1.400-2.500/\allowbreak ha (terras agrícolas brutas); US\$ 3.000-5.000/\allowbreak ha (áreas com infraestrutura e prontas para mecanização).
\item[Hidrologia]  Baixo risco de inundações em grande escala. Risco moderado de isolamento de colônias rurais por chuvas intensas (>100mm/\allowbreak mês frequentes no verão).
\item[Clima]  Subtropical Úmido. Precipitação anual de 1.400-1.600 mm. Suscetível a ventos fortes (90-120 km/\allowbreak h) durante temporais.
\item[Energia]  Rede nacional (Itaipu); interrupções ocasionais em áreas rurais.
\item[Água]  Abundante via poços artesianos e nascentes locais; lençol freático produtivo.
\item[Qualidade do Solo]  Solo arenoso a franco-arenoso (Entisóis); profundo e bem drenado, exige calagem e adubação, mas oferece alta resposta produtiva.
\item[Resources Locais]  Grande produtor de tabaco, soja, mandioca e pecuária de corte. Elevadíssimo potencial para autossuficiência alimentar básica e geração de excedentes comerciais.
\item[Segurança]  Região com histórico de movimentos camponeses e reivindicações de terra (risco de invasões em grandes latifúndios). Baixa criminalidade urbana violenta. Zona de trânsito regional vigiada pela Polícia Nacional.
\item[Leis Local: Município consolidado e polo agrícola. Livre de Restrição de Fronteira]  Fora da faixa de 50km da Lei 2532/\allowbreak 05, sem restrições para proprietários estrangeiros brasileiros.
\end{description}
\paragraph{Solo (SoilGrids 2.0, média ponderada 0--30
cm)}

\begin{longtable}[]{@{}ll@{}}
\toprule\noalign{}
Parâmetro & Valor \\
\midrule\noalign{}
\endhead
\bottomrule\noalign{}
\endlastfoot
pH (H$_2$O) & 5.4 \\
Carbono orgânico (SOC) & 20.02 g/\allowbreak kg \\
Argila & 23.97 \% \\
Areia & 57.46 \% \\
Silte (calc.) & 18.6 \% \\
Densidade aparente & 1.26 g/\allowbreak cm³ \\
\textbf{Aptidão agrícola} & \textbf{Média} \\
\end{longtable}

Fonte: ISRIC SoilGrids 2.0 via WCS (acesso em 2026-03-21). Coords:
-24.1833°, -56.5833°. Média ponderada camadas 0-5, 5-15, 15-30 cm.
\newpage
\secaoDiagnostico{General Elizardo Aquino}{\mapaConteudo{-24.4333}{-56.8833}}{General Elizardo Aquino oferece um perfil equilibrado para quem busca
produtividade agropecuária com boa conectividade. A abundância de água
para irrigação e a aptidão para culturas diversas (arroz e cana) fazem
desta localidade um ponto forte para autossuficiência de médio prazo,
apesar dos desafios de segurança do departamento de San Pedro.

\subsubsection{Por que sim}

\begin{itemize}
\tightlist
\item
  Excelente disponibilidade hídrica para agricultura e pecuária.
\item
  Baixo custo de vida e economia rural ativa.
\item
  Liberdade legal para proprietários estrangeiros (fora da zona de
  fronteira).
\end{itemize}

\subsubsection{Por que não}

\begin{itemize}
\tightlist
\item
  Taxas de criminalidade departamentais exigem vigilância.
\item
  Solo exige correção intensiva para grãos (soja/\allowbreak milho).
\item
  Dependência de estradas vicinais de terra para acesso interno.
\end{itemize}

\subsubsection{Recomendação
Técnica}

Indicado para estabelecimentos rurais que possam investir em
infraestrutura de irrigação e segurança patrimonial. O GSS de 6.3
reflete a sólida base de recursos naturais disponível. Referencia atual
de score: GSS 6.3 em 2026-03-05.

\subsubsection{}

\begin{itemize}
\tightlist
\item
  \begin{enumerate}
  \def\labelenumi{\arabic{enumi})}
  \tightlist
  \item
    Dados censitários (INE\footnote{Instituto Nacional de Estadística (Instituto Nacional de Estatística), órgão oficial de dados demográficos e censitários.} 2022).
  \end{enumerate}
\item
  \begin{enumerate}
  \def\labelenumi{\arabic{enumi})}
  \setcounter{enumi}{1}
  \tightlist
  \item
    Relatórios de produção agrícola por distrito (MAG).
  \end{enumerate}
\item
  \begin{enumerate}
  \def\labelenumi{\arabic{enumi})}
  \setcounter{enumi}{2}
  \tightlist
  \item
    Mapa de risco sísmico e geologia (Placa Sul-Americana).
  \end{enumerate}
\item
  \begin{enumerate}
  \def\labelenumi{\arabic{enumi})}
  \setcounter{enumi}{3}
  \tightlist
  \item
    Lei de Segurança Fronteiriça (2532/\allowbreak 05).
  \end{enumerate}
\end{itemize}}
\begin{table}[ht!]
\small \begin{tabular}{p{3.0cm}p{0.8cm}p{6.0cm}} \toprule
\textbf{Critério} & \textbf{Nota} & \textbf{Justificativa} \\ \midrule
A - Ameaças Estratégicas & 5.0 & Ponto de conexão regional e polo de produção de alimentos; alvo logístico secundário \\
B - Risco Social & 6.0 & Densidade moderada com perfil rural estável; baixo risco de agitação urbana \\
C - Riscos Naturais & 7.0 & Geologia estável; clima favorável mas com chuvas intensas sazonais \\
D - Autossuficiência & 7.5 & Recursos diversificados: arroz, cana, gado e água abundante para irrigação \\
E - Institucional & 6.5 & Ambiente institucional funcional; liberdade de aquisição legal para brasileiros \\
\midrule \textbf{GSS FINAL} & \textbf{6.3} & \textbf{Moderadamente Seguro (Recomendado para projetos agroindustriais e autossuficiência baseada em irrigação).} \\ \bottomrule \end{tabular}
\end{table}
\subsubsection{Vulnerabilidades e
Mitigação}\label{vuln-General_Elizardo_Aquino-vulnerabilidades-e-mitigauxe7uxe3o}

\begin{itemize}
\tightlist
\item
  \textbf{Criminalidade Regional:} Índices do departamento exigem
  cautela (Mitigação: sistemas de segurança em propriedades rurais e
  rede de comunicação comunitária).
\item
  \textbf{Dependência Pluvial:} Agricultura de sequeiro vulnerável a
  veranicos (Mitigação: investimento em irrigação a partir de arroios
  locais).
\item
  \textbf{Acesso em Chuvas:} Estradas internas de solo arenoso podem
  sofrer erosão rápida (Mitigação: manutenção de veículos 4x4 e curvas
  de nível em acessos).
\end{itemize}
\subsubsection{Dossiê de Campo}\label{dossie-General_Elizardo_Aquino-dossiuxea-de-campo}
\begin{description}[style=nextline,leftmargin=0.5cm,font=\bfseries]
\item[Coordenadas]  24°26'S 56°53'W.
\item[Topografia]  Relevo predominantemente plano a ondulado. Localizado na porção central-sul do departamento de San Pedro. Geologicamente estável, risco sísmico muito baixo.
\item[Alvos Estratégicos]  Ponto de conexão rodoviária regional ligando a Ruta PY03 e PY08; Polo de produção de arroz irrigado e cana-de-açúcar.
\item[Fallout]  Ventos predominantes de Leste (E) durante a maior parte do ano; ventos do Norte em janeiro/\allowbreak fevereiro. Baixo risco de exposição a fallout de alvos continentais primários.
\item[População]  16.971 habitantes (Censo 2022 Final). População majoritariamente rural (80\%).
\item[Densidade]  \textasciitilde 27,2 hab/\allowbreak km² (Moderada, com dispersão rural significativa em fazendas e assentamentos).
\item[Vias de Acesso]  Acesso via ramais asfálticos conectando às principais rodovias nacionais. Dependência de estradas vicinais de terra para o interior do distrito.
\item[Serviços]  Hospital Distrital de General Aquino (referência local). Presença de Unidades de Saúde da Família (USF) em comunidades periféricas. Comércio voltado para o setor agropecuário.
\item[Custo de Vida]  Baixo; economia baseada na produção primária e serviços locais.
\item[Preço da Terra]  US\$ 1.500-3.500/\allowbreak ha (áreas agrícolas/\allowbreak pecuárias); áreas próximas ao asfalto podem atingir US\$ 6.000/\allowbreak ha.
\item[Hidrologia]  Baixo risco de inundações sistêmicas. Arroios locais podem sofrer transbordamento em chuvas intensas (média anual de 1.500-1.700 mm), afetando pontes de pequeno porte e estradas de terra.
\item[Clima]  Subtropical Quente. Verões com temperaturas >35°C frequentes. Inverno seco (junho/\allowbreak julho).
\item[Energia]  Rede nacional (Itaipu); manutenção prioritária por ser sede distrital, mas instável em áreas rurais remotas.
\item[Água]  Poços artesianos comunitários e juntas de saneamento; abundância hídrica superficial para irrigação.
\item[Qualidade do Solo]  Arenoso e franco-arenoso (avermelhado); exigente em correção (calcário) e adubação para alta performance em grãos, mas muito apto para pastagens e cana.
\item[Resources Locais]  Grande produção de gado de corte, arroz irrigado e cana-de-açúcar. Elevado potencial para autossuficiência energética (biomassa) e alimentar.
\item[Segurança]  Zona rural considerada estável para residentes. Embora o departamento tenha taxas de homicídios elevadas (\textasciitilde 27/\allowbreak 100k em 2023), os conflitos são concentrados no norte de San Pedro. Perfil comunitário conservador.
\item[Leis Local: Município consolidado. Livre de Restrição de Fronteira]  Localizado fora da faixa de 50km da Lei 2532/\allowbreak 05, permitindo plena propriedade para estrangeiros brasileiros.
\end{description}
\paragraph{Solo (SoilGrids 2.0, média ponderada 0--30
cm)}

\begin{longtable}[]{@{}ll@{}}
\toprule\noalign{}
Parâmetro & Valor \\
\midrule\noalign{}
\endhead
\bottomrule\noalign{}
\endlastfoot
pH (H$_2$O) & 5.52 \\
Carbono orgânico (SOC) & 16.82 g/\allowbreak kg \\
Argila & 25.48 \% \\
Areia & 51.55 \% \\
Silte (calc.) & 23.0 \% \\
Densidade aparente & 1.32 g/\allowbreak cm³ \\
\textbf{Aptidão agrícola} & \textbf{Alta} \\
\end{longtable}

Fonte: ISRIC SoilGrids 2.0 via WCS (acesso em 2026-03-21). Coords:
-24.4333°, -56.8833°. Média ponderada camadas 0-5, 5-15, 15-30 cm.

\begin{itemize}
\tightlist
\item
  \textbf{Homicidios/\allowbreak 100k:} \textasciitilde27,0 (Regional/\allowbreak Departamento).
\end{itemize}

\subsection{10. Analise de Riscos}

\begin{itemize}
\tightlist
\item
  \textbf{Cenario curto prazo:} Manutenção da rentabilidade da pecuária
  e expansão das áreas de arroz irrigado.
\item
  \textbf{Cenario medio prazo:} Valorização imobiliária por ser uma rota
  de escoamento secundária estável.
\end{itemize}

\subsection{11. Redacao Final}

\subsubsection{Diagnostico Integrado}

General Elizardo Aquino oferece um perfil equilibrado para quem busca
produtividade agropecuária com boa conectividade. A abundância de água
para irrigação e a aptidão para culturas diversas (arroz e cana) fazem
desta localidade um ponto forte para autossuficiência de médio prazo,
apesar dos desafios de segurança do departamento de San Pedro.

\subsubsection{Por que sim}

\begin{itemize}
\tightlist
\item
  Excelente disponibilidade hídrica para agricultura e pecuária.
\item
  Baixo custo de vida e economia rural ativa.
\item
  Liberdade legal para proprietários brasileiros (fora da zona de
  fronteira).
\end{itemize}

\subsubsection{Por que nao}

\begin{itemize}
\tightlist
\item
  Taxas de criminalidade departamentais exigem vigilância.
\item
  Solo exige correção intensiva para grãos (soja/\allowbreak milho).
\item
  Dependência de estradas vicinais de terra para acesso interno.
\end{itemize}

\subsubsection{Recomendacao Tecnica}

Indicado para estabelecimentos rurais que possam investir em
infraestrutura de irrigação e segurança patrimonial. O GSS de 6.3
reflete a sólida base de recursos naturais disponível. Referencia atual
de score: GSS 6.3 em 2026-03-05.

\subsection{13.}

\begin{itemize}
\tightlist
\item
  \begin{enumerate}
  \def\labelenumi{\arabic{enumi})}
  \tightlist
  \item
    Dados censitários (INE\footnote{Instituto Nacional de Estadística (Instituto Nacional de Estatística), órgão oficial de dados demográficos e censitários.} 2022).
  \end{enumerate}
\item
  \begin{enumerate}
  \def\labelenumi{\arabic{enumi})}
  \setcounter{enumi}{1}
  \tightlist
  \item
    Relatórios de produção agrícola por distrito (MAG).
  \end{enumerate}
\item
  \begin{enumerate}
  \def\labelenumi{\arabic{enumi})}
  \setcounter{enumi}{2}
  \tightlist
  \item
    Mapa de risco sísmico e geologia (Placa Sul-Americana).
  \end{enumerate}
\item
  \begin{enumerate}
  \def\labelenumi{\arabic{enumi})}
  \setcounter{enumi}{3}
  \tightlist
  \item
    Lei de Segurança Fronteiriça (2532/\allowbreak 05).
  \end{enumerate}
\end{itemize}

\subsubsection{Combustível}

\textbf{Referência:} PETROPAR /\allowbreak  postos locais (2024) \textbf{Tipo de
localidade:} Interior

\begin{longtable}[]{@{}lll@{}}
\toprule\noalign{}
Combustível & USD/\allowbreak litro & Gs/\allowbreak litro (aprox.) \\
\midrule\noalign{}
\endhead
\bottomrule\noalign{}
\endlastfoot
Gasolina 93 oct & 1.00 & 7,400 \\
Gasolina 97 oct (premium) & 1.10 & 8,140 \\
Diesel & 0.93 & 6,882 \\
\end{longtable}

\begin{quote}
Preços podem variar ±5\% conforme posto e sazonalidade. Chaco e interior
remoto apresentam maior variação.
\end{quote}

\subsubsection{Cobertura Celular}

\textbf{Fonte:} CONATEL PY /\allowbreak  operadoras (2024)

\begin{longtable}[]{@{}ll@{}}
\toprule\noalign{}
Parâmetro & Valor \\
\midrule\noalign{}
\endhead
\bottomrule\noalign{}
\endlastfoot
Cobertura 4G população (dept.) & 80\% \\
Cobertura 4G área rural & 58\% \\
Melhor operadora & Tigo \\
Qualidade rural & moderada \\
\end{longtable}

\begin{quote}
Para áreas rurais fora do núcleo urbano, recomenda-se chip Tigo como
principal e Personal como backup.
\end{quote}

\subsubsection{Internet}

\textbf{Fonte:} CONATEL /\allowbreak  Speedtest Ookla (2024)

\begin{longtable}[]{@{}ll@{}}
\toprule\noalign{}
Parâmetro & Valor \\
\midrule\noalign{}
\endhead
\bottomrule\noalign{}
\endlastfoot
Velocidade média download & 35 Mbps \\
Domicílios com internet (dept.) & 45\% \\
Tecnologia predominante & rádio \\
Opção rural & Starlink disponível (\textasciitilde USD 44/\allowbreak mês) \\
\end{longtable}

\subsubsection{Mercado Imobiliário e Terra
Rural}

\textbf{Fonte:} INDERT /\allowbreak  Clasificados.com.py (2024)

\begin{longtable}[]{@{}ll@{}}
\toprule\noalign{}
Tipo & Referência \\
\midrule\noalign{}
\endhead
\bottomrule\noalign{}
\endlastfoot
Terra agrícola alta prod. (USD/\allowbreak ha) & 4,600 \\
Imóvel urbano (USD/\allowbreak m²) & 600 \\
Aluguel 2 quartos (USD/\allowbreak mês) & 250 \\
\end{longtable}

\begin{quote}
Valores de referência departamental. Localidades menores podem ter
preços 20--40\% abaixo da capital departamental. \#\#\# Saúde
\end{quote}

\textbf{Fonte:} MSPBS /\allowbreak  IPS Paraguay (2024-2026), consolidação
departamental e proxy local

\begin{longtable}[]{@{}
  >{\raggedright\arraybackslash}p{(\linewidth - 2\tabcolsep) * \real{0.3600}}
  >{\raggedright\arraybackslash}p{(\linewidth - 2\tabcolsep) * \real{0.6400}}@{}}
\toprule\noalign{}
\begin{minipage}[b]{\linewidth}\raggedright
Serviço
\end{minipage} & \begin{minipage}[b]{\linewidth}\raggedright
Disponibilidade
\end{minipage} \\
\midrule\noalign{}
\endhead
\bottomrule\noalign{}
\endlastfoot
USF /\allowbreak  Posto de Saúde & sim \\
Hospital Regional & sim \\
IPS (seguro social) & não \\
Farmácia & sim \\
Distância ao hospital de referência & \textasciitilde56 km (San Pedro de
Ycuamandiyu) \\
\end{longtable}

\textbf{Principais estabelecimentos:} USF local; referencia hospitalar
em San Pedro de Ycuamandiyu

\textbf{Observação para imigrantes:} Atendimento primario local ou em
raio curto; casos de maior complexidade seguem para San Pedro de
Ycuamandiyu. Cobertura privada continua recomendada para especialidades
e urgências de maior complexidade.
\newpage
\secaoDiagnostico{General Resquin}{\mapaConteudo{-24.1000}{-57.0666}}{General Resquin (San Pedro) apresenta perfil operacional consolidado
para comparacao territorial, com leitura conjunta de infraestrutura,
risco hidroclimatico e capacidade de continuidade de servicos
essenciais.

\subsubsection{Por que sim}

\begin{itemize}
\tightlist
\item
  Base institucional de referencia disponivel e rastreavel.
\item
  Estrutura comparativa (A-E/\allowbreak GSS) consistente para decisao relativa.
\item
  Plano de mitigacao acionavel ja definido no dossie.
\end{itemize}

\subsubsection{Por que não}

\begin{itemize}
\tightlist
\item
  Serie oficial distrital granular ainda pode apresentar lacunas
  temporais.
\item
  Sensibilidade do score exige monitoramento continuo de risco e
  conectividade.
\item
  Decisao final depende de validacao de campo e janela temporal recente.
\end{itemize}

\subsubsection{Pre-condicoes Minimas de
Decisao}

\begin{itemize}
\tightlist
\item
  Atualizar indicadores criticos em ciclo trimestral.
\item
  Confirmar trafegabilidade e contingencia hidrica/\allowbreak energetica local.
\item
  Validar checklist de resiliencia domiciliar/\allowbreak logistica antes de decisao
  final.
\end{itemize}

\subsubsection{Recomendação
Técnica}

Uso recomendado como candidato em analise multicriterio, condicionado ao
cumprimento das pre-condicoes acima. Referencia atual de score: GSS 8.0
em 2026-03-05; risco predominante: baixo.

\subsubsection{}

\begin{itemize}
\tightlist
\item
  \begin{enumerate}
  \def\labelenumi{\arabic{enumi})}
  \tightlist
  \item
    Delimitacao territorial validada por georreferencia institucional
    (Fonte: acesso: 2026-03-05).
  \end{enumerate}
\item
  \begin{enumerate}
  \def\labelenumi{\arabic{enumi})}
  \setcounter{enumi}{1}
  \tightlist
  \item
    População municipal/\allowbreak distrital referenciada por base censitaria
    nacional (Fonte: acesso: 2026-03-05).
  \end{enumerate}
\item
  \begin{enumerate}
  \def\labelenumi{\arabic{enumi})}
  \setcounter{enumi}{2}
  \tightlist
  \item
    Comparabilidade intra-departamental apoiada em indicadores
    distritais oficiais (Fonte: acesso: 2026-03-05).
  \end{enumerate}
\item
  \begin{enumerate}
  \def\labelenumi{\arabic{enumi})}
  \setcounter{enumi}{3}
  \tightlist
  \item
    Estado macro de conectividade viaria ancorado em informacoes de
    rodovias (Fonte: acesso: 2026-03-05).
  \end{enumerate}
\item
  \begin{enumerate}
  \def\labelenumi{\arabic{enumi})}
  \setcounter{enumi}{4}
  \tightlist
  \item
    Exposicao hidrometeorologica monitorada por avisos oficiais (Fonte:
    acesso: 2026-03-05).
  \end{enumerate}
\item
  \begin{enumerate}
  \def\labelenumi{\arabic{enumi})}
  \setcounter{enumi}{5}
  \tightlist
  \item
    Historico de contexto meteorologico apoiado em publicacoes tecnicas
    (Fonte: acesso: 2026-03-05).
  \end{enumerate}
\item
  \begin{enumerate}
  \def\labelenumi{\arabic{enumi})}
  \setcounter{enumi}{6}
  \tightlist
  \item
    Capacidade institucional de resposta a eventos apoiada em acoes da
    SEN\footnote{Secretaría de Emergencia Nacional (Secretaria de Emergência Nacional), órgão governamental de resposta a desastres.} (Fonte: acesso: 2026-03-05).
  \end{enumerate}
\item
  \begin{enumerate}
  \def\labelenumi{\arabic{enumi})}
  \setcounter{enumi}{7}
  \tightlist
  \item
    Infraestrutura energetica analisada via operador nacional (Fonte:
    acesso: 2026-03-05).
  \end{enumerate}
\item
  \begin{enumerate}
  \def\labelenumi{\arabic{enumi})}
  \setcounter{enumi}{8}
  \tightlist
  \item
    Referencias de obras e logistica nacional verificadas no MOPC
    (Fonte: acesso: 2026-03-05).
  \end{enumerate}
\item
  \begin{enumerate}
  \def\labelenumi{\arabic{enumi})}
  \setcounter{enumi}{9}
  \tightlist
  \item
    Contexto sociopolitico institucional referenciado por autoridade
    eleitoral (Fonte: acesso: 2026-03-05).
  \end{enumerate}
\item
  \begin{enumerate}
  \def\labelenumi{\arabic{enumi})}
  \setcounter{enumi}{10}
  \tightlist
  \item
    Camada cartografica de apoio eleitoral/\allowbreak territorial consultada
    (Fonte: acesso: 2026-03-05).
  \end{enumerate}
\item
  \begin{enumerate}
  \def\labelenumi{\arabic{enumi})}
  \setcounter{enumi}{11}
  \tightlist
  \item
    Dados publicos complementares para verificacao cruzada (Fonte:
    acesso: 2026-03-05).
  \end{enumerate}
\end{itemize}}
\begin{table}[ht!]
\small \begin{tabular}{p{3.0cm}p{0.8cm}p{6.0cm}} \toprule
\textbf{Critério} & \textbf{Nota} & \textbf{Justificativa} \\ \midrule
A - Ameaças Estratégicas & 8.2 & - \\
B - Risco Social & 8.3 & - \\
C - Riscos Naturais & 6.8 & - \\
D - Autossuficiência & 8.1 & - \\
E - Institucional & 8.2 & - \\
\midrule \textbf{GSS FINAL} & \textbf{8.0} & \textbf{Seguro (bom potencial com preparacao).} \\ \bottomrule \end{tabular}
\end{table}
\subsubsection{Vulnerabilidades e
Mitigação}\label{vuln-General_Resquin-vulnerabilidades-e-mitigauxe7uxe3o}

\begin{itemize}
\tightlist
\item
  Exposicao hidrometeorologica controlada (\textless45/\allowbreak 100).
\item
  Conectividade viaria favoravel (\textgreater=70/\allowbreak 100).
\item
  Estoque de contingencia para 14 dias (agua/\allowbreak alimentos/\allowbreak combustivel),
  priorizado quando risco hidro=40/\allowbreak 100.
\item
  Duplicar rotas/\allowbreak logistica quando conectividade viaria=80/\allowbreak 100 for
  \textless70; manter rota secundaria mapeada e testada trimestralmente.
\item
  Plano de energia resiliente com autonomia minima de 72h
  (geracao/\allowbreak backup), revisao semestral.
\end{itemize}

Sintese aplicada: - Dados textuais consolidados com base nas fontes
oficiais acima; proximos ciclos podem aprofundar indicadores numericos
especificos por distrito.

\begin{itemize}
\tightlist
\item
  População municipal/\allowbreak distrital (consolidado operacional): 175427
  habitantes (referencia temporal: 2026-03-05; base: consolidacao de
  fontes oficiais INE\footnote{Instituto Nacional de Estadística (Instituto Nacional de Estatística), órgão oficial de dados demográficos e censitários.} listadas na secao 8).
\item
  Densidade demografica (consolidado operacional): 260 hab/\allowbreak km2
  (referencia temporal: 2026-03-05; base: consolidacao territorial
  oficial).
\item
  Conectividade viaria funcional (consolidado operacional): 80/\allowbreak 100
  (referencia temporal: 2026-03-05; base: MOPC estado de rutas).
\item
  Exposicao hidrometeorologica relativa (consolidado operacional):
  40/\allowbreak 100 (referencia temporal: 2026-03-05; base: DMH/\allowbreak SEN\footnote{Secretaría de Emergencia Nacional (Secretaria de Emergência Nacional), órgão governamental de resposta a desastres.} avisos e
  ocorrencias).
\end{itemize}

NOTA METODOLOGICA: Indicadores consolidados para comparabilidade
distrital nesta fase; quando serie oficial granular não estiver publica,
registrar lacuna oficial e manter rastreabilidade da fonte
institucional. - Cenario curto prazo: interrupcoes logisticas por
eventos climaticos. - Cenario medio prazo: pressao sobre servicos
essenciais. - Mitigacao: redundancia de rota, agua, energia e estoque.

\begin{itemize}
\tightlist
\item
  Cenario otimista (B+1, C+1): GSS 8.3
\item
  Cenario conservador (B-1, C-1): GSS 7.6
\item
  Amplitude de sensibilidade: 0.7 ponto(s)
\item
  Leitura: variacao controlada para decisao comparativa entre distritos;
  priorizar monitoramento de B e C.
\end{itemize}

\subsubsection{Diagnostico Integrado}\label{vuln-General_Resquin-diagnostico-integrado}

General Resquin (San Pedro) apresenta perfil operacional consolidado
para comparacao territorial, com leitura conjunta de infraestrutura,
risco hidroclimatico e capacidade de continuidade de servicos
essenciais.

\subsubsection{Por que sim}\label{vuln-General_Resquin-por-que-sim}

\begin{itemize}
\tightlist
\item
  Base institucional de referencia disponivel e rastreavel.
\item
  Estrutura comparativa (A-E/\allowbreak GSS) consistente para decisao relativa.
\item
  Plano de mitigacao acionavel ja definido no dossie.
\end{itemize}

\subsubsection{Por que não}\label{vuln-General_Resquin-por-que-nuxe3o}

\begin{itemize}
\tightlist
\item
  Serie oficial distrital granular ainda pode apresentar lacunas
  temporais.
\item
  Sensibilidade do score exige monitoramento continuo de risco e
  conectividade.
\item
  Decisao final depende de validacao de campo e janela temporal recente.
\end{itemize}

\subsubsection{Pre-condicoes Minimas de
Decisao}\label{vuln-General_Resquin-pre-condicoes-minimas-de-decisao}

\begin{itemize}
\tightlist
\item
  Atualizar indicadores criticos em ciclo trimestral.
\item
  Confirmar trafegabilidade e contingencia hidrica/\allowbreak energetica local.
\item
  Validar checklist de resiliencia domiciliar/\allowbreak logistica antes de decisao
  final.
\end{itemize}

\subsubsection{Recomendação
Técnica}\label{vuln-General_Resquin-recomendauxe7uxe3o-tuxe9cnica}

Uso recomendado como candidato em analise multicriterio, condicionado ao
cumprimento das pre-condicoes acima. Referencia atual de score: GSS 8.0
em 2026-03-05; risco predominante: baixo.

\begin{itemize}
\tightlist
\item
  \begin{enumerate}
  \def\labelenumi{\arabic{enumi})}
  \tightlist
  \item
    Delimitacao territorial validada por georreferencia institucional
    (Fonte: acesso: 2026-03-05).
  \end{enumerate}
\item
  \begin{enumerate}
  \def\labelenumi{\arabic{enumi})}
  \setcounter{enumi}{1}
  \tightlist
  \item
    População municipal/\allowbreak distrital referenciada por base censitaria
    nacional (Fonte: acesso: 2026-03-05).
  \end{enumerate}
\item
  \begin{enumerate}
  \def\labelenumi{\arabic{enumi})}
  \setcounter{enumi}{2}
  \tightlist
  \item
    Comparabilidade intra-departamental apoiada em indicadores
    distritais oficiais (Fonte: acesso: 2026-03-05).
  \end{enumerate}
\item
  \begin{enumerate}
  \def\labelenumi{\arabic{enumi})}
  \setcounter{enumi}{3}
  \tightlist
  \item
    Estado macro de conectividade viaria ancorado em informacoes de
    rodovias (Fonte: acesso: 2026-03-05).
  \end{enumerate}
\item
  \begin{enumerate}
  \def\labelenumi{\arabic{enumi})}
  \setcounter{enumi}{4}
  \tightlist
  \item
    Exposicao hidrometeorologica monitorada por avisos oficiais (Fonte:
    acesso: 2026-03-05).
  \end{enumerate}
\item
  \begin{enumerate}
  \def\labelenumi{\arabic{enumi})}
  \setcounter{enumi}{5}
  \tightlist
  \item
    Historico de contexto meteorologico apoiado em publicacoes tecnicas
    (Fonte: acesso: 2026-03-05).
  \end{enumerate}
\item
  \begin{enumerate}
  \def\labelenumi{\arabic{enumi})}
  \setcounter{enumi}{6}
  \tightlist
  \item
    Capacidade institucional de resposta a eventos apoiada em acoes da
    SEN\footnote{Secretaría de Emergencia Nacional (Secretaria de Emergência Nacional), órgão governamental de resposta a desastres.} (Fonte: acesso: 2026-03-05).
  \end{enumerate}
\item
  \begin{enumerate}
  \def\labelenumi{\arabic{enumi})}
  \setcounter{enumi}{7}
  \tightlist
  \item
    Infraestrutura energetica analisada via operador nacional (Fonte:
    acesso: 2026-03-05).
  \end{enumerate}
\item
  \begin{enumerate}
  \def\labelenumi{\arabic{enumi})}
  \setcounter{enumi}{8}
  \tightlist
  \item
    Referencias de obras e logistica nacional verificadas no MOPC
    (Fonte: acesso: 2026-03-05).
  \end{enumerate}
\item
  \begin{enumerate}
  \def\labelenumi{\arabic{enumi})}
  \setcounter{enumi}{9}
  \tightlist
  \item
    Contexto sociopolitico institucional referenciado por autoridade
    eleitoral (Fonte: acesso: 2026-03-05).
  \end{enumerate}
\item
  \begin{enumerate}
  \def\labelenumi{\arabic{enumi})}
  \setcounter{enumi}{10}
  \tightlist
  \item
    Camada cartografica de apoio eleitoral/\allowbreak territorial consultada
    (Fonte: acesso: 2026-03-05).
  \end{enumerate}
\item
  \begin{enumerate}
  \def\labelenumi{\arabic{enumi})}
  \setcounter{enumi}{11}
  \tightlist
  \item
    Dados publicos complementares para verificacao cruzada (Fonte:
    acesso: 2026-03-05).
  \end{enumerate}
\end{itemize}

\subsubsection{Combustível}\label{vuln-General_Resquin-combustuxedvel}

\textbf{Referência:} PETROPAR /\allowbreak  postos locais (2024) \textbf{Tipo de
localidade:} Interior

\begin{longtable}[]{@{}lll@{}}
\toprule\noalign{}
Combustível & USD/\allowbreak litro & Gs/\allowbreak litro (aprox.) \\
\midrule\noalign{}
\endhead
\bottomrule\noalign{}
\endlastfoot
Gasolina 93 oct & 1.00 & 7,400 \\
Gasolina 97 oct (premium) & 1.10 & 8,140 \\
Diesel & 0.93 & 6,882 \\
\end{longtable}

\begin{quote}
Preços podem variar ±5\% conforme posto e sazonalidade. Chaco e interior
remoto apresentam maior variação.
\end{quote}

\subsubsection{Cobertura Celular}\label{vuln-General_Resquin-cobertura-celular}

\textbf{Fonte:} CONATEL PY /\allowbreak  operadoras (2024)

\begin{longtable}[]{@{}ll@{}}
\toprule\noalign{}
Parâmetro & Valor \\
\midrule\noalign{}
\endhead
\bottomrule\noalign{}
\endlastfoot
Cobertura 4G população (dept.) & 80\% \\
Cobertura 4G área rural & 58\% \\
Melhor operadora & Tigo \\
Qualidade rural & moderada \\
\end{longtable}

\begin{quote}
Para áreas rurais fora do núcleo urbano, recomenda-se chip Tigo como
principal e Personal como backup.
\end{quote}

\subsubsection{Internet}\label{vuln-General_Resquin-internet}

\textbf{Fonte:} CONATEL /\allowbreak  Speedtest Ookla (2024)

\begin{longtable}[]{@{}ll@{}}
\toprule\noalign{}
Parâmetro & Valor \\
\midrule\noalign{}
\endhead
\bottomrule\noalign{}
\endlastfoot
Velocidade média download & 35 Mbps \\
Domicílios com internet (dept.) & 45\% \\
Tecnologia predominante & rádio \\
Opção rural & Starlink disponível (\textasciitilde USD 44/\allowbreak mês) \\
\end{longtable}

\subsubsection{Mercado Imobiliário e Terra
Rural}\label{vuln-General_Resquin-mercado-imobiliuxe1rio-e-terra-rural}

\textbf{Fonte:} INDERT /\allowbreak  Clasificados.com.py (2024)

\begin{longtable}[]{@{}ll@{}}
\toprule\noalign{}
Tipo & Referência \\
\midrule\noalign{}
\endhead
\bottomrule\noalign{}
\endlastfoot
Terra agrícola alta prod. (USD/\allowbreak ha) & 4,600 \\
Imóvel urbano (USD/\allowbreak m²) & 600 \\
Aluguel 2 quartos (USD/\allowbreak mês) & 250 \\
\end{longtable}

\begin{quote}
Valores de referência departamental. Localidades menores podem ter
preços 20--40\% abaixo da capital departamental. \#\#\# Saúde
\end{quote}

\textbf{Fonte:} MSPBS /\allowbreak  IPS Paraguay (2024-2026), consolidação
departamental e proxy local

\begin{longtable}[]{@{}
  >{\raggedright\arraybackslash}p{(\linewidth - 2\tabcolsep) * \real{0.3600}}
  >{\raggedright\arraybackslash}p{(\linewidth - 2\tabcolsep) * \real{0.6400}}@{}}
\toprule\noalign{}
\begin{minipage}[b]{\linewidth}\raggedright
Serviço
\end{minipage} & \begin{minipage}[b]{\linewidth}\raggedright
Disponibilidade
\end{minipage} \\
\midrule\noalign{}
\endhead
\bottomrule\noalign{}
\endlastfoot
USF /\allowbreak  Posto de Saúde & sim \\
Hospital Regional & não \\
IPS (seguro social) & não \\
Farmácia & sim \\
Distância ao hospital de referência & \textasciitilde2 km (San Pedro de
Ycuamandiyu) \\
\end{longtable}

\textbf{Principais estabelecimentos:} USF local; referencia hospitalar
em San Pedro de Ycuamandiyu

\textbf{Observação para imigrantes:} Atendimento primario local ou em
raio curto; casos de maior complexidade seguem para San Pedro de
Ycuamandiyu. Cobertura privada continua recomendada para especialidades
e urgências de maior complexidade.
\subsubsection{Dossiê de Campo}\label{dossie-General_Resquin-dossiuxea-de-campo}
\begin{description}[style=nextline,leftmargin=0.5cm,font=\bfseries]
\item[Coordenadas]  24°06'S 57°04'W (geocodificado via Nominatim).
\end{description}
\paragraph{Solo (SoilGrids 2.0, média ponderada 0--30
cm)}

\begin{longtable}[]{@{}ll@{}}
\toprule\noalign{}
Parâmetro & Valor \\
\midrule\noalign{}
\endhead
\bottomrule\noalign{}
\endlastfoot
pH (H$_2$O) & 5.7 \\
Carbono orgânico (SOC) & 17.85 g/\allowbreak kg \\
Argila & 26.05 \% \\
Areia & 49.35 \% \\
Silte (calc.) & 24.6 \% \\
Densidade aparente & 1.29 g/\allowbreak cm³ \\
\textbf{Aptidão agrícola} & \textbf{Alta} \\
\end{longtable}

Fonte: ISRIC SoilGrids 2.0 via WCS (acesso em 2026-03-21). Coords:
-24.1°, -57.0667°. Média ponderada camadas 0-5, 5-15, 15-30 cm.
\newpage
\secaoDiagnostico{Guayaibi}{\mapaConteudo{-24.5333}{-56.4000}}{Guayaibi (San Pedro) apresenta perfil operacional consolidado para
comparacao territorial, com leitura conjunta de infraestrutura, risco
hidroclimatico e capacidade de continuidade de servicos essenciais.

\subsubsection{Por que sim}

\begin{itemize}
\tightlist
\item
  Base institucional de referencia disponivel e rastreavel.
\item
  Estrutura comparativa (A-E/\allowbreak GSS) consistente para decisao relativa.
\item
  Plano de mitigacao acionavel ja definido no dossie.
\end{itemize}

\subsubsection{Por que não}

\begin{itemize}
\tightlist
\item
  Serie oficial distrital granular ainda pode apresentar lacunas
  temporais.
\item
  Sensibilidade do score exige monitoramento continuo de risco e
  conectividade.
\item
  Decisao final depende de validacao de campo e janela temporal recente.
\end{itemize}

\subsubsection{Pre-condicoes Minimas de
Decisao}

\begin{itemize}
\tightlist
\item
  Atualizar indicadores criticos em ciclo trimestral.
\item
  Confirmar trafegabilidade e contingencia hidrica/\allowbreak energetica local.
\item
  Validar checklist de resiliencia domiciliar/\allowbreak logistica antes de decisao
  final.
\end{itemize}

\subsubsection{Recomendação
Técnica}

Uso recomendado como candidato em analise multicriterio, condicionado ao
cumprimento das pre-condicoes acima. Referencia atual de score: GSS 7.2
em 2026-03-05; risco predominante: baixo.

\subsubsection{}

\begin{itemize}
\tightlist
\item
  \begin{enumerate}
  \def\labelenumi{\arabic{enumi})}
  \tightlist
  \item
    Delimitacao territorial validada por georreferencia institucional
    (Fonte: acesso: 2026-03-05).
  \end{enumerate}
\item
  \begin{enumerate}
  \def\labelenumi{\arabic{enumi})}
  \setcounter{enumi}{1}
  \tightlist
  \item
    População municipal/\allowbreak distrital referenciada por base censitaria
    nacional (Fonte: acesso: 2026-03-05).
  \end{enumerate}
\item
  \begin{enumerate}
  \def\labelenumi{\arabic{enumi})}
  \setcounter{enumi}{2}
  \tightlist
  \item
    Comparabilidade intra-departamental apoiada em indicadores
    distritais oficiais (Fonte: acesso: 2026-03-05).
  \end{enumerate}
\item
  \begin{enumerate}
  \def\labelenumi{\arabic{enumi})}
  \setcounter{enumi}{3}
  \tightlist
  \item
    Estado macro de conectividade viaria ancorado em informacoes de
    rodovias (Fonte: acesso: 2026-03-05).
  \end{enumerate}
\item
  \begin{enumerate}
  \def\labelenumi{\arabic{enumi})}
  \setcounter{enumi}{4}
  \tightlist
  \item
    Exposicao hidrometeorologica monitorada por avisos oficiais (Fonte:
    acesso: 2026-03-05).
  \end{enumerate}
\item
  \begin{enumerate}
  \def\labelenumi{\arabic{enumi})}
  \setcounter{enumi}{5}
  \tightlist
  \item
    Historico de contexto meteorologico apoiado em publicacoes tecnicas
    (Fonte: acesso: 2026-03-05).
  \end{enumerate}
\item
  \begin{enumerate}
  \def\labelenumi{\arabic{enumi})}
  \setcounter{enumi}{6}
  \tightlist
  \item
    Capacidade institucional de resposta a eventos apoiada em acoes da
    SEN\footnote{Secretaría de Emergencia Nacional (Secretaria de Emergência Nacional), órgão governamental de resposta a desastres.} (Fonte: acesso: 2026-03-05).
  \end{enumerate}
\item
  \begin{enumerate}
  \def\labelenumi{\arabic{enumi})}
  \setcounter{enumi}{7}
  \tightlist
  \item
    Infraestrutura energetica analisada via operador nacional (Fonte:
    acesso: 2026-03-05).
  \end{enumerate}
\item
  \begin{enumerate}
  \def\labelenumi{\arabic{enumi})}
  \setcounter{enumi}{8}
  \tightlist
  \item
    Referencias de obras e logistica nacional verificadas no MOPC
    (Fonte: acesso: 2026-03-05).
  \end{enumerate}
\item
  \begin{enumerate}
  \def\labelenumi{\arabic{enumi})}
  \setcounter{enumi}{9}
  \tightlist
  \item
    Contexto sociopolitico institucional referenciado por autoridade
    eleitoral (Fonte: acesso: 2026-03-05).
  \end{enumerate}
\item
  \begin{enumerate}
  \def\labelenumi{\arabic{enumi})}
  \setcounter{enumi}{10}
  \tightlist
  \item
    Camada cartografica de apoio eleitoral/\allowbreak territorial consultada
    (Fonte: acesso: 2026-03-05).
  \end{enumerate}
\item
  \begin{enumerate}
  \def\labelenumi{\arabic{enumi})}
  \setcounter{enumi}{11}
  \tightlist
  \item
    Dados publicos complementares para verificacao cruzada (Fonte:
    acesso: 2026-03-05).
  \end{enumerate}
\end{itemize}}
\begin{table}[ht!]
\small \begin{tabular}{p{3.0cm}p{0.8cm}p{6.0cm}} \toprule
\textbf{Critério} & \textbf{Nota} & \textbf{Justificativa} \\ \midrule
A - Ameaças Estratégicas & 8.1 & - \\
B - Risco Social & 5.6 & - \\
C - Riscos Naturais & 8.0 & - \\
D - Autossuficiência & 7.6 & - \\
E - Institucional & 6.9 & - \\
\midrule \textbf{GSS FINAL} & \textbf{7.2} & \textbf{Seguro (bom potencial com preparacao).} \\ \bottomrule \end{tabular}
\end{table}
\subsubsection{Vulnerabilidades e
Mitigação}\label{vuln-Guayaibi-vulnerabilidades-e-mitigauxe7uxe3o}

\begin{itemize}
\tightlist
\item
  Exposicao hidrometeorologica controlada (\textless45/\allowbreak 100).
\item
  Conectividade viaria intermediaria (50-69/\allowbreak 100).
\item
  Estoque de contingencia para 14 dias (agua/\allowbreak alimentos/\allowbreak combustivel),
  priorizado quando risco hidro=25/\allowbreak 100.
\item
  Duplicar rotas/\allowbreak logistica quando conectividade viaria=61/\allowbreak 100 for
  \textless70; manter rota secundaria mapeada e testada trimestralmente.
\item
  Plano de energia resiliente com autonomia minima de 72h
  (geracao/\allowbreak backup), revisao semestral.
\end{itemize}

Sintese aplicada: - Dados textuais consolidados com base nas fontes
oficiais acima; proximos ciclos podem aprofundar indicadores numericos
especificos por distrito.

\begin{itemize}
\tightlist
\item
  População municipal/\allowbreak distrital (consolidado operacional): 16517
  habitantes (referencia temporal: 2026-03-05; base: consolidacao de
  fontes oficiais INE\footnote{Instituto Nacional de Estadística (Instituto Nacional de Estatística), órgão oficial de dados demográficos e censitários.} listadas na secao 8).
\item
  Densidade demografica (consolidado operacional): 134 hab/\allowbreak km2
  (referencia temporal: 2026-03-05; base: consolidacao territorial
  oficial).
\item
  Conectividade viaria funcional (consolidado operacional): 61/\allowbreak 100
  (referencia temporal: 2026-03-05; base: MOPC estado de rutas).
\item
  Exposicao hidrometeorologica relativa (consolidado operacional):
  25/\allowbreak 100 (referencia temporal: 2026-03-05; base: DMH/\allowbreak SEN\footnote{Secretaría de Emergencia Nacional (Secretaria de Emergência Nacional), órgão governamental de resposta a desastres.} avisos e
  ocorrencias).
\end{itemize}

NOTA METODOLOGICA: Indicadores consolidados para comparabilidade
distrital nesta fase; quando serie oficial granular não estiver publica,
registrar lacuna oficial e manter rastreabilidade da fonte
institucional. - Cenario curto prazo: interrupcoes logisticas por
eventos climaticos. - Cenario medio prazo: pressao sobre servicos
essenciais. - Mitigacao: redundancia de rota, agua, energia e estoque.

\begin{itemize}
\tightlist
\item
  Cenario otimista (B+1, C+1): GSS 7.6
\item
  Cenario conservador (B-1, C-1): GSS 6.9
\item
  Amplitude de sensibilidade: 0.7 ponto(s)
\item
  Leitura: variacao controlada para decisao comparativa entre distritos;
  priorizar monitoramento de B e C.
\end{itemize}

\subsubsection{Diagnostico Integrado}\label{vuln-Guayaibi-diagnostico-integrado}

Guayaibi (San Pedro) apresenta perfil operacional consolidado para
comparacao territorial, com leitura conjunta de infraestrutura, risco
hidroclimatico e capacidade de continuidade de servicos essenciais.

\subsubsection{Por que sim}\label{vuln-Guayaibi-por-que-sim}

\begin{itemize}
\tightlist
\item
  Base institucional de referencia disponivel e rastreavel.
\item
  Estrutura comparativa (A-E/\allowbreak GSS) consistente para decisao relativa.
\item
  Plano de mitigacao acionavel ja definido no dossie.
\end{itemize}

\subsubsection{Por que não}\label{vuln-Guayaibi-por-que-nuxe3o}

\begin{itemize}
\tightlist
\item
  Serie oficial distrital granular ainda pode apresentar lacunas
  temporais.
\item
  Sensibilidade do score exige monitoramento continuo de risco e
  conectividade.
\item
  Decisao final depende de validacao de campo e janela temporal recente.
\end{itemize}

\subsubsection{Pre-condicoes Minimas de
Decisao}\label{vuln-Guayaibi-pre-condicoes-minimas-de-decisao}

\begin{itemize}
\tightlist
\item
  Atualizar indicadores criticos em ciclo trimestral.
\item
  Confirmar trafegabilidade e contingencia hidrica/\allowbreak energetica local.
\item
  Validar checklist de resiliencia domiciliar/\allowbreak logistica antes de decisao
  final.
\end{itemize}

\subsubsection{Recomendação
Técnica}\label{vuln-Guayaibi-recomendauxe7uxe3o-tuxe9cnica}

Uso recomendado como candidato em analise multicriterio, condicionado ao
cumprimento das pre-condicoes acima. Referencia atual de score: GSS 7.2
em 2026-03-05; risco predominante: baixo.

\begin{itemize}
\tightlist
\item
  \begin{enumerate}
  \def\labelenumi{\arabic{enumi})}
  \tightlist
  \item
    Delimitacao territorial validada por georreferencia institucional
    (Fonte: acesso: 2026-03-05).
  \end{enumerate}
\item
  \begin{enumerate}
  \def\labelenumi{\arabic{enumi})}
  \setcounter{enumi}{1}
  \tightlist
  \item
    População municipal/\allowbreak distrital referenciada por base censitaria
    nacional (Fonte: acesso: 2026-03-05).
  \end{enumerate}
\item
  \begin{enumerate}
  \def\labelenumi{\arabic{enumi})}
  \setcounter{enumi}{2}
  \tightlist
  \item
    Comparabilidade intra-departamental apoiada em indicadores
    distritais oficiais (Fonte: acesso: 2026-03-05).
  \end{enumerate}
\item
  \begin{enumerate}
  \def\labelenumi{\arabic{enumi})}
  \setcounter{enumi}{3}
  \tightlist
  \item
    Estado macro de conectividade viaria ancorado em informacoes de
    rodovias (Fonte: acesso: 2026-03-05).
  \end{enumerate}
\item
  \begin{enumerate}
  \def\labelenumi{\arabic{enumi})}
  \setcounter{enumi}{4}
  \tightlist
  \item
    Exposicao hidrometeorologica monitorada por avisos oficiais (Fonte:
    acesso: 2026-03-05).
  \end{enumerate}
\item
  \begin{enumerate}
  \def\labelenumi{\arabic{enumi})}
  \setcounter{enumi}{5}
  \tightlist
  \item
    Historico de contexto meteorologico apoiado em publicacoes tecnicas
    (Fonte: acesso: 2026-03-05).
  \end{enumerate}
\item
  \begin{enumerate}
  \def\labelenumi{\arabic{enumi})}
  \setcounter{enumi}{6}
  \tightlist
  \item
    Capacidade institucional de resposta a eventos apoiada em acoes da
    SEN\footnote{Secretaría de Emergencia Nacional (Secretaria de Emergência Nacional), órgão governamental de resposta a desastres.} (Fonte: acesso: 2026-03-05).
  \end{enumerate}
\item
  \begin{enumerate}
  \def\labelenumi{\arabic{enumi})}
  \setcounter{enumi}{7}
  \tightlist
  \item
    Infraestrutura energetica analisada via operador nacional (Fonte:
    acesso: 2026-03-05).
  \end{enumerate}
\item
  \begin{enumerate}
  \def\labelenumi{\arabic{enumi})}
  \setcounter{enumi}{8}
  \tightlist
  \item
    Referencias de obras e logistica nacional verificadas no MOPC
    (Fonte: acesso: 2026-03-05).
  \end{enumerate}
\item
  \begin{enumerate}
  \def\labelenumi{\arabic{enumi})}
  \setcounter{enumi}{9}
  \tightlist
  \item
    Contexto sociopolitico institucional referenciado por autoridade
    eleitoral (Fonte: acesso: 2026-03-05).
  \end{enumerate}
\item
  \begin{enumerate}
  \def\labelenumi{\arabic{enumi})}
  \setcounter{enumi}{10}
  \tightlist
  \item
    Camada cartografica de apoio eleitoral/\allowbreak territorial consultada
    (Fonte: acesso: 2026-03-05).
  \end{enumerate}
\item
  \begin{enumerate}
  \def\labelenumi{\arabic{enumi})}
  \setcounter{enumi}{11}
  \tightlist
  \item
    Dados publicos complementares para verificacao cruzada (Fonte:
    acesso: 2026-03-05).
  \end{enumerate}
\end{itemize}

\subsubsection{Combustível}\label{vuln-Guayaibi-combustuxedvel}

\textbf{Referência:} PETROPAR /\allowbreak  postos locais (2024) \textbf{Tipo de
localidade:} Interior

\begin{longtable}[]{@{}lll@{}}
\toprule\noalign{}
Combustível & USD/\allowbreak litro & Gs/\allowbreak litro (aprox.) \\
\midrule\noalign{}
\endhead
\bottomrule\noalign{}
\endlastfoot
Gasolina 93 oct & 1.00 & 7,400 \\
Gasolina 97 oct (premium) & 1.10 & 8,140 \\
Diesel & 0.93 & 6,882 \\
\end{longtable}

\begin{quote}
Preços podem variar ±5\% conforme posto e sazonalidade. Chaco e interior
remoto apresentam maior variação.
\end{quote}

\subsubsection{Cobertura Celular}\label{vuln-Guayaibi-cobertura-celular}

\textbf{Fonte:} CONATEL PY /\allowbreak  operadoras (2024)

\begin{longtable}[]{@{}ll@{}}
\toprule\noalign{}
Parâmetro & Valor \\
\midrule\noalign{}
\endhead
\bottomrule\noalign{}
\endlastfoot
Cobertura 4G população (dept.) & 80\% \\
Cobertura 4G área rural & 58\% \\
Melhor operadora & Tigo \\
Qualidade rural & moderada \\
\end{longtable}

\begin{quote}
Para áreas rurais fora do núcleo urbano, recomenda-se chip Tigo como
principal e Personal como backup.
\end{quote}

\subsubsection{Internet}\label{vuln-Guayaibi-internet}

\textbf{Fonte:} CONATEL /\allowbreak  Speedtest Ookla (2024)

\begin{longtable}[]{@{}ll@{}}
\toprule\noalign{}
Parâmetro & Valor \\
\midrule\noalign{}
\endhead
\bottomrule\noalign{}
\endlastfoot
Velocidade média download & 35 Mbps \\
Domicílios com internet (dept.) & 45\% \\
Tecnologia predominante & rádio \\
Opção rural & Starlink disponível (\textasciitilde USD 44/\allowbreak mês) \\
\end{longtable}

\subsubsection{Mercado Imobiliário e Terra
Rural}\label{vuln-Guayaibi-mercado-imobiliuxe1rio-e-terra-rural}

\textbf{Fonte:} INDERT /\allowbreak  Clasificados.com.py (2024)

\begin{longtable}[]{@{}ll@{}}
\toprule\noalign{}
Tipo & Referência \\
\midrule\noalign{}
\endhead
\bottomrule\noalign{}
\endlastfoot
Terra agrícola alta prod. (USD/\allowbreak ha) & 4,600 \\
Imóvel urbano (USD/\allowbreak m²) & 600 \\
Aluguel 2 quartos (USD/\allowbreak mês) & 250 \\
\end{longtable}

\begin{quote}
Valores de referência departamental. Localidades menores podem ter
preços 20--40\% abaixo da capital departamental. \#\#\# Saúde
\end{quote}

\textbf{Fonte:} MSPBS /\allowbreak  IPS Paraguay (2024-2026), consolidação
departamental e proxy local

\begin{longtable}[]{@{}
  >{\raggedright\arraybackslash}p{(\linewidth - 2\tabcolsep) * \real{0.3600}}
  >{\raggedright\arraybackslash}p{(\linewidth - 2\tabcolsep) * \real{0.6400}}@{}}
\toprule\noalign{}
\begin{minipage}[b]{\linewidth}\raggedright
Serviço
\end{minipage} & \begin{minipage}[b]{\linewidth}\raggedright
Disponibilidade
\end{minipage} \\
\midrule\noalign{}
\endhead
\bottomrule\noalign{}
\endlastfoot
USF /\allowbreak  Posto de Saúde & sim \\
Hospital Regional & não \\
IPS (seguro social) & não \\
Farmácia & sim \\
Distância ao hospital de referência & \textasciitilde106 km (San Pedro
de Ycuamandiyu) \\
\end{longtable}

\textbf{Principais estabelecimentos:} USF local; referencia hospitalar
em San Pedro de Ycuamandiyu

\textbf{Observação para imigrantes:} Atendimento primario local ou em
raio curto; casos de maior complexidade seguem para San Pedro de
Ycuamandiyu. Cobertura privada continua recomendada para especialidades
e urgências de maior complexidade.
\subsubsection{Dossiê de Campo}\label{dossie-Guayaibi-dossiuxea-de-campo}
\begin{description}[style=nextline,leftmargin=0.5cm,font=\bfseries]
\item[Coordenadas]  24°32'S 56°24'W (geocodificado via Nominatim).
\end{description}
\paragraph{Solo (SoilGrids 2.0, média ponderada 0--30
cm)}

\begin{longtable}[]{@{}ll@{}}
\toprule\noalign{}
Parâmetro & Valor \\
\midrule\noalign{}
\endhead
\bottomrule\noalign{}
\endlastfoot
pH (H$_2$O) & 5.3 \\
Carbono orgânico (SOC) & 19.9 g/\allowbreak kg \\
Argila & 25.23 \% \\
Areia & 56.75 \% \\
Silte (calc.) & 18.0 \% \\
Densidade aparente & 1.27 g/\allowbreak cm³ \\
\textbf{Aptidão agrícola} & \textbf{Média} \\
\end{longtable}

Fonte: ISRIC SoilGrids 2.0 via WCS (acesso em 2026-03-21). Coords:
-24.5333°, -56.4°. Média ponderada camadas 0-5, 5-15, 15-30 cm.
\newpage
\secaoDiagnostico{Itacurubi}{\mapaConteudo{-24.5333}{-56.8166}}{Itacurubi (San Pedro) apresenta perfil operacional consolidado para
comparacao territorial, com leitura conjunta de infraestrutura, risco
hidroclimatico e capacidade de continuidade de servicos essenciais.

\subsubsection{Por que sim}

\begin{itemize}
\tightlist
\item
  Base institucional de referencia disponivel e rastreavel.
\item
  Estrutura comparativa (A-E/\allowbreak GSS) consistente para decisao relativa.
\item
  Plano de mitigacao acionavel ja definido no dossie.
\end{itemize}

\subsubsection{Por que não}

\begin{itemize}
\tightlist
\item
  Serie oficial distrital granular ainda pode apresentar lacunas
  temporais.
\item
  Sensibilidade do score exige monitoramento continuo de risco e
  conectividade.
\item
  Decisao final depende de validacao de campo e janela temporal recente.
\end{itemize}

\subsubsection{Pre-condicoes Minimas de
Decisao}

\begin{itemize}
\tightlist
\item
  Atualizar indicadores criticos em ciclo trimestral.
\item
  Confirmar trafegabilidade e contingencia hidrica/\allowbreak energetica local.
\item
  Validar checklist de resiliencia domiciliar/\allowbreak logistica antes de decisao
  final.
\end{itemize}

\subsubsection{Recomendação
Técnica}

Uso recomendado como candidato em analise multicriterio, condicionado ao
cumprimento das pre-condicoes acima. Referencia atual de score: GSS 7.3
em 2026-03-05; risco predominante: baixo.

\subsubsection{}

\begin{itemize}
\tightlist
\item
  \begin{enumerate}
  \def\labelenumi{\arabic{enumi})}
  \tightlist
  \item
    Delimitacao territorial validada por georreferencia institucional
    (Fonte: acesso: 2026-03-05).
  \end{enumerate}
\item
  \begin{enumerate}
  \def\labelenumi{\arabic{enumi})}
  \setcounter{enumi}{1}
  \tightlist
  \item
    População municipal/\allowbreak distrital referenciada por base censitaria
    nacional (Fonte: acesso: 2026-03-05).
  \end{enumerate}
\item
  \begin{enumerate}
  \def\labelenumi{\arabic{enumi})}
  \setcounter{enumi}{2}
  \tightlist
  \item
    Comparabilidade intra-departamental apoiada em indicadores
    distritais oficiais (Fonte: acesso: 2026-03-05).
  \end{enumerate}
\item
  \begin{enumerate}
  \def\labelenumi{\arabic{enumi})}
  \setcounter{enumi}{3}
  \tightlist
  \item
    Estado macro de conectividade viaria ancorado em informacoes de
    rodovias (Fonte: acesso: 2026-03-05).
  \end{enumerate}
\item
  \begin{enumerate}
  \def\labelenumi{\arabic{enumi})}
  \setcounter{enumi}{4}
  \tightlist
  \item
    Exposicao hidrometeorologica monitorada por avisos oficiais (Fonte:
    acesso: 2026-03-05).
  \end{enumerate}
\item
  \begin{enumerate}
  \def\labelenumi{\arabic{enumi})}
  \setcounter{enumi}{5}
  \tightlist
  \item
    Historico de contexto meteorologico apoiado em publicacoes tecnicas
    (Fonte: acesso: 2026-03-05).
  \end{enumerate}
\item
  \begin{enumerate}
  \def\labelenumi{\arabic{enumi})}
  \setcounter{enumi}{6}
  \tightlist
  \item
    Capacidade institucional de resposta a eventos apoiada em acoes da
    SEN\footnote{Secretaría de Emergencia Nacional (Secretaria de Emergência Nacional), órgão governamental de resposta a desastres.} (Fonte: acesso: 2026-03-05).
  \end{enumerate}
\item
  \begin{enumerate}
  \def\labelenumi{\arabic{enumi})}
  \setcounter{enumi}{7}
  \tightlist
  \item
    Infraestrutura energetica analisada via operador nacional (Fonte:
    acesso: 2026-03-05).
  \end{enumerate}
\item
  \begin{enumerate}
  \def\labelenumi{\arabic{enumi})}
  \setcounter{enumi}{8}
  \tightlist
  \item
    Referencias de obras e logistica nacional verificadas no MOPC
    (Fonte: acesso: 2026-03-05).
  \end{enumerate}
\item
  \begin{enumerate}
  \def\labelenumi{\arabic{enumi})}
  \setcounter{enumi}{9}
  \tightlist
  \item
    Contexto sociopolitico institucional referenciado por autoridade
    eleitoral (Fonte: acesso: 2026-03-05).
  \end{enumerate}
\item
  \begin{enumerate}
  \def\labelenumi{\arabic{enumi})}
  \setcounter{enumi}{10}
  \tightlist
  \item
    Camada cartografica de apoio eleitoral/\allowbreak territorial consultada
    (Fonte: acesso: 2026-03-05).
  \end{enumerate}
\item
  \begin{enumerate}
  \def\labelenumi{\arabic{enumi})}
  \setcounter{enumi}{11}
  \tightlist
  \item
    Dados publicos complementares para verificacao cruzada (Fonte:
    acesso: 2026-03-05).
  \end{enumerate}
\end{itemize}}
\begin{table}[ht!]
\small \begin{tabular}{p{3.0cm}p{0.8cm}p{6.0cm}} \toprule
\textbf{Critério} & \textbf{Nota} & \textbf{Justificativa} \\ \midrule
A - Ameaças Estratégicas & 7.7 & - \\
B - Risco Social & 6.9 & - \\
C - Riscos Naturais & 7.0 & - \\
D - Autossuficiência & 7.4 & - \\
E - Institucional & 7.4 & - \\
\midrule \textbf{GSS FINAL} & \textbf{7.3} & \textbf{Seguro (bom potencial com preparacao).} \\ \bottomrule \end{tabular}
\end{table}
\subsubsection{Vulnerabilidades e
Mitigação}\label{vuln-Itacurubi-vulnerabilidades-e-mitigauxe7uxe3o}

\begin{itemize}
\tightlist
\item
  Exposicao hidrometeorologica controlada (\textless45/\allowbreak 100).
\item
  Conectividade viaria intermediaria (50-69/\allowbreak 100).
\item
  Estoque de contingencia para 14 dias (agua/\allowbreak alimentos/\allowbreak combustivel),
  priorizado quando risco hidro=38/\allowbreak 100.
\item
  Duplicar rotas/\allowbreak logistica quando conectividade viaria=61/\allowbreak 100 for
  \textless70; manter rota secundaria mapeada e testada trimestralmente.
\item
  Plano de energia resiliente com autonomia minima de 72h
  (geracao/\allowbreak backup), revisao semestral.
\end{itemize}

Sintese aplicada: - Dados textuais consolidados com base nas fontes
oficiais acima; proximos ciclos podem aprofundar indicadores numericos
especificos por distrito.

\begin{itemize}
\tightlist
\item
  População municipal/\allowbreak distrital (consolidado operacional): 120283
  habitantes (referencia temporal: 2026-03-05; base: consolidacao de
  fontes oficiais INE\footnote{Instituto Nacional de Estadística (Instituto Nacional de Estatística), órgão oficial de dados demográficos e censitários.} listadas na secao 8).
\item
  Densidade demografica (consolidado operacional): 136 hab/\allowbreak km2
  (referencia temporal: 2026-03-05; base: consolidacao territorial
  oficial).
\item
  Conectividade viaria funcional (consolidado operacional): 61/\allowbreak 100
  (referencia temporal: 2026-03-05; base: MOPC estado de rutas).
\item
  Exposicao hidrometeorologica relativa (consolidado operacional):
  38/\allowbreak 100 (referencia temporal: 2026-03-05; base: DMH/\allowbreak SEN\footnote{Secretaría de Emergencia Nacional (Secretaria de Emergência Nacional), órgão governamental de resposta a desastres.} avisos e
  ocorrencias).
\end{itemize}

NOTA METODOLOGICA: Indicadores consolidados para comparabilidade
distrital nesta fase; quando serie oficial granular não estiver publica,
registrar lacuna oficial e manter rastreabilidade da fonte
institucional. - Cenario curto prazo: interrupcoes logisticas por
eventos climaticos. - Cenario medio prazo: pressao sobre servicos
essenciais. - Mitigacao: redundancia de rota, agua, energia e estoque.

\begin{itemize}
\tightlist
\item
  Cenario otimista (B+1, C+1): GSS 7.7
\item
  Cenario conservador (B-1, C-1): GSS 7.0
\item
  Amplitude de sensibilidade: 0.7 ponto(s)
\item
  Leitura: variacao controlada para decisao comparativa entre distritos;
  priorizar monitoramento de B e C.
\end{itemize}

\subsubsection{Diagnostico Integrado}\label{vuln-Itacurubi-diagnostico-integrado}

Itacurubi (San Pedro) apresenta perfil operacional consolidado para
comparacao territorial, com leitura conjunta de infraestrutura, risco
hidroclimatico e capacidade de continuidade de servicos essenciais.

\subsubsection{Por que sim}\label{vuln-Itacurubi-por-que-sim}

\begin{itemize}
\tightlist
\item
  Base institucional de referencia disponivel e rastreavel.
\item
  Estrutura comparativa (A-E/\allowbreak GSS) consistente para decisao relativa.
\item
  Plano de mitigacao acionavel ja definido no dossie.
\end{itemize}

\subsubsection{Por que não}\label{vuln-Itacurubi-por-que-nuxe3o}

\begin{itemize}
\tightlist
\item
  Serie oficial distrital granular ainda pode apresentar lacunas
  temporais.
\item
  Sensibilidade do score exige monitoramento continuo de risco e
  conectividade.
\item
  Decisao final depende de validacao de campo e janela temporal recente.
\end{itemize}

\subsubsection{Pre-condicoes Minimas de
Decisao}\label{vuln-Itacurubi-pre-condicoes-minimas-de-decisao}

\begin{itemize}
\tightlist
\item
  Atualizar indicadores criticos em ciclo trimestral.
\item
  Confirmar trafegabilidade e contingencia hidrica/\allowbreak energetica local.
\item
  Validar checklist de resiliencia domiciliar/\allowbreak logistica antes de decisao
  final.
\end{itemize}

\subsubsection{Recomendação
Técnica}\label{vuln-Itacurubi-recomendauxe7uxe3o-tuxe9cnica}

Uso recomendado como candidato em analise multicriterio, condicionado ao
cumprimento das pre-condicoes acima. Referencia atual de score: GSS 7.3
em 2026-03-05; risco predominante: baixo.

\begin{itemize}
\tightlist
\item
  \begin{enumerate}
  \def\labelenumi{\arabic{enumi})}
  \tightlist
  \item
    Delimitacao territorial validada por georreferencia institucional
    (Fonte: acesso: 2026-03-05).
  \end{enumerate}
\item
  \begin{enumerate}
  \def\labelenumi{\arabic{enumi})}
  \setcounter{enumi}{1}
  \tightlist
  \item
    População municipal/\allowbreak distrital referenciada por base censitaria
    nacional (Fonte: acesso: 2026-03-05).
  \end{enumerate}
\item
  \begin{enumerate}
  \def\labelenumi{\arabic{enumi})}
  \setcounter{enumi}{2}
  \tightlist
  \item
    Comparabilidade intra-departamental apoiada em indicadores
    distritais oficiais (Fonte: acesso: 2026-03-05).
  \end{enumerate}
\item
  \begin{enumerate}
  \def\labelenumi{\arabic{enumi})}
  \setcounter{enumi}{3}
  \tightlist
  \item
    Estado macro de conectividade viaria ancorado em informacoes de
    rodovias (Fonte: acesso: 2026-03-05).
  \end{enumerate}
\item
  \begin{enumerate}
  \def\labelenumi{\arabic{enumi})}
  \setcounter{enumi}{4}
  \tightlist
  \item
    Exposicao hidrometeorologica monitorada por avisos oficiais (Fonte:
    acesso: 2026-03-05).
  \end{enumerate}
\item
  \begin{enumerate}
  \def\labelenumi{\arabic{enumi})}
  \setcounter{enumi}{5}
  \tightlist
  \item
    Historico de contexto meteorologico apoiado em publicacoes tecnicas
    (Fonte: acesso: 2026-03-05).
  \end{enumerate}
\item
  \begin{enumerate}
  \def\labelenumi{\arabic{enumi})}
  \setcounter{enumi}{6}
  \tightlist
  \item
    Capacidade institucional de resposta a eventos apoiada em acoes da
    SEN\footnote{Secretaría de Emergencia Nacional (Secretaria de Emergência Nacional), órgão governamental de resposta a desastres.} (Fonte: acesso: 2026-03-05).
  \end{enumerate}
\item
  \begin{enumerate}
  \def\labelenumi{\arabic{enumi})}
  \setcounter{enumi}{7}
  \tightlist
  \item
    Infraestrutura energetica analisada via operador nacional (Fonte:
    acesso: 2026-03-05).
  \end{enumerate}
\item
  \begin{enumerate}
  \def\labelenumi{\arabic{enumi})}
  \setcounter{enumi}{8}
  \tightlist
  \item
    Referencias de obras e logistica nacional verificadas no MOPC
    (Fonte: acesso: 2026-03-05).
  \end{enumerate}
\item
  \begin{enumerate}
  \def\labelenumi{\arabic{enumi})}
  \setcounter{enumi}{9}
  \tightlist
  \item
    Contexto sociopolitico institucional referenciado por autoridade
    eleitoral (Fonte: acesso: 2026-03-05).
  \end{enumerate}
\item
  \begin{enumerate}
  \def\labelenumi{\arabic{enumi})}
  \setcounter{enumi}{10}
  \tightlist
  \item
    Camada cartografica de apoio eleitoral/\allowbreak territorial consultada
    (Fonte: acesso: 2026-03-05).
  \end{enumerate}
\item
  \begin{enumerate}
  \def\labelenumi{\arabic{enumi})}
  \setcounter{enumi}{11}
  \tightlist
  \item
    Dados publicos complementares para verificacao cruzada (Fonte:
    acesso: 2026-03-05).
  \end{enumerate}
\end{itemize}

\subsubsection{Combustível}\label{vuln-Itacurubi-combustuxedvel}

\textbf{Referência:} PETROPAR /\allowbreak  postos locais (2024) \textbf{Tipo de
localidade:} Interior

\begin{longtable}[]{@{}lll@{}}
\toprule\noalign{}
Combustível & USD/\allowbreak litro & Gs/\allowbreak litro (aprox.) \\
\midrule\noalign{}
\endhead
\bottomrule\noalign{}
\endlastfoot
Gasolina 93 oct & 1.00 & 7,400 \\
Gasolina 97 oct (premium) & 1.10 & 8,140 \\
Diesel & 0.93 & 6,882 \\
\end{longtable}

\begin{quote}
Preços podem variar ±5\% conforme posto e sazonalidade. Chaco e interior
remoto apresentam maior variação.
\end{quote}

\subsubsection{Cobertura Celular}\label{vuln-Itacurubi-cobertura-celular}

\textbf{Fonte:} CONATEL PY /\allowbreak  operadoras (2024)

\begin{longtable}[]{@{}ll@{}}
\toprule\noalign{}
Parâmetro & Valor \\
\midrule\noalign{}
\endhead
\bottomrule\noalign{}
\endlastfoot
Cobertura 4G população (dept.) & 80\% \\
Cobertura 4G área rural & 58\% \\
Melhor operadora & Tigo \\
Qualidade rural & moderada \\
\end{longtable}

\begin{quote}
Para áreas rurais fora do núcleo urbano, recomenda-se chip Tigo como
principal e Personal como backup.
\end{quote}

\subsubsection{Internet}\label{vuln-Itacurubi-internet}

\textbf{Fonte:} CONATEL /\allowbreak  Speedtest Ookla (2024)

\begin{longtable}[]{@{}ll@{}}
\toprule\noalign{}
Parâmetro & Valor \\
\midrule\noalign{}
\endhead
\bottomrule\noalign{}
\endlastfoot
Velocidade média download & 35 Mbps \\
Domicílios com internet (dept.) & 45\% \\
Tecnologia predominante & rádio \\
Opção rural & Starlink disponível (\textasciitilde USD 44/\allowbreak mês) \\
\end{longtable}

\subsubsection{Mercado Imobiliário e Terra
Rural}\label{vuln-Itacurubi-mercado-imobiliuxe1rio-e-terra-rural}

\textbf{Fonte:} INDERT /\allowbreak  Clasificados.com.py (2024)

\begin{longtable}[]{@{}ll@{}}
\toprule\noalign{}
Tipo & Referência \\
\midrule\noalign{}
\endhead
\bottomrule\noalign{}
\endlastfoot
Terra agrícola alta prod. (USD/\allowbreak ha) & 4,600 \\
Imóvel urbano (USD/\allowbreak m²) & 600 \\
Aluguel 2 quartos (USD/\allowbreak mês) & 250 \\
\end{longtable}

\begin{quote}
Valores de referência departamental. Localidades menores podem ter
preços 20--40\% abaixo da capital departamental. \#\#\# Saúde
\end{quote}

\textbf{Fonte:} MSPBS /\allowbreak  IPS Paraguay (2024-2026), consolidação
departamental e proxy local

\begin{longtable}[]{@{}
  >{\raggedright\arraybackslash}p{(\linewidth - 2\tabcolsep) * \real{0.3600}}
  >{\raggedright\arraybackslash}p{(\linewidth - 2\tabcolsep) * \real{0.6400}}@{}}
\toprule\noalign{}
\begin{minipage}[b]{\linewidth}\raggedright
Serviço
\end{minipage} & \begin{minipage}[b]{\linewidth}\raggedright
Disponibilidade
\end{minipage} \\
\midrule\noalign{}
\endhead
\bottomrule\noalign{}
\endlastfoot
USF /\allowbreak  Posto de Saúde & sim \\
Hospital Regional & não \\
IPS (seguro social) & não \\
Farmácia & sim \\
Distância ao hospital de referência & \textasciitilde71 km (San Pedro de
Ycuamandiyu) \\
\end{longtable}

\textbf{Principais estabelecimentos:} USF local; referencia hospitalar
em San Pedro de Ycuamandiyu

\textbf{Observação para imigrantes:} Atendimento primario local ou em
raio curto; casos de maior complexidade seguem para San Pedro de
Ycuamandiyu. Cobertura privada continua recomendada para especialidades
e urgências de maior complexidade.
\subsubsection{Dossiê de Campo}\label{dossie-Itacurubi-dossiuxea-de-campo}
\begin{description}[style=nextline,leftmargin=0.5cm,font=\bfseries]
\item[Coordenadas]  24°32'S 56°49'W (geocodificado via Nominatim).
\end{description}
\paragraph{Solo (SoilGrids 2.0, média ponderada 0--30
cm)}

\begin{longtable}[]{@{}ll@{}}
\toprule\noalign{}
Parâmetro & Valor \\
\midrule\noalign{}
\endhead
\bottomrule\noalign{}
\endlastfoot
pH (H$_2$O) & 5.4 \\
Carbono orgânico (SOC) & 15.82 g/\allowbreak kg \\
Argila & 26.52 \% \\
Areia & 55.02 \% \\
Silte (calc.) & 18.5 \% \\
Densidade aparente & 1.34 g/\allowbreak cm³ \\
\textbf{Aptidão agrícola} & \textbf{Média} \\
\end{longtable}

Fonte: ISRIC SoilGrids 2.0 via WCS (acesso em 2026-03-21). Coords:
-24.5333°, -56.8167°. Média ponderada camadas 0-5, 5-15, 15-30 cm.
\newpage
\secaoDiagnostico{Itacurubi del Rosario}{\mapaConteudo{-24.5333}{-56.8166}}{Itacurubi del Rosario (San Pedro) apresenta perfil operacional
consolidado para comparacao territorial, com leitura conjunta de
infraestrutura, risco hidroclimatico e capacidade de continuidade de
servicos essenciais.

\subsubsection{Por que sim}

\begin{itemize}
\tightlist
\item
  Base institucional de referencia disponivel e rastreavel.
\item
  Estrutura comparativa (A-E/\allowbreak GSS) consistente para decisao relativa.
\item
  Plano de mitigacao acionavel ja definido no dossie.
\end{itemize}

\subsubsection{Por que não}

\begin{itemize}
\tightlist
\item
  Serie oficial distrital granular ainda pode apresentar lacunas
  temporais.
\item
  Sensibilidade do score exige monitoramento continuo de risco e
  conectividade.
\item
  Decisao final depende de validacao de campo e janela temporal recente.
\end{itemize}

\subsubsection{Pre-condicoes Minimas de
Decisao}

\begin{itemize}
\tightlist
\item
  Atualizar indicadores criticos em ciclo trimestral.
\item
  Confirmar trafegabilidade e contingencia hidrica/\allowbreak energetica local.
\item
  Validar checklist de resiliencia domiciliar/\allowbreak logistica antes de decisao
  final.
\end{itemize}

\subsubsection{Recomendação
Técnica}

Uso recomendado como candidato em analise multicriterio, condicionado ao
cumprimento das pre-condicoes acima. Referencia atual de score: GSS 8.0
em 2026-03-05; risco predominante: baixo.

\subsubsection{}

\begin{itemize}
\tightlist
\item
  \begin{enumerate}
  \def\labelenumi{\arabic{enumi})}
  \tightlist
  \item
    Delimitacao territorial validada por georreferencia institucional
    (Fonte: acesso: 2026-03-05).
  \end{enumerate}
\item
  \begin{enumerate}
  \def\labelenumi{\arabic{enumi})}
  \setcounter{enumi}{1}
  \tightlist
  \item
    População municipal/\allowbreak distrital referenciada por base censitaria
    nacional (Fonte: acesso: 2026-03-05).
  \end{enumerate}
\item
  \begin{enumerate}
  \def\labelenumi{\arabic{enumi})}
  \setcounter{enumi}{2}
  \tightlist
  \item
    Comparabilidade intra-departamental apoiada em indicadores
    distritais oficiais (Fonte: acesso: 2026-03-05).
  \end{enumerate}
\item
  \begin{enumerate}
  \def\labelenumi{\arabic{enumi})}
  \setcounter{enumi}{3}
  \tightlist
  \item
    Estado macro de conectividade viaria ancorado em informacoes de
    rodovias (Fonte: acesso: 2026-03-05).
  \end{enumerate}
\item
  \begin{enumerate}
  \def\labelenumi{\arabic{enumi})}
  \setcounter{enumi}{4}
  \tightlist
  \item
    Exposicao hidrometeorologica monitorada por avisos oficiais (Fonte:
    acesso: 2026-03-05).
  \end{enumerate}
\item
  \begin{enumerate}
  \def\labelenumi{\arabic{enumi})}
  \setcounter{enumi}{5}
  \tightlist
  \item
    Historico de contexto meteorologico apoiado em publicacoes tecnicas
    (Fonte: acesso: 2026-03-05).
  \end{enumerate}
\item
  \begin{enumerate}
  \def\labelenumi{\arabic{enumi})}
  \setcounter{enumi}{6}
  \tightlist
  \item
    Capacidade institucional de resposta a eventos apoiada em acoes da
    SEN\footnote{Secretaría de Emergencia Nacional (Secretaria de Emergência Nacional), órgão governamental de resposta a desastres.} (Fonte: acesso: 2026-03-05).
  \end{enumerate}
\item
  \begin{enumerate}
  \def\labelenumi{\arabic{enumi})}
  \setcounter{enumi}{7}
  \tightlist
  \item
    Infraestrutura energetica analisada via operador nacional (Fonte:
    acesso: 2026-03-05).
  \end{enumerate}
\item
  \begin{enumerate}
  \def\labelenumi{\arabic{enumi})}
  \setcounter{enumi}{8}
  \tightlist
  \item
    Referencias de obras e logistica nacional verificadas no MOPC
    (Fonte: acesso: 2026-03-05).
  \end{enumerate}
\item
  \begin{enumerate}
  \def\labelenumi{\arabic{enumi})}
  \setcounter{enumi}{9}
  \tightlist
  \item
    Contexto sociopolitico institucional referenciado por autoridade
    eleitoral (Fonte: acesso: 2026-03-05).
  \end{enumerate}
\item
  \begin{enumerate}
  \def\labelenumi{\arabic{enumi})}
  \setcounter{enumi}{10}
  \tightlist
  \item
    Camada cartografica de apoio eleitoral/\allowbreak territorial consultada
    (Fonte: acesso: 2026-03-05).
  \end{enumerate}
\item
  \begin{enumerate}
  \def\labelenumi{\arabic{enumi})}
  \setcounter{enumi}{11}
  \tightlist
  \item
    Dados publicos complementares para verificacao cruzada (Fonte:
    acesso: 2026-03-05).
  \end{enumerate}
\end{itemize}}
\begin{table}[ht!]
\small \begin{tabular}{p{3.0cm}p{0.8cm}p{6.0cm}} \toprule
\textbf{Critério} & \textbf{Nota} & \textbf{Justificativa} \\ \midrule
A - Ameaças Estratégicas & 9.0 & - \\
B - Risco Social & 7.3 & - \\
C - Riscos Naturais & 8.2 & - \\
D - Autossuficiência & 7.7 & - \\
E - Institucional & 7.8 & - \\
\midrule \textbf{GSS FINAL} & \textbf{8.0} & \textbf{Seguro (bom potencial com preparacao).} \\ \bottomrule \end{tabular}
\end{table}
\subsubsection{Vulnerabilidades e
Mitigação}\label{vuln-Itacurubi_del_Rosario-vulnerabilidades-e-mitigauxe7uxe3o}

\begin{itemize}
\tightlist
\item
  Exposicao hidrometeorologica controlada (\textless45/\allowbreak 100).
\item
  Conectividade viaria favoravel (\textgreater=70/\allowbreak 100).
\item
  Estoque de contingencia para 14 dias (agua/\allowbreak alimentos/\allowbreak combustivel),
  priorizado quando risco hidro=23/\allowbreak 100.
\item
  Duplicar rotas/\allowbreak logistica quando conectividade viaria=88/\allowbreak 100 for
  \textless70; manter rota secundaria mapeada e testada trimestralmente.
\item
  Plano de energia resiliente com autonomia minima de 72h
  (geracao/\allowbreak backup), revisao semestral.
\end{itemize}

Sintese aplicada: - Dados textuais consolidados com base nas fontes
oficiais acima; proximos ciclos podem aprofundar indicadores numericos
especificos por distrito.

\begin{itemize}
\tightlist
\item
  População municipal/\allowbreak distrital (consolidado operacional): 66854
  habitantes (referencia temporal: 2026-03-05; base: consolidacao de
  fontes oficiais INE\footnote{Instituto Nacional de Estadística (Instituto Nacional de Estatística), órgão oficial de dados demográficos e censitários.} listadas na secao 8).
\item
  Densidade demografica (consolidado operacional): 367 hab/\allowbreak km2
  (referencia temporal: 2026-03-05; base: consolidacao territorial
  oficial).
\item
  Conectividade viaria funcional (consolidado operacional): 88/\allowbreak 100
  (referencia temporal: 2026-03-05; base: MOPC estado de rutas).
\item
  Exposicao hidrometeorologica relativa (consolidado operacional):
  23/\allowbreak 100 (referencia temporal: 2026-03-05; base: DMH/\allowbreak SEN\footnote{Secretaría de Emergencia Nacional (Secretaria de Emergência Nacional), órgão governamental de resposta a desastres.} avisos e
  ocorrencias).
\end{itemize}

NOTA METODOLOGICA: Indicadores consolidados para comparabilidade
distrital nesta fase; quando serie oficial granular não estiver publica,
registrar lacuna oficial e manter rastreabilidade da fonte
institucional. - Cenario curto prazo: interrupcoes logisticas por
eventos climaticos. - Cenario medio prazo: pressao sobre servicos
essenciais. - Mitigacao: redundancia de rota, agua, energia e estoque.

\begin{itemize}
\tightlist
\item
  Cenario otimista (B+1, C+1): GSS 8.4
\item
  Cenario conservador (B-1, C-1): GSS 7.7
\item
  Amplitude de sensibilidade: 0.7 ponto(s)
\item
  Leitura: variacao controlada para decisao comparativa entre distritos;
  priorizar monitoramento de B e C.
\end{itemize}

\subsubsection{Diagnostico Integrado}\label{vuln-Itacurubi_del_Rosario-diagnostico-integrado}

Itacurubi del Rosario (San Pedro) apresenta perfil operacional
consolidado para comparacao territorial, com leitura conjunta de
infraestrutura, risco hidroclimatico e capacidade de continuidade de
servicos essenciais.

\subsubsection{Por que sim}\label{vuln-Itacurubi_del_Rosario-por-que-sim}

\begin{itemize}
\tightlist
\item
  Base institucional de referencia disponivel e rastreavel.
\item
  Estrutura comparativa (A-E/\allowbreak GSS) consistente para decisao relativa.
\item
  Plano de mitigacao acionavel ja definido no dossie.
\end{itemize}

\subsubsection{Por que não}\label{vuln-Itacurubi_del_Rosario-por-que-nuxe3o}

\begin{itemize}
\tightlist
\item
  Serie oficial distrital granular ainda pode apresentar lacunas
  temporais.
\item
  Sensibilidade do score exige monitoramento continuo de risco e
  conectividade.
\item
  Decisao final depende de validacao de campo e janela temporal recente.
\end{itemize}

\subsubsection{Pre-condicoes Minimas de
Decisao}\label{vuln-Itacurubi_del_Rosario-pre-condicoes-minimas-de-decisao}

\begin{itemize}
\tightlist
\item
  Atualizar indicadores criticos em ciclo trimestral.
\item
  Confirmar trafegabilidade e contingencia hidrica/\allowbreak energetica local.
\item
  Validar checklist de resiliencia domiciliar/\allowbreak logistica antes de decisao
  final.
\end{itemize}

\subsubsection{Recomendação
Técnica}\label{vuln-Itacurubi_del_Rosario-recomendauxe7uxe3o-tuxe9cnica}

Uso recomendado como candidato em analise multicriterio, condicionado ao
cumprimento das pre-condicoes acima. Referencia atual de score: GSS 8.0
em 2026-03-05; risco predominante: baixo.

\begin{itemize}
\tightlist
\item
  \begin{enumerate}
  \def\labelenumi{\arabic{enumi})}
  \tightlist
  \item
    Delimitacao territorial validada por georreferencia institucional
    (Fonte: acesso: 2026-03-05).
  \end{enumerate}
\item
  \begin{enumerate}
  \def\labelenumi{\arabic{enumi})}
  \setcounter{enumi}{1}
  \tightlist
  \item
    População municipal/\allowbreak distrital referenciada por base censitaria
    nacional (Fonte: acesso: 2026-03-05).
  \end{enumerate}
\item
  \begin{enumerate}
  \def\labelenumi{\arabic{enumi})}
  \setcounter{enumi}{2}
  \tightlist
  \item
    Comparabilidade intra-departamental apoiada em indicadores
    distritais oficiais (Fonte: acesso: 2026-03-05).
  \end{enumerate}
\item
  \begin{enumerate}
  \def\labelenumi{\arabic{enumi})}
  \setcounter{enumi}{3}
  \tightlist
  \item
    Estado macro de conectividade viaria ancorado em informacoes de
    rodovias (Fonte: acesso: 2026-03-05).
  \end{enumerate}
\item
  \begin{enumerate}
  \def\labelenumi{\arabic{enumi})}
  \setcounter{enumi}{4}
  \tightlist
  \item
    Exposicao hidrometeorologica monitorada por avisos oficiais (Fonte:
    acesso: 2026-03-05).
  \end{enumerate}
\item
  \begin{enumerate}
  \def\labelenumi{\arabic{enumi})}
  \setcounter{enumi}{5}
  \tightlist
  \item
    Historico de contexto meteorologico apoiado em publicacoes tecnicas
    (Fonte: acesso: 2026-03-05).
  \end{enumerate}
\item
  \begin{enumerate}
  \def\labelenumi{\arabic{enumi})}
  \setcounter{enumi}{6}
  \tightlist
  \item
    Capacidade institucional de resposta a eventos apoiada em acoes da
    SEN\footnote{Secretaría de Emergencia Nacional (Secretaria de Emergência Nacional), órgão governamental de resposta a desastres.} (Fonte: acesso: 2026-03-05).
  \end{enumerate}
\item
  \begin{enumerate}
  \def\labelenumi{\arabic{enumi})}
  \setcounter{enumi}{7}
  \tightlist
  \item
    Infraestrutura energetica analisada via operador nacional (Fonte:
    acesso: 2026-03-05).
  \end{enumerate}
\item
  \begin{enumerate}
  \def\labelenumi{\arabic{enumi})}
  \setcounter{enumi}{8}
  \tightlist
  \item
    Referencias de obras e logistica nacional verificadas no MOPC
    (Fonte: acesso: 2026-03-05).
  \end{enumerate}
\item
  \begin{enumerate}
  \def\labelenumi{\arabic{enumi})}
  \setcounter{enumi}{9}
  \tightlist
  \item
    Contexto sociopolitico institucional referenciado por autoridade
    eleitoral (Fonte: acesso: 2026-03-05).
  \end{enumerate}
\item
  \begin{enumerate}
  \def\labelenumi{\arabic{enumi})}
  \setcounter{enumi}{10}
  \tightlist
  \item
    Camada cartografica de apoio eleitoral/\allowbreak territorial consultada
    (Fonte: acesso: 2026-03-05).
  \end{enumerate}
\item
  \begin{enumerate}
  \def\labelenumi{\arabic{enumi})}
  \setcounter{enumi}{11}
  \tightlist
  \item
    Dados publicos complementares para verificacao cruzada (Fonte:
    acesso: 2026-03-05).
  \end{enumerate}
\end{itemize}

\subsubsection{Combustível}\label{vuln-Itacurubi_del_Rosario-combustuxedvel}

\textbf{Referência:} PETROPAR /\allowbreak  postos locais (2024) \textbf{Tipo de
localidade:} Interior

\begin{longtable}[]{@{}lll@{}}
\toprule\noalign{}
Combustível & USD/\allowbreak litro & Gs/\allowbreak litro (aprox.) \\
\midrule\noalign{}
\endhead
\bottomrule\noalign{}
\endlastfoot
Gasolina 93 oct & 1.00 & 7,400 \\
Gasolina 97 oct (premium) & 1.10 & 8,140 \\
Diesel & 0.93 & 6,882 \\
\end{longtable}

\begin{quote}
Preços podem variar ±5\% conforme posto e sazonalidade. Chaco e interior
remoto apresentam maior variação.
\end{quote}

\subsubsection{Cobertura Celular}\label{vuln-Itacurubi_del_Rosario-cobertura-celular}

\textbf{Fonte:} CONATEL PY /\allowbreak  operadoras (2024)

\begin{longtable}[]{@{}ll@{}}
\toprule\noalign{}
Parâmetro & Valor \\
\midrule\noalign{}
\endhead
\bottomrule\noalign{}
\endlastfoot
Cobertura 4G população (dept.) & 80\% \\
Cobertura 4G área rural & 58\% \\
Melhor operadora & Tigo \\
Qualidade rural & moderada \\
\end{longtable}

\begin{quote}
Para áreas rurais fora do núcleo urbano, recomenda-se chip Tigo como
principal e Personal como backup.
\end{quote}

\subsubsection{Internet}\label{vuln-Itacurubi_del_Rosario-internet}

\textbf{Fonte:} CONATEL /\allowbreak  Speedtest Ookla (2024)

\begin{longtable}[]{@{}ll@{}}
\toprule\noalign{}
Parâmetro & Valor \\
\midrule\noalign{}
\endhead
\bottomrule\noalign{}
\endlastfoot
Velocidade média download & 35 Mbps \\
Domicílios com internet (dept.) & 45\% \\
Tecnologia predominante & rádio \\
Opção rural & Starlink disponível (\textasciitilde USD 44/\allowbreak mês) \\
\end{longtable}

\subsubsection{Mercado Imobiliário e Terra
Rural}\label{vuln-Itacurubi_del_Rosario-mercado-imobiliuxe1rio-e-terra-rural}

\textbf{Fonte:} INDERT /\allowbreak  Clasificados.com.py (2024)

\begin{longtable}[]{@{}ll@{}}
\toprule\noalign{}
Tipo & Referência \\
\midrule\noalign{}
\endhead
\bottomrule\noalign{}
\endlastfoot
Terra agrícola alta prod. (USD/\allowbreak ha) & 4,600 \\
Imóvel urbano (USD/\allowbreak m²) & 600 \\
Aluguel 2 quartos (USD/\allowbreak mês) & 250 \\
\end{longtable}

\begin{quote}
Valores de referência departamental. Localidades menores podem ter
preços 20--40\% abaixo da capital departamental. \#\#\# Saúde
\end{quote}

\textbf{Fonte:} MSPBS /\allowbreak  IPS Paraguay (2024-2026), consolidação
departamental e proxy local

\begin{longtable}[]{@{}
  >{\raggedright\arraybackslash}p{(\linewidth - 2\tabcolsep) * \real{0.3600}}
  >{\raggedright\arraybackslash}p{(\linewidth - 2\tabcolsep) * \real{0.6400}}@{}}
\toprule\noalign{}
\begin{minipage}[b]{\linewidth}\raggedright
Serviço
\end{minipage} & \begin{minipage}[b]{\linewidth}\raggedright
Disponibilidade
\end{minipage} \\
\midrule\noalign{}
\endhead
\bottomrule\noalign{}
\endlastfoot
USF /\allowbreak  Posto de Saúde & sim \\
Hospital Regional & não \\
IPS (seguro social) & não \\
Farmácia & sim \\
Distância ao hospital de referência & \textasciitilde71 km (San Pedro de
Ycuamandiyu) \\
\end{longtable}

\textbf{Principais estabelecimentos:} USF local; referencia hospitalar
em San Pedro de Ycuamandiyu

\textbf{Observação para imigrantes:} Atendimento primario local ou em
raio curto; casos de maior complexidade seguem para San Pedro de
Ycuamandiyu. Cobertura privada continua recomendada para especialidades
e urgências de maior complexidade.
\subsubsection{Dossiê de Campo}\label{dossie-Itacurubi_del_Rosario-dossiuxea-de-campo}
\begin{description}[style=nextline,leftmargin=0.5cm,font=\bfseries]
\item[Coordenadas]  24°32'S 56°49'W (geocodificado via Nominatim).
\end{description}
\paragraph{Solo (SoilGrids 2.0, média ponderada 0--30
cm)}

\begin{longtable}[]{@{}ll@{}}
\toprule\noalign{}
Parâmetro & Valor \\
\midrule\noalign{}
\endhead
\bottomrule\noalign{}
\endlastfoot
pH (H$_2$O) & 5.4 \\
Carbono orgânico (SOC) & 15.82 g/\allowbreak kg \\
Argila & 26.52 \% \\
Areia & 55.02 \% \\
Silte (calc.) & 18.5 \% \\
Densidade aparente & 1.34 g/\allowbreak cm³ \\
\textbf{Aptidão agrícola} & \textbf{Média} \\
\end{longtable}

Fonte: ISRIC SoilGrids 2.0 via WCS (acesso em 2026-03-21). Coords:
-24.5333°, -56.8167°. Média ponderada camadas 0-5, 5-15, 15-30 cm.
\newpage
\secaoDiagnostico{Liberacion}{\mapaConteudo{-24.3500}{-56.5000}}{Liberación é um nó logístico central em San Pedro. Oferece conectividade
rodoviária ímpar e solo apto para uma grande diversidade de culturas
(grãos, chia e frutas). Sua força reside na logística e abundância de
recursos, enquanto sua principal fraqueza é a alta visibilidade por
estar no encontro de duas das principais rodovias do país.

\subsubsection{Por que sim}

\begin{itemize}
\tightlist
\item
  Conectividade estratégica facilitada para todas as regiões do
  Paraguai.
\item
  Solo versátil e clima favorável para autossuficiência e excedente
  comercial.
\item
  Liberdade legal completa para estrangeiros em propriedades privadas
  tituladas.
\end{itemize}

\subsubsection{Por que não}

\begin{itemize}
\tightlist
\item
  Alvo logístico óbvio em cenários de instabilidade nacional.
\item
  Exposição social elevada no núcleo urbano do Cruce.
\item
  Dependência de serviços hospitalares especializados em cidades
  vizinhas.
\end{itemize}

\subsubsection{Recomendação
Técnica}

Indicado para bases produtivas que requeiram facilidade de escoamento e
acesso a suprimentos diversificados. Exige um planejamento de segurança
que considere o recuo em relação às rodovias principais. Referencia
atual de score: GSS 6.3 em 2026-03-05.

\subsubsection{}

\begin{itemize}
\tightlist
\item
  \begin{enumerate}
  \def\labelenumi{\arabic{enumi})}
  \tightlist
  \item
    Dados censitários (INE\footnote{Instituto Nacional de Estadística (Instituto Nacional de Estatística), órgão oficial de dados demográficos e censitários.} 2022).
  \end{enumerate}
\item
  \begin{enumerate}
  \def\labelenumi{\arabic{enumi})}
  \setcounter{enumi}{1}
  \tightlist
  \item
    Registro de produção agrícola e óleos essenciais (MAG).
  \end{enumerate}
\item
  \begin{enumerate}
  \def\labelenumi{\arabic{enumi})}
  \setcounter{enumi}{2}
  \tightlist
  \item
    Mapa de conectividade e fluxos rodoviários (MOPC 2026).
  \end{enumerate}
\item
  \begin{enumerate}
  \def\labelenumi{\arabic{enumi})}
  \setcounter{enumi}{3}
  \tightlist
  \item
    Jurisprudência sobre terras INDERT vs Títulos Privados.
  \end{enumerate}
\end{itemize}}
\begin{table}[ht!]
\small \begin{tabular}{p{3.0cm}p{0.8cm}p{6.0cm}} \toprule
\textbf{Critério} & \textbf{Nota} & \textbf{Justificativa} \\ \midrule
A - Ameaças Estratégicas & 5.0 & Entroncamento rodoviário vital; alvo tático e logístico de alta visibilidade \\
B - Risco Social & 5.5 & População em crescimento e alta exposição em rotas nacionais de trânsito \\
C - Riscos Naturais & 7.0 & Geologia muito estável e baixo risco de grandes desastres naturais \\
D - Autossuficiência & 7.5 & Recursos diversificados: soja, grãos especiais, citricultura e água abundante \\
E - Institucional & 7.0 & Ambiente institucional funcional; total liberdade jurídica para brasileiros em terras privadas \\
\midrule \textbf{GSS FINAL} & \textbf{6.3} & \textbf{Moderadamente Seguro (Ideal para quem busca base logística e produtividade agrícola diversificada).} \\ \bottomrule \end{tabular}
\end{table}
\subsubsection{Vulnerabilidades e
Mitigação}\label{vuln-Liberacion-vulnerabilidades-e-mitigauxe7uxe3o}

\begin{itemize}
\tightlist
\item
  \textbf{Exposição Social:} O Cruce Liberación atrai fluxos
  populacionais intensos em crises (Mitigação: residência rural afastada
  do eixo asfáltico principal).
\item
  \textbf{Dependência Elétrica:} Instabilidade em áreas rurais
  (Mitigação: sistemas de energia solar redundantes).
\item
  \textbf{Segurança da Terra:} Aquisição de áreas destinadas à reforma
  agrária (INDERT) é proibida para estrangeiros (Mitigação: verificação
  rigorosa de títulos de propriedade privada antes da compra).
\end{itemize}
\subsubsection{Dossiê de Campo}\label{dossie-Liberacion-dossiuxea-de-campo}
\begin{description}[style=nextline,leftmargin=0.5cm,font=\bfseries]
\item[Coordenadas]  24°21'S 56°30'W.
\item[Topografia]  Relevo predominantemente plano. Localizado em um ponto de convergência logística no centro de San Pedro. Geologicamente estável, risco sísmico muito baixo.
\item[Alvos Estratégicos]  Cruce Liberación (Entroncamento vital das Rutas PY03 e PY08); Hub logístico para o transporte de grãos e gado do Norte em direção aos portos do Sul e Leste.
\item[Fallout]  Ventos predominantes do Norte e Nordeste (NE), quentes e úmidos. Baixo risco de fallout direto de alvos continentais primários.
\item[População]  \textasciitilde 19.100 habitantes (Censo 2022 Final). Perfil urbano-rural dinâmico concentrado no Cruce Liberación.
\item[Densidade]  \textasciitilde 32,5 hab/\allowbreak km² (Moderada, com alta concentração comercial ao longo do eixo rodoviário).
\item[Vias de Acesso]  Ponto de encontro da Ruta PY03 (Assunção-Norte) e PY08 (Eixo Norte-Sul). Conectividade asfáltica de excelência em tempos de paz, mas risco elevado de controle/\allowbreak bloqueio em crises.
\item[Serviços]  Centro de Saúde "Doroteo Colmán Navarro" (modernizado recentemente). Dependência de San Estanislao (Santaní) ou Santa Rosa del Aguaray para atendimentos de alta complexidade.
\item[Custo de Vida]  Baixo; focado em produtos de agricultura familiar local e serviços rodoviários.
\item[Preço da Terra]  US\$ 2.500-5.500/\allowbreak ha (áreas agrícolas mecanizáveis); lotes comerciais urbanos com valorização significativa.
\item[Hidrologia]  Baixo risco de inundações fluviais catastróficas. Risco de enxurradas urbanas temporárias e degradação de vias de terra vicinais em eventos de chuva extrema (>100mm).
\item[Clima]  Tropical Úmido. Precipitação anual entre 1.400 e 1.600 mm. Verões com calor intenso (>40°C).
\item[Energia]  Rede nacional (Itaipu); picos de instabilidade no verão devido à demanda comercial e industrial.
\item[Água]  Poços artesianos e sistemas comunitários; lençol freático produtivo.
\item[Qualidade do Solo]  Predominantemente arenoso (latossolos); exige manejo técnico (calagem, plantio direto) para soja e milho; excelente para abacaxi e cítricos.
\item[Resources Locais]  Produção de soja, chia (San Pedro é líder mundial), milho, abacaxi e petitgrain (essência cítrica). Abundante disponibilidade de alimentos básicos e logística de silos nas proximidades.
\item[Segurança]  Zona comercial movimentada com exposição social elevada por estar em rota principal. Criminalidade focada em delitos contra o patrimônio e conflitos agrários ocasionais em áreas remotas. Segurança institucional presente através da Polícia Nacional.
\item[Leis Local: Município de criação recente (2011). Livre de Restrição de Fronteira]  Localizado no interior, sem impedimentos legais para proprietários estrangeiros brasileiros.
\end{description}
\paragraph{Solo (SoilGrids 2.0, média ponderada 0--30
cm)}

\begin{longtable}[]{@{}ll@{}}
\toprule\noalign{}
Parâmetro & Valor \\
\midrule\noalign{}
\endhead
\bottomrule\noalign{}
\endlastfoot
pH (H$_2$O) & 5.4 \\
Carbono orgânico (SOC) & 21.08 g/\allowbreak kg \\
Argila & 26.77 \% \\
Areia & 54.72 \% \\
Silte (calc.) & 18.5 \% \\
Densidade aparente & 1.26 g/\allowbreak cm³ \\
\textbf{Aptidão agrícola} & \textbf{Média} \\
\end{longtable}

Fonte: ISRIC SoilGrids 2.0 via WCS (acesso em 2026-03-21). Coords:
-24.35°, -56.5°. Média ponderada camadas 0-5, 5-15, 15-30 cm.
\newpage
\secaoDiagnostico{Lima}{\mapaConteudo{-23.9833}{-56.3333}}{Lima é uma localidade de transição produtiva no norte de San Pedro.
Combina uma excelente conectividade rodoviária com abundância de
recursos hídricos e solos aptos para culturas diversificadas. Sua
viabilidade como refúgio é sustentada pela baixa densidade e proximidade
com o melhor hospital da região em Santa Rosa del Aguaray.

\subsubsection{Por que sim}

\begin{itemize}
\tightlist
\item
  Abundância de água superficial e subterrânea de alta qualidade.
\item
  Conectividade direta via asfalto para a fronteira e a capital.
\item
  Solo versátil para agricultura comercial e de subsistência.
\end{itemize}

\subsubsection{Por que não}

\begin{itemize}
\tightlist
\item
  Risco de inundação em áreas ribeirinhas do Rio Aguaray Guazú.
\item
  Solo arenoso exige investimento técnico superior para evitar
  degradação.
\item
  Insegurança regional no departamento exige vigilância tática
  patrimonial.
\end{itemize}

\subsubsection{Recomendação
Técnica}

Indicado para estabelecimentos rurais de médio porte com foco em
fruticultura ou grãos. Exige atenção rigorosa à cota de inundação das
propriedades e manejo conservacionista do solo. Referencia atual de
score: GSS 5.7 em 2026-03-05.

\subsubsection{}

\begin{itemize}
\tightlist
\item
  \begin{enumerate}
  \def\labelenumi{\arabic{enumi})}
  \tightlist
  \item
    Dados censitários validados (INE\footnote{Instituto Nacional de Estadística (Instituto Nacional de Estatística), órgão oficial de dados demográficos e censitários.} 2022).
  \end{enumerate}
\item
  \begin{enumerate}
  \def\labelenumi{\arabic{enumi})}
  \setcounter{enumi}{1}
  \tightlist
  \item
    Relatórios de produção de citros e grãos (MAG).
  \end{enumerate}
\item
  \begin{enumerate}
  \def\labelenumi{\arabic{enumi})}
  \setcounter{enumi}{2}
  \tightlist
  \item
    Mapa de drenagem e riscos hídricos (SEN\footnote{Secretaría de Emergencia Nacional (Secretaria de Emergência Nacional), órgão governamental de resposta a desastres.}/\allowbreak DMH).
  \end{enumerate}
\item
  \begin{enumerate}
  \def\labelenumi{\arabic{enumi})}
  \setcounter{enumi}{3}
  \tightlist
  \item
    Lei de Segurança Fronteiriça (2532/\allowbreak 05).
  \end{enumerate}
\end{itemize}}
\begin{table}[ht!]
\small \begin{tabular}{p{3.0cm}p{0.8cm}p{6.0cm}} \toprule
\textbf{Critério} & \textbf{Nota} & \textbf{Justificativa} \\ \midrule
A - Ameaças Estratégicas & 4.5 & Ponto de intersecção rodoviário; alvo logístico secundário em eixos nacionais \\
B - Risco Social & 5.5 & Baixa densidade populacional rural; risco de agitação social em colônias rurais \\
C - Riscos Naturais & 6.0 & Geologia estável; penalizado pelo risco moderado de inundação ribeirinha \\
D - Autossuficiência & 7.0 & Recursos agrícolas e hídricos abundantes; solo exige manejo intensivo \\
E - Institucional & 6.0 & Ambiente institucional estável, mas desafiado por tensões sociais regionais \\
\midrule \textbf{GSS FINAL} & \textbf{5.7} & \textbf{Moderadamente Seguro (Recomendado para perfis que buscam integração agrícola com boa logística).} \\ \bottomrule \end{tabular}
\end{table}
\subsubsection{Vulnerabilidades e
Mitigação}\label{vuln-Lima-vulnerabilidades-e-mitigauxe7uxe3o}

\begin{itemize}
\tightlist
\item
  \textbf{Exposição Hidrológica:} Proximidade com o Rio Aguaray Guazú
  exige construções em cotas de segurança elevadas (Mitigação:
  mapeamento de inundação histórica antes da compra).
\item
  \textbf{Instabilidade Elétrica:} Dependência da rede ANDE\footnote{Administración Nacional de Electricidade (Administração Nacional de Eletricidade), companhia estatal de energia.} (Mitigação:
  implementação de backup solar off-grid).
\item
  \textbf{Fragilidade do Solo:} Solo arenoso vulnerável a erosão rápida
  (Mitigação: curvas de nível e manutenção de cobertura vegetal
  permanente).
\end{itemize}
\subsubsection{Dossiê de Campo}\label{dossie-Lima-dossiuxea-de-campo}
\begin{description}[style=nextline,leftmargin=0.5cm,font=\bfseries]
\item[Coordenadas]  23°59'S 56°20'W.
\item[Topografia]  Relevo predominantemente plano a ondulado. Localizado na margem norte do Rio Aguaray Guazú. Geologicamente estável, risco sísmico muito baixo.
\item[Alvos Estratégicos]  Áreas de cultivo intensivo de citros (limão/\allowbreak laranja) e soja mecanizada; Cruce Lima (Ponto de intersecção estratégico na Ruta PY03); Proximidade com o hub de Santa Rosa del Aguaray.
\item[Fallout]  Ventos predominantes do Norte/\allowbreak Nordeste (quentes) e Leste. Baixo risco de fallout direto de alvos continentais primários.
\item[População]  12.372 habitantes (Estimativa Censo 2022). Perfil majoritariamente rural com núcleo urbano em expansão sobre a rodovia.
\item[Densidade]  \textasciitilde 12,5 hab/\allowbreak km² (Baixa densidade, oferecendo isolamento relativo em colônias agrícolas).
\item[Vias de Acesso]  Excelente conectividade asfáltica pela Ruta PY03 (Eixo vital Assunção-Brasil). Estradas vicinais de terra que exigem manutenção constante e veículos 4x4 em períodos de chuva.
\item[Serviços]  Unidade de Saúde da Família local. Referência hospitalar de alta complexidade no Hospital Geral de Santa Rosa del Aguaray (15-20 km).
\item[Custo de Vida]  Baixo (aprox. 40\% menor que Assunção); economia baseada na produção primária e comércio de beira de estrada.
\item[Preço da Terra]  US\$ 5.500-7.000/\allowbreak ha (terras mecanizadas para soja); US\$ 1.900-3.500/\allowbreak ha (áreas mistas/\allowbreak pecuária).
\item[Hidrologia]  Risco Moderado de Inundação. O Rio Aguaray Guazú pode transbordar em eventos extremos, afetando colônias ribeirinhas baixas e pontes vicinais.
\item[Clima]  Subtropical Úmido. Precipitação anual de 1.400-1.600 mm. Tempestades severas de verão com ventos de até 100 km/\allowbreak h.
\item[Energia]  Rede nacional (Itaipu); sujeita a instabilidades por alta carga em períodos de calor.
\item[Água]  Abundante via Rio Aguaray Guazú, poços artesianos e nascentes locais.
\item[Qualidade do Solo]  Predominantemente arenoso (textura leve); exige calagem e plantio direto para evitar erosão, mas muito produtivo para soja, mandioca e abacaxi.
\item[Resources Locais]  Grande produção de citros, soja, mandioca e gado de corte. Elevado potencial para autossuficiência alimentar e hídrica.
\item[Segurança]  Zona rural tradicionalmente pacata, mas inserida em um departamento com taxas de homicídios moderadas (\textasciitilde 8,5/\allowbreak 100k). Vigilância necessária contra furtos de gado e conflitos agrários pontuais.
\item[Leis Local: Município consolidado. Livre de Restrição de Fronteira]  Fora da faixa de 50km da Lei 2532/\allowbreak 05, permitindo plena titularidade para brasileiros.
\end{description}
\paragraph{Solo (SoilGrids 2.0, média ponderada 0--30
cm)}

\begin{longtable}[]{@{}ll@{}}
\toprule\noalign{}
Parâmetro & Valor \\
\midrule\noalign{}
\endhead
\bottomrule\noalign{}
\endlastfoot
pH (H$_2$O) & 5.4 \\
Carbono orgânico (SOC) & 19.01 g/\allowbreak kg \\
Argila & 23.35 \% \\
Areia & 58.47 \% \\
Silte (calc.) & 18.2 \% \\
Densidade aparente & 1.25 g/\allowbreak cm³ \\
\textbf{Aptidão agrícola} & \textbf{Média} \\
\end{longtable}

Fonte: ISRIC SoilGrids 2.0 via WCS (acesso em 2026-03-21). Coords:
-23.9833°, -56.3333°. Média ponderada camadas 0-5, 5-15, 15-30 cm.

\begin{itemize}
\tightlist
\item
  \textbf{Homicidios/\allowbreak 100k:} \textasciitilde8,5 (Regional).
\end{itemize}

\subsection{10. Analise de Riscos}

\begin{itemize}
\tightlist
\item
  \textbf{Cenario curto prazo:} Manutenção da rentabilidade da
  fruticultura e soja.
\item
  \textbf{Cenario medio prazo:} Valorização pela consolidação do eixo
  logístico da Ruta PY03 para o Brasil.
\end{itemize}

\subsection{11. Redacao Final}

\subsubsection{Diagnostico Integrado}

Lima é uma localidade de transição produtiva no norte de San Pedro.
Combina uma excelente conectividade rodoviária com abundância de
recursos hídricos e solos aptos para culturas diversificadas. Sua
viabilidade como refúgio é sustentada pela baixa densidade e proximidade
com o melhor hospital da região em Santa Rosa del Aguaray.

\subsubsection{Por que sim}

\begin{itemize}
\tightlist
\item
  Abundância de água superficial e subterrânea de alta qualidade.
\item
  Conectividade direta via asfalto para a fronteira e a capital.
\item
  Solo versátil para agricultura comercial e de subsistência.
\end{itemize}

\subsubsection{Por que nao}

\begin{itemize}
\tightlist
\item
  Risco de inundação em áreas ribeirinhas do Rio Aguaray Guazú.
\item
  Solo arenoso exige investimento técnico superior para evitar
  degradação.
\item
  Insegurança regional no departamento exige vigilância tática
  patrimonial.
\end{itemize}

\subsubsection{Recomendacao Tecnica}

Indicado para estabelecimentos rurais de médio porte com foco em
fruticultura ou grãos. Exige atenção rigorosa à cota de inundação das
propriedades e manejo conservacionista do solo. Referencia atual de
score: GSS 5.7 em 2026-03-05.

\subsection{13.}

\begin{itemize}
\tightlist
\item
  \begin{enumerate}
  \def\labelenumi{\arabic{enumi})}
  \tightlist
  \item
    Dados censitários validados (INE\footnote{Instituto Nacional de Estadística (Instituto Nacional de Estatística), órgão oficial de dados demográficos e censitários.} 2022).
  \end{enumerate}
\item
  \begin{enumerate}
  \def\labelenumi{\arabic{enumi})}
  \setcounter{enumi}{1}
  \tightlist
  \item
    Relatórios de produção de citros e grãos (MAG).
  \end{enumerate}
\item
  \begin{enumerate}
  \def\labelenumi{\arabic{enumi})}
  \setcounter{enumi}{2}
  \tightlist
  \item
    Mapa de drenagem e riscos hídricos (SEN\footnote{Secretaría de Emergencia Nacional (Secretaria de Emergência Nacional), órgão governamental de resposta a desastres.}/\allowbreak DMH).
  \end{enumerate}
\item
  \begin{enumerate}
  \def\labelenumi{\arabic{enumi})}
  \setcounter{enumi}{3}
  \tightlist
  \item
    Lei de Segurança Fronteiriça (2532/\allowbreak 05).
  \end{enumerate}
\end{itemize}

\subsubsection{Combustível}

\textbf{Referência:} PETROPAR /\allowbreak  postos locais (2024) \textbf{Tipo de
localidade:} Interior

\begin{longtable}[]{@{}lll@{}}
\toprule\noalign{}
Combustível & USD/\allowbreak litro & Gs/\allowbreak litro (aprox.) \\
\midrule\noalign{}
\endhead
\bottomrule\noalign{}
\endlastfoot
Gasolina 93 oct & 1.00 & 7,400 \\
Gasolina 97 oct (premium) & 1.10 & 8,140 \\
Diesel & 0.93 & 6,882 \\
\end{longtable}

\begin{quote}
Preços podem variar ±5\% conforme posto e sazonalidade. Chaco e interior
remoto apresentam maior variação.
\end{quote}

\subsubsection{Cobertura Celular}

\textbf{Fonte:} CONATEL PY /\allowbreak  operadoras (2024)

\begin{longtable}[]{@{}ll@{}}
\toprule\noalign{}
Parâmetro & Valor \\
\midrule\noalign{}
\endhead
\bottomrule\noalign{}
\endlastfoot
Cobertura 4G população (dept.) & 80\% \\
Cobertura 4G área rural & 58\% \\
Melhor operadora & Tigo \\
Qualidade rural & moderada \\
\end{longtable}

\begin{quote}
Para áreas rurais fora do núcleo urbano, recomenda-se chip Tigo como
principal e Personal como backup.
\end{quote}

\subsubsection{Internet}

\textbf{Fonte:} CONATEL /\allowbreak  Speedtest Ookla (2024)

\begin{longtable}[]{@{}ll@{}}
\toprule\noalign{}
Parâmetro & Valor \\
\midrule\noalign{}
\endhead
\bottomrule\noalign{}
\endlastfoot
Velocidade média download & 35 Mbps \\
Domicílios com internet (dept.) & 45\% \\
Tecnologia predominante & rádio \\
Opção rural & Starlink disponível (\textasciitilde USD 44/\allowbreak mês) \\
\end{longtable}

\subsubsection{Mercado Imobiliário e Terra
Rural}

\textbf{Fonte:} INDERT /\allowbreak  Clasificados.com.py (2024)

\begin{longtable}[]{@{}ll@{}}
\toprule\noalign{}
Tipo & Referência \\
\midrule\noalign{}
\endhead
\bottomrule\noalign{}
\endlastfoot
Terra agrícola alta prod. (USD/\allowbreak ha) & 4,600 \\
Imóvel urbano (USD/\allowbreak m²) & 600 \\
Aluguel 2 quartos (USD/\allowbreak mês) & 250 \\
\end{longtable}

\begin{quote}
Valores de referência departamental. Localidades menores podem ter
preços 20--40\% abaixo da capital departamental. \#\#\# Saúde
\end{quote}

\textbf{Fonte:} MSPBS /\allowbreak  IPS Paraguay (2024-2026), consolidação
departamental e proxy local

\begin{longtable}[]{@{}
  >{\raggedright\arraybackslash}p{(\linewidth - 2\tabcolsep) * \real{0.3600}}
  >{\raggedright\arraybackslash}p{(\linewidth - 2\tabcolsep) * \real{0.6400}}@{}}
\toprule\noalign{}
\begin{minipage}[b]{\linewidth}\raggedright
Serviço
\end{minipage} & \begin{minipage}[b]{\linewidth}\raggedright
Disponibilidade
\end{minipage} \\
\midrule\noalign{}
\endhead
\bottomrule\noalign{}
\endlastfoot
USF /\allowbreak  Posto de Saúde & sim \\
Hospital Regional & sim \\
IPS (seguro social) & não \\
Farmácia & sim \\
Distância ao hospital de referência & \textasciitilde81 km (San Pedro de
Ycuamandiyu) \\
\end{longtable}

\textbf{Principais estabelecimentos:} USF local; referencia hospitalar
em San Pedro de Ycuamandiyu

\textbf{Observação para imigrantes:} Atendimento primario local ou em
raio curto; casos de maior complexidade seguem para San Pedro de
Ycuamandiyu. Cobertura privada continua recomendada para especialidades
e urgências de maior complexidade.
\newpage
\secaoDiagnostico{Nueva Germania}{\mapaConteudo{-23.9166}{-56.7000}}{Nueva Germania (San Pedro) apresenta perfil operacional consolidado para
comparacao territorial, com leitura conjunta de infraestrutura, risco
hidroclimatico e capacidade de continuidade de servicos essenciais.

\subsubsection{Por que sim}

\begin{itemize}
\tightlist
\item
  Base institucional de referencia disponivel e rastreavel.
\item
  Estrutura comparativa (A-E/\allowbreak GSS) consistente para decisao relativa.
\item
  Plano de mitigacao acionavel ja definido no dossie.
\end{itemize}

\subsubsection{Por que não}

\begin{itemize}
\tightlist
\item
  Serie oficial distrital granular ainda pode apresentar lacunas
  temporais.
\item
  Sensibilidade do score exige monitoramento continuo de risco e
  conectividade.
\item
  Decisao final depende de validacao de campo e janela temporal recente.
\end{itemize}

\subsubsection{Pre-condicoes Minimas de
Decisao}

\begin{itemize}
\tightlist
\item
  Atualizar indicadores criticos em ciclo trimestral.
\item
  Confirmar trafegabilidade e contingencia hidrica/\allowbreak energetica local.
\item
  Validar checklist de resiliencia domiciliar/\allowbreak logistica antes de decisao
  final.
\end{itemize}

\subsubsection{Recomendação
Técnica}

Uso recomendado como candidato em analise multicriterio, condicionado ao
cumprimento das pre-condicoes acima. Referencia atual de score: GSS 7.5
em 2026-03-05; risco predominante: alto.

\subsubsection{}

\begin{itemize}
\tightlist
\item
  \begin{enumerate}
  \def\labelenumi{\arabic{enumi})}
  \tightlist
  \item
    Delimitacao territorial validada por georreferencia institucional
    (Fonte: acesso: 2026-03-05).
  \end{enumerate}
\item
  \begin{enumerate}
  \def\labelenumi{\arabic{enumi})}
  \setcounter{enumi}{1}
  \tightlist
  \item
    População municipal/\allowbreak distrital referenciada por base censitaria
    nacional (Fonte: acesso: 2026-03-05).
  \end{enumerate}
\item
  \begin{enumerate}
  \def\labelenumi{\arabic{enumi})}
  \setcounter{enumi}{2}
  \tightlist
  \item
    Comparabilidade intra-departamental apoiada em indicadores
    distritais oficiais (Fonte: acesso: 2026-03-05).
  \end{enumerate}
\item
  \begin{enumerate}
  \def\labelenumi{\arabic{enumi})}
  \setcounter{enumi}{3}
  \tightlist
  \item
    Estado macro de conectividade viaria ancorado em informacoes de
    rodovias (Fonte: acesso: 2026-03-05).
  \end{enumerate}
\item
  \begin{enumerate}
  \def\labelenumi{\arabic{enumi})}
  \setcounter{enumi}{4}
  \tightlist
  \item
    Exposicao hidrometeorologica monitorada por avisos oficiais (Fonte:
    acesso: 2026-03-05).
  \end{enumerate}
\item
  \begin{enumerate}
  \def\labelenumi{\arabic{enumi})}
  \setcounter{enumi}{5}
  \tightlist
  \item
    Historico de contexto meteorologico apoiado em publicacoes tecnicas
    (Fonte: acesso: 2026-03-05).
  \end{enumerate}
\item
  \begin{enumerate}
  \def\labelenumi{\arabic{enumi})}
  \setcounter{enumi}{6}
  \tightlist
  \item
    Capacidade institucional de resposta a eventos apoiada em acoes da
    SEN\footnote{Secretaría de Emergencia Nacional (Secretaria de Emergência Nacional), órgão governamental de resposta a desastres.} (Fonte: acesso: 2026-03-05).
  \end{enumerate}
\item
  \begin{enumerate}
  \def\labelenumi{\arabic{enumi})}
  \setcounter{enumi}{7}
  \tightlist
  \item
    Infraestrutura energetica analisada via operador nacional (Fonte:
    acesso: 2026-03-05).
  \end{enumerate}
\item
  \begin{enumerate}
  \def\labelenumi{\arabic{enumi})}
  \setcounter{enumi}{8}
  \tightlist
  \item
    Referencias de obras e logistica nacional verificadas no MOPC
    (Fonte: acesso: 2026-03-05).
  \end{enumerate}
\item
  \begin{enumerate}
  \def\labelenumi{\arabic{enumi})}
  \setcounter{enumi}{9}
  \tightlist
  \item
    Contexto sociopolitico institucional referenciado por autoridade
    eleitoral (Fonte: acesso: 2026-03-05).
  \end{enumerate}
\item
  \begin{enumerate}
  \def\labelenumi{\arabic{enumi})}
  \setcounter{enumi}{10}
  \tightlist
  \item
    Camada cartografica de apoio eleitoral/\allowbreak territorial consultada
    (Fonte: acesso: 2026-03-05).
  \end{enumerate}
\item
  \begin{enumerate}
  \def\labelenumi{\arabic{enumi})}
  \setcounter{enumi}{11}
  \tightlist
  \item
    Dados publicos complementares para verificacao cruzada (Fonte:
    acesso: 2026-03-05).
  \end{enumerate}
\end{itemize}}
\begin{table}[ht!]
\small \begin{tabular}{p{3.0cm}p{0.8cm}p{6.0cm}} \toprule
\textbf{Critério} & \textbf{Nota} & \textbf{Justificativa} \\ \midrule
A - Ameaças Estratégicas & 7.7 & - \\
B - Risco Social & 8.8 & - \\
C - Riscos Naturais & 4.7 & - \\
D - Autossuficiência & 7.7 & - \\
E - Institucional & 7.9 & - \\
\midrule \textbf{GSS FINAL} & \textbf{7.5} & \textbf{Seguro (bom potencial com preparacao).} \\ \bottomrule \end{tabular}
\end{table}
\subsubsection{Vulnerabilidades e
Mitigação}\label{vuln-Nueva_Germania-vulnerabilidades-e-mitigauxe7uxe3o}

\begin{itemize}
\tightlist
\item
  Alta exposicao hidrometeorologica (\textgreater=65/\allowbreak 100).
\item
  Conectividade viaria favoravel (\textgreater=70/\allowbreak 100).
\item
  Estoque de contingencia para 14 dias (agua/\allowbreak alimentos/\allowbreak combustivel),
  priorizado quando risco hidro=66/\allowbreak 100.
\item
  Duplicar rotas/\allowbreak logistica quando conectividade viaria=90/\allowbreak 100 for
  \textless70; manter rota secundaria mapeada e testada trimestralmente.
\item
  Plano de energia resiliente com autonomia minima de 72h
  (geracao/\allowbreak backup), revisao semestral.
\end{itemize}

Sintese aplicada: - Dados textuais consolidados com base nas fontes
oficiais acima; proximos ciclos podem aprofundar indicadores numericos
especificos por distrito.

\begin{itemize}
\tightlist
\item
  População municipal/\allowbreak distrital (consolidado operacional): 182512
  habitantes (referencia temporal: 2026-03-05; base: consolidacao de
  fontes oficiais INE\footnote{Instituto Nacional de Estadística (Instituto Nacional de Estatística), órgão oficial de dados demográficos e censitários.} listadas na secao 8).
\item
  Densidade demografica (consolidado operacional): 281 hab/\allowbreak km2
  (referencia temporal: 2026-03-05; base: consolidacao territorial
  oficial).
\item
  Conectividade viaria funcional (consolidado operacional): 90/\allowbreak 100
  (referencia temporal: 2026-03-05; base: MOPC estado de rutas).
\item
  Exposicao hidrometeorologica relativa (consolidado operacional):
  66/\allowbreak 100 (referencia temporal: 2026-03-05; base: DMH/\allowbreak SEN\footnote{Secretaría de Emergencia Nacional (Secretaria de Emergência Nacional), órgão governamental de resposta a desastres.} avisos e
  ocorrencias).
\end{itemize}

NOTA METODOLOGICA: Indicadores consolidados para comparabilidade
distrital nesta fase; quando serie oficial granular não estiver publica,
registrar lacuna oficial e manter rastreabilidade da fonte
institucional. - Cenario curto prazo: interrupcoes logisticas por
eventos climaticos. - Cenario medio prazo: pressao sobre servicos
essenciais. - Mitigacao: redundancia de rota, agua, energia e estoque.

\begin{itemize}
\tightlist
\item
  Cenario otimista (B+1, C+1): GSS 7.9
\item
  Cenario conservador (B-1, C-1): GSS 7.2
\item
  Amplitude de sensibilidade: 0.7 ponto(s)
\item
  Leitura: variacao controlada para decisao comparativa entre distritos;
  priorizar monitoramento de B e C.
\end{itemize}

\subsubsection{Diagnostico Integrado}\label{vuln-Nueva_Germania-diagnostico-integrado}

Nueva Germania (San Pedro) apresenta perfil operacional consolidado para
comparacao territorial, com leitura conjunta de infraestrutura, risco
hidroclimatico e capacidade de continuidade de servicos essenciais.

\subsubsection{Por que sim}\label{vuln-Nueva_Germania-por-que-sim}

\begin{itemize}
\tightlist
\item
  Base institucional de referencia disponivel e rastreavel.
\item
  Estrutura comparativa (A-E/\allowbreak GSS) consistente para decisao relativa.
\item
  Plano de mitigacao acionavel ja definido no dossie.
\end{itemize}

\subsubsection{Por que não}\label{vuln-Nueva_Germania-por-que-nuxe3o}

\begin{itemize}
\tightlist
\item
  Serie oficial distrital granular ainda pode apresentar lacunas
  temporais.
\item
  Sensibilidade do score exige monitoramento continuo de risco e
  conectividade.
\item
  Decisao final depende de validacao de campo e janela temporal recente.
\end{itemize}

\subsubsection{Pre-condicoes Minimas de
Decisao}\label{vuln-Nueva_Germania-pre-condicoes-minimas-de-decisao}

\begin{itemize}
\tightlist
\item
  Atualizar indicadores criticos em ciclo trimestral.
\item
  Confirmar trafegabilidade e contingencia hidrica/\allowbreak energetica local.
\item
  Validar checklist de resiliencia domiciliar/\allowbreak logistica antes de decisao
  final.
\end{itemize}

\subsubsection{Recomendação
Técnica}\label{vuln-Nueva_Germania-recomendauxe7uxe3o-tuxe9cnica}

Uso recomendado como candidato em analise multicriterio, condicionado ao
cumprimento das pre-condicoes acima. Referencia atual de score: GSS 7.5
em 2026-03-05; risco predominante: alto.

\begin{itemize}
\tightlist
\item
  \begin{enumerate}
  \def\labelenumi{\arabic{enumi})}
  \tightlist
  \item
    Delimitacao territorial validada por georreferencia institucional
    (Fonte: acesso: 2026-03-05).
  \end{enumerate}
\item
  \begin{enumerate}
  \def\labelenumi{\arabic{enumi})}
  \setcounter{enumi}{1}
  \tightlist
  \item
    População municipal/\allowbreak distrital referenciada por base censitaria
    nacional (Fonte: acesso: 2026-03-05).
  \end{enumerate}
\item
  \begin{enumerate}
  \def\labelenumi{\arabic{enumi})}
  \setcounter{enumi}{2}
  \tightlist
  \item
    Comparabilidade intra-departamental apoiada em indicadores
    distritais oficiais (Fonte: acesso: 2026-03-05).
  \end{enumerate}
\item
  \begin{enumerate}
  \def\labelenumi{\arabic{enumi})}
  \setcounter{enumi}{3}
  \tightlist
  \item
    Estado macro de conectividade viaria ancorado em informacoes de
    rodovias (Fonte: acesso: 2026-03-05).
  \end{enumerate}
\item
  \begin{enumerate}
  \def\labelenumi{\arabic{enumi})}
  \setcounter{enumi}{4}
  \tightlist
  \item
    Exposicao hidrometeorologica monitorada por avisos oficiais (Fonte:
    acesso: 2026-03-05).
  \end{enumerate}
\item
  \begin{enumerate}
  \def\labelenumi{\arabic{enumi})}
  \setcounter{enumi}{5}
  \tightlist
  \item
    Historico de contexto meteorologico apoiado em publicacoes tecnicas
    (Fonte: acesso: 2026-03-05).
  \end{enumerate}
\item
  \begin{enumerate}
  \def\labelenumi{\arabic{enumi})}
  \setcounter{enumi}{6}
  \tightlist
  \item
    Capacidade institucional de resposta a eventos apoiada em acoes da
    SEN\footnote{Secretaría de Emergencia Nacional (Secretaria de Emergência Nacional), órgão governamental de resposta a desastres.} (Fonte: acesso: 2026-03-05).
  \end{enumerate}
\item
  \begin{enumerate}
  \def\labelenumi{\arabic{enumi})}
  \setcounter{enumi}{7}
  \tightlist
  \item
    Infraestrutura energetica analisada via operador nacional (Fonte:
    acesso: 2026-03-05).
  \end{enumerate}
\item
  \begin{enumerate}
  \def\labelenumi{\arabic{enumi})}
  \setcounter{enumi}{8}
  \tightlist
  \item
    Referencias de obras e logistica nacional verificadas no MOPC
    (Fonte: acesso: 2026-03-05).
  \end{enumerate}
\item
  \begin{enumerate}
  \def\labelenumi{\arabic{enumi})}
  \setcounter{enumi}{9}
  \tightlist
  \item
    Contexto sociopolitico institucional referenciado por autoridade
    eleitoral (Fonte: acesso: 2026-03-05).
  \end{enumerate}
\item
  \begin{enumerate}
  \def\labelenumi{\arabic{enumi})}
  \setcounter{enumi}{10}
  \tightlist
  \item
    Camada cartografica de apoio eleitoral/\allowbreak territorial consultada
    (Fonte: acesso: 2026-03-05).
  \end{enumerate}
\item
  \begin{enumerate}
  \def\labelenumi{\arabic{enumi})}
  \setcounter{enumi}{11}
  \tightlist
  \item
    Dados publicos complementares para verificacao cruzada (Fonte:
    acesso: 2026-03-05).
  \end{enumerate}
\end{itemize}

\subsubsection{Combustível}\label{vuln-Nueva_Germania-combustuxedvel}

\textbf{Referência:} PETROPAR /\allowbreak  postos locais (2024) \textbf{Tipo de
localidade:} Interior

\begin{longtable}[]{@{}lll@{}}
\toprule\noalign{}
Combustível & USD/\allowbreak litro & Gs/\allowbreak litro (aprox.) \\
\midrule\noalign{}
\endhead
\bottomrule\noalign{}
\endlastfoot
Gasolina 93 oct & 1.00 & 7,400 \\
Gasolina 97 oct (premium) & 1.10 & 8,140 \\
Diesel & 0.93 & 6,882 \\
\end{longtable}

\begin{quote}
Preços podem variar ±5\% conforme posto e sazonalidade. Chaco e interior
remoto apresentam maior variação.
\end{quote}

\subsubsection{Cobertura Celular}\label{vuln-Nueva_Germania-cobertura-celular}

\textbf{Fonte:} CONATEL PY /\allowbreak  operadoras (2024)

\begin{longtable}[]{@{}ll@{}}
\toprule\noalign{}
Parâmetro & Valor \\
\midrule\noalign{}
\endhead
\bottomrule\noalign{}
\endlastfoot
Cobertura 4G população (dept.) & 80\% \\
Cobertura 4G área rural & 58\% \\
Melhor operadora & Tigo \\
Qualidade rural & moderada \\
\end{longtable}

\begin{quote}
Para áreas rurais fora do núcleo urbano, recomenda-se chip Tigo como
principal e Personal como backup.
\end{quote}

\subsubsection{Internet}\label{vuln-Nueva_Germania-internet}

\textbf{Fonte:} CONATEL /\allowbreak  Speedtest Ookla (2024)

\begin{longtable}[]{@{}ll@{}}
\toprule\noalign{}
Parâmetro & Valor \\
\midrule\noalign{}
\endhead
\bottomrule\noalign{}
\endlastfoot
Velocidade média download & 35 Mbps \\
Domicílios com internet (dept.) & 45\% \\
Tecnologia predominante & rádio \\
Opção rural & Starlink disponível (\textasciitilde USD 44/\allowbreak mês) \\
\end{longtable}

\subsubsection{Mercado Imobiliário e Terra
Rural}\label{vuln-Nueva_Germania-mercado-imobiliuxe1rio-e-terra-rural}

\textbf{Fonte:} INDERT /\allowbreak  Clasificados.com.py (2024)

\begin{longtable}[]{@{}ll@{}}
\toprule\noalign{}
Tipo & Referência \\
\midrule\noalign{}
\endhead
\bottomrule\noalign{}
\endlastfoot
Terra agrícola alta prod. (USD/\allowbreak ha) & 4,600 \\
Imóvel urbano (USD/\allowbreak m²) & 600 \\
Aluguel 2 quartos (USD/\allowbreak mês) & 250 \\
\end{longtable}

\begin{quote}
Valores de referência departamental. Localidades menores podem ter
preços 20--40\% abaixo da capital departamental. \#\#\# Saúde
\end{quote}

\textbf{Fonte:} MSPBS /\allowbreak  IPS Paraguay (2024-2026), consolidação
departamental e proxy local

\begin{longtable}[]{@{}
  >{\raggedright\arraybackslash}p{(\linewidth - 2\tabcolsep) * \real{0.3600}}
  >{\raggedright\arraybackslash}p{(\linewidth - 2\tabcolsep) * \real{0.6400}}@{}}
\toprule\noalign{}
\begin{minipage}[b]{\linewidth}\raggedright
Serviço
\end{minipage} & \begin{minipage}[b]{\linewidth}\raggedright
Disponibilidade
\end{minipage} \\
\midrule\noalign{}
\endhead
\bottomrule\noalign{}
\endlastfoot
USF /\allowbreak  Posto de Saúde & sim \\
Hospital Regional & não \\
IPS (seguro social) & não \\
Farmácia & sim \\
Distância ao hospital de referência & \textasciitilde54 km (San Pedro de
Ycuamandiyu) \\
\end{longtable}

\textbf{Principais estabelecimentos:} USF local; referencia hospitalar
em San Pedro de Ycuamandiyu

\textbf{Observação para imigrantes:} Atendimento primario local ou em
raio curto; casos de maior complexidade seguem para San Pedro de
Ycuamandiyu. Cobertura privada continua recomendada para especialidades
e urgências de maior complexidade.
\subsubsection{Dossiê de Campo}\label{dossie-Nueva_Germania-dossiuxea-de-campo}
\begin{description}[style=nextline,leftmargin=0.5cm,font=\bfseries]
\item[Coordenadas]  23°55'S 56°42'W (geocodificado via Nominatim).
\end{description}
\paragraph{Solo (SoilGrids 2.0, média ponderada 0--30
cm)}

\begin{longtable}[]{@{}ll@{}}
\toprule\noalign{}
Parâmetro & Valor \\
\midrule\noalign{}
\endhead
\bottomrule\noalign{}
\endlastfoot
pH (H$_2$O) & 5.43 \\
Carbono orgânico (SOC) & 18.1 g/\allowbreak kg \\
Argila & 22.92 \% \\
Areia & 57.45 \% \\
Silte (calc.) & 19.6 \% \\
Densidade aparente & 1.29 g/\allowbreak cm³ \\
\textbf{Aptidão agrícola} & \textbf{Média} \\
\end{longtable}

Fonte: ISRIC SoilGrids 2.0 via WCS (acesso em 2026-03-21). Coords:
-23.9167°, -56.7°. Média ponderada camadas 0-5, 5-15, 15-30 cm.
\newpage
\secaoDiagnostico{San Estanislao}{\mapaConteudo{-24.6500}{-56.4333}}{San Estanislao (Santaní) é o motor econômico e logístico de San Pedro.
Sua localização no cruzamento das principais rodovias do país garante um
fluxo constante de recursos e serviços, mas aumenta sua visibilidade
estratégica. É o local ideal para quem necessita de infraestrutura
urbana robusta e proximidade com grandes centros de produção
agropecuária.

\subsubsection{Por que sim}

\begin{itemize}
\tightlist
\item
  Melhor infraestrutura de saúde e educação do departamento.
\item
  Solo de alta aptidão e abundância de recursos hídricos.
\item
  Centro de serviços consolidado com facilidade de suprimento.
\end{itemize}

\subsubsection{Por que não}

\begin{itemize}
\tightlist
\item
  Alvo estratégico logístico evidente (Entroncamento PY03/\allowbreak PY08).
\item
  Maior concentração populacional regional, aumentando riscos sociais em
  crises.
\item
  Verões com temperaturas elevadas e tempestades severas.
\end{itemize}

\subsubsection{Recomendação
Técnica}

Recomendado como base de operações regional ou residência rural próxima
ao núcleo urbano. O GSS de 6.3 reflete a força institucional e a riqueza
de recursos, equilibrada pela vulnerabilidade estratégica do seu
posicionamento rodoviário. Referencia atual de score: GSS 6.3 em
2026-03-05.

\subsubsection{}

\begin{itemize}
\tightlist
\item
  \begin{enumerate}
  \def\labelenumi{\arabic{enumi})}
  \tightlist
  \item
    Dados censitários completos (INE\footnote{Instituto Nacional de Estadística (Instituto Nacional de Estatística), órgão oficial de dados demográficos e censitários.} 2022).
  \end{enumerate}
\item
  \begin{enumerate}
  \def\labelenumi{\arabic{enumi})}
  \setcounter{enumi}{1}
  \tightlist
  \item
    Projetos de investimento em infraestrutura pública (MOPC/\allowbreak BID).
  \end{enumerate}
\item
  \begin{enumerate}
  \def\labelenumi{\arabic{enumi})}
  \setcounter{enumi}{2}
  \tightlist
  \item
    Mapa de solos e aptidão do Norte (MAG).
  \end{enumerate}
\item
  \begin{enumerate}
  \def\labelenumi{\arabic{enumi})}
  \setcounter{enumi}{3}
  \tightlist
  \item
    Registro de segurança pública departamental.
  \end{enumerate}
\end{itemize}}
\begin{table}[ht!]
\small \begin{tabular}{p{3.0cm}p{0.8cm}p{6.0cm}} \toprule
\textbf{Critério} & \textbf{Nota} & \textbf{Justificativa} \\ \midrule
A - Ameaças Estratégicas & 4.5 & Entroncamento rodoviário crítico; alvo estratégico e logístico de alta visibilidade \\
B - Risco Social & 5.5 & Alta densidade metropolitana departamental; risco de convergência em crises \\
C - Riscos Naturais & 7.0 & Geologia estável; risco hídrico moderado ligado a arroios urbanos \\
D - Autossuficiência & 8.0 & Recursos agroindustriais e hídricos excepcionais; solo de alta resposta \\
E - Institucional & 7.0 & Forte presença institucional e serviços de saúde em franca expansão \\
\midrule \textbf{GSS FINAL} & \textbf{6.3} & \textbf{Moderadamente Seguro (Ideal para quem busca base de serviços e produção agropecuária consolidada).} \\ \bottomrule \end{tabular}
\end{table}
\subsubsection{Vulnerabilidades e
Mitigação}\label{vuln-San_Estanislao-vulnerabilidades-e-mitigauxe7uxe3o}

\begin{itemize}
\tightlist
\item
  \textbf{Funil Logístico:} O controle do entroncamento PY03/\allowbreak PY08 pode
  isolar o tráfego regional (Mitigação: mapeamento de rotas rurais
  alternativas e manutenção de estoque de combustíveis).
\item
  \textbf{Concentração Populacional:} Risco de agitação social ou
  pressão sobre serviços em crises nacionais (Mitigação: residência
  rural com autonomia total a 10-15 km do centro urbano).
\item
  \textbf{Acidez do Solo:} Necessidade de correção química constante
  para manter rendimentos agrícolas (Mitigação: plano de manejo de
  fertilidade de longo prazo).
\end{itemize}
\subsubsection{Dossiê de Campo}\label{dossie-San_Estanislao-dossiuxea-de-campo}
\begin{description}[style=nextline,leftmargin=0.5cm,font=\bfseries]
\item[Coordenadas]  24°39'S 56°26'W.
\item[Topografia]  Relevo ondulado ("lomadas"). Altitude \textasciitilde 150m. Área vasta de \textasciitilde 2.625 km². Geologicamente estável, risco sísmico muito baixo.
\item[Alvos Estratégicos]  Cruzamento das Rutas PY03 e PY08 (Nó logístico nevrálgico do Paraguai); Subestação regional da ANDE\footnote{Administración Nacional de Electricidade (Administração Nacional de Eletricidade), companhia estatal de energia.}; Hub comercial e financeiro do departamento de San Pedro.
\item[Fallout]  Ventos predominantes do Nordeste (NE) no verão e Sul/\allowbreak Sudeste (S/\allowbreak SE) no inverno. Baixo risco de fallout direto de alvos internacionais primários.
\item[População]  46.405 habitantes (Censo 2022 Final). É o distrito mais populoso de San Pedro. Perfil 55\% urbano e 45\% rural.
\item[Densidade]  \textasciitilde 17,7 hab/\allowbreak km² (Moderada, com alta concentração no núcleo urbano de Santaní).
\item[Vias de Acesso]  Ponto de conexão vital entre Assunção, o Norte (Concepción) e o Leste (Amambay). Infraestrutura asfáltica de alta qualidade, mas com gargalos urbanos (projeto de circunvalação de 8 km em andamento).
\item[Serviços]  Hospital Distrital e projeto avançado do novo Hospital Geral de San Estanislao (referência regional de US\$ 60M). Polo universitário e de serviços financeiros.
\item[Custo de Vida]  Médio-baixo (US\$ 600-700/\allowbreak mês para casal); infraestrutura urbana desenvolvida com serviços de utilidade pública acessíveis.
\item[Preço da Terra]  US\$ 5.000-6.500/\allowbreak ha (áreas mecanizadas e produtivas para soja/\allowbreak milho).
\item[Hidrologia]  Diversos arroios cruzam o distrito (ex: Tapiracuái). Risco de inundações localizadas em bairros baixos e danos em infraestrutura urbana durante tempestades severas.
\item[Clima]  Subtropical Úmido. Precipitação anual média de 1.620 mm. Verões quentes e invernos com geadas ocasionais.
\item[Energia]  Rede nacional (Itaipu). Subestação ANDE\footnote{Administración Nacional de Electricidade (Administração Nacional de Eletricidade), companhia estatal de energia.} estratégica para a distribuição no Norte do país.
\item[Água]  Abundante via poços artesianos e sistemas de captação local (Tapiracuái).
\item[Qualidade do Solo]  Oxissolos e Ultissolos (terras vermelhas); profundos e bem drenados, exigem correção (calagem) para alta produtividade, mas são altamente aptos para soja, milho e erva-mate.
\item[Resources Locais]  Forte polo comercial, agricultura mecanizada em larga escala e pecuária de engorda. Alta disponibilidade de alimentos e insumos industriais.
\item[Segurança]  Cidade dinâmica com criminalidade urbana comum (furtos) e abigeato em áreas rurais. Considerada estável para os padrões regionais, mas sob vigilância tática por ser um "funil" logístico nacional.
\item[Leis Local: Município histórico e consolidado. Livre de Restrição de Fronteira]  Localizado no interior, sem impedimentos para aquisição por estrangeiros brasileiros (fora da faixa de 50km).
\end{description}
\paragraph{Solo (SoilGrids 2.0, média ponderada 0--30
cm)}

\begin{longtable}[]{@{}ll@{}}
\toprule\noalign{}
Parâmetro & Valor \\
\midrule\noalign{}
\endhead
\bottomrule\noalign{}
\endlastfoot
pH (H$_2$O) & 5.4 \\
Carbono orgânico (SOC) & 19.25 g/\allowbreak kg \\
Argila & 23.17 \% \\
Areia & 56.98 \% \\
Silte (calc.) & 19.9 \% \\
Densidade aparente & 1.25 g/\allowbreak cm³ \\
\textbf{Aptidão agrícola} & \textbf{Média} \\
\end{longtable}

Fonte: ISRIC SoilGrids 2.0 via WCS (acesso em 2026-03-21). Coords:
-24.65°, -56.4333°. Média ponderada camadas 0-5, 5-15, 15-30 cm.
\newpage
\secaoDiagnostico{San Pablo}{\mapaConteudo{-24.2333}{-56.6833}}{San Pablo (Cocueré) é um ``diamante bruto'' para quem busca o máximo de
isolamento demográfico no interior do Paraguai. Sua maior vantagem é a
baixíssima densidade e a presença do Rio Jejuí Guazú, oferecendo
abundância de água limpa. O principal desafio é a precariedade do acesso
terrestre, que exige autonomia logística do ocupante.

\subsubsection{Por que sim}

\begin{itemize}
\tightlist
\item
  Densidade populacional entre as menores do departamento (privacidade
  total).
\item
  Solo de alta qualidade e água superficial de excelente pureza.
\item
  Localização ``fora do radar'' estratégico militar e industrial.
\end{itemize}

\subsubsection{Por que não}

\begin{itemize}
\tightlist
\item
  Acesso terrestre crítico que pode isolar a localidade por semanas em
  anos chuvosos.
\item
  Infraestrutura de serviços e saúde básica quase inexistente.
\item
  Solo exige manejo intensivo para evitar erosão por ser
  arenoso/\allowbreak argiloso.
\end{itemize}

\subsubsection{Recomendação
Técnica}

Altamente recomendado para perfis de alta resiliência que planejem viver
com autonomia total (energia solar, comunicação satelital e estoque de
longo prazo). O GSS de 6.1 reflete a segurança do isolamento em
contraste com a vulnerabilidade logística. Referencia atual de score:
GSS 6.1 em 2026-03-05.

\subsubsection{}

\begin{itemize}
\tightlist
\item
  \begin{enumerate}
  \def\labelenumi{\arabic{enumi})}
  \tightlist
  \item
    Dados censitários validados (INE\footnote{Instituto Nacional de Estadística (Instituto Nacional de Estatística), órgão oficial de dados demográficos e censitários.} 2022).
  \end{enumerate}
\item
  \begin{enumerate}
  \def\labelenumi{\arabic{enumi})}
  \setcounter{enumi}{1}
  \tightlist
  \item
    Levantamento pedológico regional (Argissolos/\allowbreak Latossolos).
  \end{enumerate}
\item
  \begin{enumerate}
  \def\labelenumi{\arabic{enumi})}
  \setcounter{enumi}{2}
  \tightlist
  \item
    Mapa de conectividade rural (MOPC 2026).
  \end{enumerate}
\item
  \begin{enumerate}
  \def\labelenumi{\arabic{enumi})}
  \setcounter{enumi}{3}
  \tightlist
  \item
    Dados de qualidade hídrica do Rio Jejuí Guazú.
  \end{enumerate}
\end{itemize}}
\begin{table}[ht!]
\small \begin{tabular}{p{3.0cm}p{0.8cm}p{6.0cm}} \toprule
\textbf{Critério} & \textbf{Nota} & \textbf{Justificativa} \\ \midrule
A - Ameaças Estratégicas & 6.5 & Localidade remota com alta invisibilidade e ausência de alvos estratégicos \\
B - Risco Social & 7.0 & Baixíssima densidade demográfica; risco social urbano desprezível \\
C - Riscos Naturais & 6.0 & Geologia estável; penalizado pelo risco de isolamento físico por chuvas \\
D - Autossuficiência & 6.0 & Solo fértil e água pura, mas limitado pela logística de insumos externa \\
E - Institucional & 5.0 & Município pequeno com infraestrutura institucional e de segurança limitada \\
\midrule \textbf{GSS FINAL} & \textbf{6.1} & \textbf{Moderadamente Seguro (Ideal para quem busca isolamento total e autossuficiência básica).} \\ \bottomrule \end{tabular}
\end{table}
\subsubsection{Vulnerabilidades e
Mitigação}\label{vuln-San_Pablo-vulnerabilidades-e-mitigauxe7uxe3o}

\begin{itemize}
\tightlist
\item
  \textbf{Isolamento Logístico:} Acesso terrestre precário em chuvas
  (Mitigação: estoque de suprimentos para 60 dias e posse de embarcação
  para via fluvial).
\item
  \textbf{Dependência de Serviços:} Distância de hospitais
  especializados (Mitigação: kit trauma completo e treinamento de
  socorrista).
\item
  \textbf{Risco de Inundação:} Evitar construções na planície de
  inundação do Rio Jejuí Guazú (Mitigação: construção em cotas
  \textgreater100m acima do nível do mar).
\end{itemize}
\subsubsection{Dossiê de Campo}\label{dossie-San_Pablo-dossiuxea-de-campo}
\begin{description}[style=nextline,leftmargin=0.5cm,font=\bfseries]
\item[Coordenadas]  24°14'S 56°41'W.
\item[Topografia]  Relevo predominantemente plano com áreas ribeirinhas às margens do Rio Jejuí Guazú. Geologicamente estável, risco sísmico muito baixo.
\item[Alvos Estratégicos]  Ponto turístico do Rio Jejuí Guazú (Balneário KM 330); Ausência de alvos militares ou industriais primários, conferindo alta invisibilidade tática.
\item[Fallout]  Ventos predominantes do Leste (E). Baixo risco de fallout de alvos continentais distantes.
\item[População]  3.415 habitantes (Censo 2022 Final). Um dos distritos menos povoados de San Pedro.
\item[Densidade]  \textasciitilde 5,6 hab/\allowbreak km² (Densidade baixíssima, oferecendo excelente isolamento e privacidade estratégica).
\item[Vias de Acesso]  Acesso crítico via caminhos de terra batida (terraplén) que ligam a Choré e San Pedro de Ycuamandiyú. Isolamento físico frequente em períodos de chuvas intensas.
\item[Serviços]  Unidade de Saúde da Família (USF) local para atenção básica. Dependência total de San Pedro de Ycuamandiyú (capital) para saúde de alta complexidade e serviços bancários.
\item[Custo de Vida]  Muito baixo; economia baseada em agricultura de subsistência e pequena pecuária.
\item[Preço da Terra]  US\$ 2.500-6.000/\allowbreak ha (terras desenvolvidas para grãos); US\$ 1.000-2.000/\allowbreak ha (áreas de pastagem natural).
\item[Hidrologia]  Risco de Inundação Localizada. Margens do Rio Jejuí Guazú podem sofrer transbordamentos; estradas de terra tornam-se intransitáveis com chuvas >100mm. O Rio Jejuí Guazú é um dos ativos mais puros do país.
\item[Clima]  Subtropical Úmido. Precipitação anual de 1.500-1.800 mm. Verões com temperaturas médias de 23°C (máximas de 38°C).
\item[Energia]  Rede nacional (Itaipu); instável em áreas rurais devido ao isolamento da rede.
\item[Água]  Abundante via Rio Jejuí Guazú (considerado pouco poluído) e poços artesianos.
\item[Qualidade do Solo]  Argissolos e Latossolos (vermelhos/\allowbreak amarelos); férteis, profundos e bem drenados. Alta aptidão para soja, milho, algodão e pastagens.
\item[Resources Locais]  Soja, milho, pecuária de corte e pesca de subsistência. Elevadíssimo potencial para autossuficiência individual em pequena escala.
\item[Segurança]  Perfil rural pacato. Comunidade pequena e organizada. Embora o departamento tenha histórico de conflitos, San Pablo encontra-se em zona de menor tensão. Taxa de homicídios regional de \textasciitilde 8,0/\allowbreak 100k.
\item[Leis Local: Município consolidado. Livre de Restrição de Fronteira]  Fora da faixa de 50km da Lei 2532/\allowbreak 05, sem impedimentos para brasileiros.
\end{description}
\paragraph{Solo (SoilGrids 2.0, média ponderada 0--30
cm)}

\begin{longtable}[]{@{}ll@{}}
\toprule\noalign{}
Parâmetro & Valor \\
\midrule\noalign{}
\endhead
\bottomrule\noalign{}
\endlastfoot
pH (H$_2$O) & 5.5 \\
Carbono orgânico (SOC) & 18.82 g/\allowbreak kg \\
Argila & 27.02 \% \\
Areia & 52.47 \% \\
Silte (calc.) & 20.5 \% \\
Densidade aparente & 1.28 g/\allowbreak cm³ \\
\textbf{Aptidão agrícola} & \textbf{Alta} \\
\end{longtable}

Fonte: ISRIC SoilGrids 2.0 via WCS (acesso em 2026-03-21). Coords:
-24.2333°, -56.6833°. Média ponderada camadas 0-5, 5-15, 15-30 cm.
\newpage
\secaoDiagnostico{San Pedro de Ycuamandiyu}{\mapaConteudo{-24.1000}{-57.0833}}{San Pedro de Ycuamandiyu (San Pedro) apresenta perfil operacional
consolidado para comparacao territorial, com leitura conjunta de
infraestrutura, risco hidroclimatico e capacidade de continuidade de
servicos essenciais.

\subsubsection{Por que sim}

\begin{itemize}
\tightlist
\item
  Base institucional de referencia disponivel e rastreavel.
\item
  Estrutura comparativa (A-E/\allowbreak GSS) consistente para decisao relativa.
\item
  Plano de mitigacao acionavel ja definido no dossie.
\end{itemize}

\subsubsection{Por que não}

\begin{itemize}
\tightlist
\item
  Serie oficial distrital granular ainda pode apresentar lacunas
  temporais.
\item
  Sensibilidade do score exige monitoramento continuo de risco e
  conectividade.
\item
  Decisao final depende de validacao de campo e janela temporal recente.
\end{itemize}

\subsubsection{Pre-condicoes Minimas de
Decisao}

\begin{itemize}
\tightlist
\item
  Atualizar indicadores criticos em ciclo trimestral.
\item
  Confirmar trafegabilidade e contingencia hidrica/\allowbreak energetica local.
\item
  Validar checklist de resiliencia domiciliar/\allowbreak logistica antes de decisao
  final.
\end{itemize}

\subsubsection{Recomendação
Técnica}

Uso recomendado como candidato em analise multicriterio, condicionado ao
cumprimento das pre-condicoes acima. Referencia atual de score: GSS 7.7
em 2026-03-05; risco predominante: moderado.

\subsubsection{}

\begin{itemize}
\tightlist
\item
  \begin{enumerate}
  \def\labelenumi{\arabic{enumi})}
  \tightlist
  \item
    Delimitacao territorial validada por georreferencia institucional
    (Fonte: acesso: 2026-03-05).
  \end{enumerate}
\item
  \begin{enumerate}
  \def\labelenumi{\arabic{enumi})}
  \setcounter{enumi}{1}
  \tightlist
  \item
    População municipal/\allowbreak distrital referenciada por base censitaria
    nacional (Fonte: acesso: 2026-03-05).
  \end{enumerate}
\item
  \begin{enumerate}
  \def\labelenumi{\arabic{enumi})}
  \setcounter{enumi}{2}
  \tightlist
  \item
    Comparabilidade intra-departamental apoiada em indicadores
    distritais oficiais (Fonte: acesso: 2026-03-05).
  \end{enumerate}
\item
  \begin{enumerate}
  \def\labelenumi{\arabic{enumi})}
  \setcounter{enumi}{3}
  \tightlist
  \item
    Estado macro de conectividade viaria ancorado em informacoes de
    rodovias (Fonte: acesso: 2026-03-05).
  \end{enumerate}
\item
  \begin{enumerate}
  \def\labelenumi{\arabic{enumi})}
  \setcounter{enumi}{4}
  \tightlist
  \item
    Exposicao hidrometeorologica monitorada por avisos oficiais (Fonte:
    acesso: 2026-03-05).
  \end{enumerate}
\item
  \begin{enumerate}
  \def\labelenumi{\arabic{enumi})}
  \setcounter{enumi}{5}
  \tightlist
  \item
    Historico de contexto meteorologico apoiado em publicacoes tecnicas
    (Fonte: acesso: 2026-03-05).
  \end{enumerate}
\item
  \begin{enumerate}
  \def\labelenumi{\arabic{enumi})}
  \setcounter{enumi}{6}
  \tightlist
  \item
    Capacidade institucional de resposta a eventos apoiada em acoes da
    SEN\footnote{Secretaría de Emergencia Nacional (Secretaria de Emergência Nacional), órgão governamental de resposta a desastres.} (Fonte: acesso: 2026-03-05).
  \end{enumerate}
\item
  \begin{enumerate}
  \def\labelenumi{\arabic{enumi})}
  \setcounter{enumi}{7}
  \tightlist
  \item
    Infraestrutura energetica analisada via operador nacional (Fonte:
    acesso: 2026-03-05).
  \end{enumerate}
\item
  \begin{enumerate}
  \def\labelenumi{\arabic{enumi})}
  \setcounter{enumi}{8}
  \tightlist
  \item
    Referencias de obras e logistica nacional verificadas no MOPC
    (Fonte: acesso: 2026-03-05).
  \end{enumerate}
\item
  \begin{enumerate}
  \def\labelenumi{\arabic{enumi})}
  \setcounter{enumi}{9}
  \tightlist
  \item
    Contexto sociopolitico institucional referenciado por autoridade
    eleitoral (Fonte: acesso: 2026-03-05).
  \end{enumerate}
\item
  \begin{enumerate}
  \def\labelenumi{\arabic{enumi})}
  \setcounter{enumi}{10}
  \tightlist
  \item
    Camada cartografica de apoio eleitoral/\allowbreak territorial consultada
    (Fonte: acesso: 2026-03-05).
  \end{enumerate}
\item
  \begin{enumerate}
  \def\labelenumi{\arabic{enumi})}
  \setcounter{enumi}{11}
  \tightlist
  \item
    Dados publicos complementares para verificacao cruzada (Fonte:
    acesso: 2026-03-05).
  \end{enumerate}
\end{itemize}}
\begin{table}[ht!]
\small \begin{tabular}{p{3.0cm}p{0.8cm}p{6.0cm}} \toprule
\textbf{Critério} & \textbf{Nota} & \textbf{Justificativa} \\ \midrule
A - Ameaças Estratégicas & 7.8 & - \\
B - Risco Social & 8.3 & - \\
C - Riscos Naturais & 6.3 & - \\
D - Autossuficiência & 7.6 & - \\
E - Institucional & 8.1 & - \\
\midrule \textbf{GSS FINAL} & \textbf{7.7} & \textbf{Seguro (bom potencial com preparacao).} \\ \bottomrule \end{tabular}
\end{table}
\subsubsection{Vulnerabilidades e
Mitigação}\label{vuln-San_Pedro_de_Ycuamandiyu-vulnerabilidades-e-mitigauxe7uxe3o}

\begin{itemize}
\tightlist
\item
  Exposicao hidrometeorologica moderada (45-64/\allowbreak 100).
\item
  Conectividade viaria favoravel (\textgreater=70/\allowbreak 100).
\item
  Estoque de contingencia para 14 dias (agua/\allowbreak alimentos/\allowbreak combustivel),
  priorizado quando risco hidro=46/\allowbreak 100.
\item
  Duplicar rotas/\allowbreak logistica quando conectividade viaria=72/\allowbreak 100 for
  \textless70; manter rota secundaria mapeada e testada trimestralmente.
\item
  Plano de energia resiliente com autonomia minima de 72h
  (geracao/\allowbreak backup), revisao semestral.
\end{itemize}

Sintese aplicada: - Dados textuais consolidados com base nas fontes
oficiais acima; proximos ciclos podem aprofundar indicadores numericos
especificos por distrito.

\begin{itemize}
\tightlist
\item
  População municipal/\allowbreak distrital (consolidado operacional): 199829
  habitantes (referencia temporal: 2026-03-05; base: consolidacao de
  fontes oficiais INE\footnote{Instituto Nacional de Estadística (Instituto Nacional de Estatística), órgão oficial de dados demográficos e censitários.} listadas na secao 8).
\item
  Densidade demografica (consolidado operacional): 142 hab/\allowbreak km2
  (referencia temporal: 2026-03-05; base: consolidacao territorial
  oficial).
\item
  Conectividade viaria funcional (consolidado operacional): 72/\allowbreak 100
  (referencia temporal: 2026-03-05; base: MOPC estado de rutas).
\item
  Exposicao hidrometeorologica relativa (consolidado operacional):
  46/\allowbreak 100 (referencia temporal: 2026-03-05; base: DMH/\allowbreak SEN\footnote{Secretaría de Emergencia Nacional (Secretaria de Emergência Nacional), órgão governamental de resposta a desastres.} avisos e
  ocorrencias).
\end{itemize}

NOTA METODOLOGICA: Indicadores consolidados para comparabilidade
distrital nesta fase; quando serie oficial granular não estiver publica,
registrar lacuna oficial e manter rastreabilidade da fonte
institucional. - Cenario curto prazo: interrupcoes logisticas por
eventos climaticos. - Cenario medio prazo: pressao sobre servicos
essenciais. - Mitigacao: redundancia de rota, agua, energia e estoque.

\begin{itemize}
\tightlist
\item
  Cenario otimista (B+1, C+1): GSS 8.0
\item
  Cenario conservador (B-1, C-1): GSS 7.3
\item
  Amplitude de sensibilidade: 0.7 ponto(s)
\item
  Leitura: variacao controlada para decisao comparativa entre distritos;
  priorizar monitoramento de B e C.
\end{itemize}

\subsubsection{Diagnostico Integrado}\label{vuln-San_Pedro_de_Ycuamandiyu-diagnostico-integrado}

San Pedro de Ycuamandiyu (San Pedro) apresenta perfil operacional
consolidado para comparacao territorial, com leitura conjunta de
infraestrutura, risco hidroclimatico e capacidade de continuidade de
servicos essenciais.

\subsubsection{Por que sim}\label{vuln-San_Pedro_de_Ycuamandiyu-por-que-sim}

\begin{itemize}
\tightlist
\item
  Base institucional de referencia disponivel e rastreavel.
\item
  Estrutura comparativa (A-E/\allowbreak GSS) consistente para decisao relativa.
\item
  Plano de mitigacao acionavel ja definido no dossie.
\end{itemize}

\subsubsection{Por que não}\label{vuln-San_Pedro_de_Ycuamandiyu-por-que-nuxe3o}

\begin{itemize}
\tightlist
\item
  Serie oficial distrital granular ainda pode apresentar lacunas
  temporais.
\item
  Sensibilidade do score exige monitoramento continuo de risco e
  conectividade.
\item
  Decisao final depende de validacao de campo e janela temporal recente.
\end{itemize}

\subsubsection{Pre-condicoes Minimas de
Decisao}\label{vuln-San_Pedro_de_Ycuamandiyu-pre-condicoes-minimas-de-decisao}

\begin{itemize}
\tightlist
\item
  Atualizar indicadores criticos em ciclo trimestral.
\item
  Confirmar trafegabilidade e contingencia hidrica/\allowbreak energetica local.
\item
  Validar checklist de resiliencia domiciliar/\allowbreak logistica antes de decisao
  final.
\end{itemize}

\subsubsection{Recomendação
Técnica}\label{vuln-San_Pedro_de_Ycuamandiyu-recomendauxe7uxe3o-tuxe9cnica}

Uso recomendado como candidato em analise multicriterio, condicionado ao
cumprimento das pre-condicoes acima. Referencia atual de score: GSS 7.7
em 2026-03-05; risco predominante: moderado.

\begin{itemize}
\tightlist
\item
  \begin{enumerate}
  \def\labelenumi{\arabic{enumi})}
  \tightlist
  \item
    Delimitacao territorial validada por georreferencia institucional
    (Fonte: acesso: 2026-03-05).
  \end{enumerate}
\item
  \begin{enumerate}
  \def\labelenumi{\arabic{enumi})}
  \setcounter{enumi}{1}
  \tightlist
  \item
    População municipal/\allowbreak distrital referenciada por base censitaria
    nacional (Fonte: acesso: 2026-03-05).
  \end{enumerate}
\item
  \begin{enumerate}
  \def\labelenumi{\arabic{enumi})}
  \setcounter{enumi}{2}
  \tightlist
  \item
    Comparabilidade intra-departamental apoiada em indicadores
    distritais oficiais (Fonte: acesso: 2026-03-05).
  \end{enumerate}
\item
  \begin{enumerate}
  \def\labelenumi{\arabic{enumi})}
  \setcounter{enumi}{3}
  \tightlist
  \item
    Estado macro de conectividade viaria ancorado em informacoes de
    rodovias (Fonte: acesso: 2026-03-05).
  \end{enumerate}
\item
  \begin{enumerate}
  \def\labelenumi{\arabic{enumi})}
  \setcounter{enumi}{4}
  \tightlist
  \item
    Exposicao hidrometeorologica monitorada por avisos oficiais (Fonte:
    acesso: 2026-03-05).
  \end{enumerate}
\item
  \begin{enumerate}
  \def\labelenumi{\arabic{enumi})}
  \setcounter{enumi}{5}
  \tightlist
  \item
    Historico de contexto meteorologico apoiado em publicacoes tecnicas
    (Fonte: acesso: 2026-03-05).
  \end{enumerate}
\item
  \begin{enumerate}
  \def\labelenumi{\arabic{enumi})}
  \setcounter{enumi}{6}
  \tightlist
  \item
    Capacidade institucional de resposta a eventos apoiada em acoes da
    SEN\footnote{Secretaría de Emergencia Nacional (Secretaria de Emergência Nacional), órgão governamental de resposta a desastres.} (Fonte: acesso: 2026-03-05).
  \end{enumerate}
\item
  \begin{enumerate}
  \def\labelenumi{\arabic{enumi})}
  \setcounter{enumi}{7}
  \tightlist
  \item
    Infraestrutura energetica analisada via operador nacional (Fonte:
    acesso: 2026-03-05).
  \end{enumerate}
\item
  \begin{enumerate}
  \def\labelenumi{\arabic{enumi})}
  \setcounter{enumi}{8}
  \tightlist
  \item
    Referencias de obras e logistica nacional verificadas no MOPC
    (Fonte: acesso: 2026-03-05).
  \end{enumerate}
\item
  \begin{enumerate}
  \def\labelenumi{\arabic{enumi})}
  \setcounter{enumi}{9}
  \tightlist
  \item
    Contexto sociopolitico institucional referenciado por autoridade
    eleitoral (Fonte: acesso: 2026-03-05).
  \end{enumerate}
\item
  \begin{enumerate}
  \def\labelenumi{\arabic{enumi})}
  \setcounter{enumi}{10}
  \tightlist
  \item
    Camada cartografica de apoio eleitoral/\allowbreak territorial consultada
    (Fonte: acesso: 2026-03-05).
  \end{enumerate}
\item
  \begin{enumerate}
  \def\labelenumi{\arabic{enumi})}
  \setcounter{enumi}{11}
  \tightlist
  \item
    Dados publicos complementares para verificacao cruzada (Fonte:
    acesso: 2026-03-05).
  \end{enumerate}
\end{itemize}

\subsubsection{Combustível}\label{vuln-San_Pedro_de_Ycuamandiyu-combustuxedvel}

\textbf{Referência:} PETROPAR /\allowbreak  postos locais (2024) \textbf{Tipo de
localidade:} Capital departamental

\begin{longtable}[]{@{}lll@{}}
\toprule\noalign{}
Combustível & USD/\allowbreak litro & Gs/\allowbreak litro (aprox.) \\
\midrule\noalign{}
\endhead
\bottomrule\noalign{}
\endlastfoot
Gasolina 93 oct & 1.00 & 7,400 \\
Gasolina 97 oct (premium) & 1.10 & 8,140 \\
Diesel & 0.93 & 6,882 \\
\end{longtable}

\begin{quote}
Preços podem variar ±5\% conforme posto e sazonalidade. Chaco e interior
remoto apresentam maior variação.
\end{quote}

\subsubsection{Cobertura Celular}\label{vuln-San_Pedro_de_Ycuamandiyu-cobertura-celular}

\textbf{Fonte:} CONATEL PY /\allowbreak  operadoras (2024)

\begin{longtable}[]{@{}ll@{}}
\toprule\noalign{}
Parâmetro & Valor \\
\midrule\noalign{}
\endhead
\bottomrule\noalign{}
\endlastfoot
Cobertura 4G população (dept.) & 80\% \\
Cobertura 4G área rural & 58\% \\
Melhor operadora & Tigo \\
Qualidade rural & moderada \\
\end{longtable}

\begin{quote}
Para áreas rurais fora do núcleo urbano, recomenda-se chip Tigo como
principal e Personal como backup.
\end{quote}

\subsubsection{Internet}\label{vuln-San_Pedro_de_Ycuamandiyu-internet}

\textbf{Fonte:} CONATEL /\allowbreak  Speedtest Ookla (2024)

\begin{longtable}[]{@{}ll@{}}
\toprule\noalign{}
Parâmetro & Valor \\
\midrule\noalign{}
\endhead
\bottomrule\noalign{}
\endlastfoot
Velocidade média download & 35 Mbps \\
Domicílios com internet (dept.) & 45\% \\
Tecnologia predominante & rádio \\
Opção rural & Starlink disponível (\textasciitilde USD 44/\allowbreak mês) \\
\end{longtable}

\subsubsection{Mercado Imobiliário e Terra
Rural}\label{vuln-San_Pedro_de_Ycuamandiyu-mercado-imobiliuxe1rio-e-terra-rural}

\textbf{Fonte:} INDERT /\allowbreak  Clasificados.com.py (2024)

\begin{longtable}[]{@{}ll@{}}
\toprule\noalign{}
Tipo & Referência \\
\midrule\noalign{}
\endhead
\bottomrule\noalign{}
\endlastfoot
Terra agrícola alta prod. (USD/\allowbreak ha) & 4,600 \\
Imóvel urbano (USD/\allowbreak m²) & 690 \\
Aluguel 2 quartos (USD/\allowbreak mês) & 300 \\
\end{longtable}

\begin{quote}
Valores de referência departamental. Localidades menores podem ter
preços 20--40\% abaixo da capital departamental. \#\#\# Saúde
\end{quote}

\textbf{Fonte:} MSPBS /\allowbreak  IPS Paraguay (2024-2026), consolidação
departamental e proxy local

\begin{longtable}[]{@{}ll@{}}
\toprule\noalign{}
Serviço & Disponibilidade \\
\midrule\noalign{}
\endhead
\bottomrule\noalign{}
\endlastfoot
USF /\allowbreak  Posto de Saúde & sim \\
Hospital Regional & sim \\
IPS (seguro social) & sim \\
Farmácia & sim \\
Distância ao hospital de referência & local \\
\end{longtable}

\textbf{Principais estabelecimentos:} Hospital distrital /\allowbreak  regional de
San Pedro de Ycuamandiyu; rede de USF e farmacias locais

\textbf{Observação para imigrantes:} Centro de referencia do
departamento. Acesso a atendimento primario e, em geral, a maior oferta
publica e privada da area. Cobertura privada continua recomendada para
especialidades e urgências de maior complexidade.
\subsubsection{Dossiê de Campo}\label{dossie-San_Pedro_de_Ycuamandiyu-dossiuxea-de-campo}
\begin{description}[style=nextline,leftmargin=0.5cm,font=\bfseries]
\item[Coordenadas]  24°06'S 57°05'W (geocodificado via Nominatim).
\item[Irradiância solar global — ALLSKY\_SFC\_SW\_DWN (kWh/\allowbreak m²/\allowbreak dia)]  Jan 6.64 | Fev 6.13 | Mar 5.48 | Abr 4.54 | Mai 3.38 | Jun 2.87 | Jul 3.21 | Ago 3.93 | Set 4.62 | Out 5.45 | Nov 6.37 | Dez 6.67. Média anual: 4.93 kWh/\allowbreak m²/\allowbreak dia. Mínimo: Jun (2.87) | Máximo: Dez (6.67). Fonte: NASA POWER Climatology API, período 2001-2020 (acesso em 2026-03-20).
\item[Precipitação — PRECTOTCORR (mm/\allowbreak mês)]  Jan 4.39 | Fev 5.38 | Mar 3.67 | Abr 4.73 | Mai 4.50 | Jun 2.66 | Jul 1.89 | Ago 1.18 | Set 2.77 | Out 5.36 | Nov 6.53 | Dez 5.43. Média anual: \textasciitilde 1.471 mm/\allowbreak ano. Estação chuvosa Out-Mai; seca Jun-Set (mínimo em Ago). Fonte: NASA POWER (acesso em 2026-03-20).
\item[Inclinação recomendada para placas solares]  24° voltado para o Norte (anual); 34° no inverno (Jun-Ago); 14° no verão (Nov-Jan). Base: latitude local -24.10°S.
\end{description}
\paragraph{Solo (SoilGrids 2.0, média ponderada 0--30
cm)}

\begin{longtable}[]{@{}ll@{}}
\toprule\noalign{}
Parâmetro & Valor \\
\midrule\noalign{}
\endhead
\bottomrule\noalign{}
\endlastfoot
pH (H$_2$O) & 5.7 \\
Carbono orgânico (SOC) & 17.03 g/\allowbreak kg \\
Argila & 26.45 \% \\
Areia & 49.98 \% \\
Silte (calc.) & 23.6 \% \\
Densidade aparente & 1.31 g/\allowbreak cm³ \\
\textbf{Aptidão agrícola} & \textbf{Alta} \\
\end{longtable}

Fonte: ISRIC SoilGrids 2.0 via WCS (acesso em 2026-03-21). Coords:
-24.1°, -57.0833°. Média ponderada camadas 0-5, 5-15, 15-30 cm.
\newpage
\secaoDiagnostico{Santa Rosa del Aguaray}{\mapaConteudo{-23.8833}{-56.4833}}{Santa Rosa del Aguaray é a joia produtiva do Norte paraguaio. Oferece a
melhor combinação de infraestrutura hospitalar, qualidade de solo e
conectividade rodoviária de toda a região. Sua viabilidade como refúgio
de nível superior é baseada na abundância de recursos e na força
institucional, contrabalançando os desafios de segurança tática do
departamento.

\subsubsection{Por que sim}

\begin{itemize}
\tightlist
\item
  Hospital de alta complexidade moderno e bem equipado.
\item
  Solos de altíssima produtividade (``Terra Vermelha'').
\item
  Abundância de águas puras (Laguna Blanca e aquíferos).
\item
  Sem restrições legais para proprietários estrangeiros.
\end{itemize}

\subsubsection{Por que não}

\begin{itemize}
\tightlist
\item
  Hub logístico de alta visibilidade estratégica em cenários de
  conflito.
\item
  Presença de riscos sociais/\allowbreak insurgentes no departamento exige
  vigilância.
\item
  Crescimento populacional urbano acelerado.
\end{itemize}

\subsubsection{Recomendação
Técnica}

Ideal para estabelecimentos que buscam integrar alta produtividade
agrícola com proximidade a serviços essenciais de saúde. O GSS de 6.6
reflete a excepcional base de recursos naturais e institucionais.
Referencia atual de score: GSS 6.6 em 2026-03-05.

\subsubsection{}

\begin{itemize}
\tightlist
\item
  \begin{enumerate}
  \def\labelenumi{\arabic{enumi})}
  \tightlist
  \item
    Dados censitários (INE\footnote{Instituto Nacional de Estadística (Instituto Nacional de Estatística), órgão oficial de dados demográficos e censitários.} 2022).
  \end{enumerate}
\item
  \begin{enumerate}
  \def\labelenumi{\arabic{enumi})}
  \setcounter{enumi}{1}
  \tightlist
  \item
    Relatórios de cooperação técnica KOICA/\allowbreak MSPBS.
  \end{enumerate}
\item
  \begin{enumerate}
  \def\labelenumi{\arabic{enumi})}
  \setcounter{enumi}{2}
  \tightlist
  \item
    Mapa de uso da terra e pedologia (MAG).
  \end{enumerate}
\item
  \begin{enumerate}
  \def\labelenumi{\arabic{enumi})}
  \setcounter{enumi}{3}
  \tightlist
  \item
    Dados turísticos e ambientais da Laguna Blanca (MADES).
  \end{enumerate}
\end{itemize}}
\begin{table}[ht!]
\small \begin{tabular}{p{3.0cm}p{0.8cm}p{6.0cm}} \toprule
\textbf{Critério} & \textbf{Nota} & \textbf{Justificativa} \\ \midrule
A - Ameaças Estratégicas & 5.0 & Hub logístico principal e alvo estratégico de saúde; alta visibilidade tática \\
B - Risco Social & 6.0 & População em rápido crescimento com riscos sociais regionais conhecidos \\
C - Riscos Naturais & 7.0 & Geologia estável e riscos naturais climáticos gerenciáveis \\
D - Autossuficiência & 8.5 & Riqueza de recursos excepcional: solos de elite, gergelim, grãos e água pura \\
E - Institucional & 7.0 & Forte presença institucional e apoio técnico internacional; liberdade jurídica total \\
\midrule \textbf{GSS FINAL} & \textbf{6.6} & \textbf{Moderadamente Seguro (Altamente recomendado como base produtiva e de serviços para o Norte paraguaio).} \\ \bottomrule \end{tabular}
\end{table}
\subsubsection{Vulnerabilidades e
Mitigação}\label{vuln-Santa_Rosa_del_Aguaray-vulnerabilidades-e-mitigauxe7uxe3o}

\begin{itemize}
\tightlist
\item
  \textbf{Exposição Tática:} O hospital e o cruzamento rodoviário são
  alvos logísticos em crises (Mitigação: residência rural afastada 15-20
  km do nó urbano).
\item
  \textbf{Dependência de Exportação:} Economia vinculada a commodities
  (Mitigação: manutenção de diversidade produtiva para subsistência).
\item
  \textbf{Conflitos Agrários Regionais:} Vigilância necessária em
  grandes extensões rurais remotas (Mitigação: segurança privada e
  integração comunitária).
\end{itemize}
\subsubsection{Dossiê de Campo}\label{dossie-Santa_Rosa_del_Aguaray-dossiuxea-de-campo}
\begin{description}[style=nextline,leftmargin=0.5cm,font=\bfseries]
\item[Coordenadas]  23°53'S 56°29'W.
\item[Topografia]  Relevo predominantemente plano a ondulado. Localizado no cruzamento estratégico das Rutas PY08 e PY11. Geologicamente estável, risco sísmico praticamente nulo.
\item[Alvos Estratégicos]  Hospital Geral de Santa Rosa del Aguaray (Referência regional de alta complexidade); Cruzamento das Rutas PY08 (Norte-Sul) e PY11 (Leste-Oeste); Hub logístico vital para Amambay e Concepción.
\item[Fallout]  Ventos predominantes do Norte (quentes e úmidos) e Sul (frentes frias). Baixo risco de fallout direto de alvos continentais primários.
\item[População]  39.643 habitantes (Censo 2022 Final). Um dos distritos com maior taxa de crescimento anual (3,3\%) no interior. Perfil 80\% rural.
\item[Densidade]  \textasciitilde 26,0 hab/\allowbreak km² (Moderada, com núcleo urbano comercial denso sobre as rodovias).
\item[Vias de Acesso]  Ponto de conexão vital entre Pedro Juan Caballero, Concepción e Assunção. Infraestrutura asfáltica de excelência, servindo como base de suprimentos para todo o Norte de San Pedro.
\item[Serviços]  Melhor infraestrutura hospitalar do interior paraguaio (Hospital Paraguai-Coreia/\allowbreak KOICA). Presença consolidada de bancos, silos e comércio agroindustrial.
\item[Custo de Vida]  Baixo (30-40\% inferior a centros urbanos brasileiros similares); economia pujante baseada na exportação agrícola.
\item[Preço da Terra]  US\$ 4.000-8.000/\allowbreak ha (terras agrícolas mecanizadas); US\$ 1.500-3.500/\allowbreak ha (áreas de pecuária/\allowbreak mistas).
\item[Hidrologia]  Baixo risco de inundações sistêmicas catastróficas. Áreas ribeirinhas baixas do Rio Aguaray Guazú vulneráveis em tempestades severas (precipitação média 1.300-1.500 mm).
\item[Clima]  Subtropical Úmido. Verões intensos; invernos com entradas bruscas de ar frio.
\item[Energia]  Rede nacional (Itaipu); subestação regional local garante maior prioridade em reparos técnicos.
\item[Água]  Abundância hídrica excepcional; proximidade com a Reserva Laguna Blanca (águas cristalinas sobre areia calcária).
\item[Qualidade do Solo]  Latossolo Vermelho (Terra Vermelha); profundo e muito fértil, excelente para agricultura mecanizada intensiva.
\item[Resources Locais]  Capital Nacional do Sésamo (Gergelim); polo produtor de soja, milho, chia e trigo. Grande excedente de proteína animal e grãos.
\item[Segurança]  Zona sob Vigilância. Historicamente inserida em região com presença de grupos insurgentes (EPP\footnote{Ejército del Pueblo Paraguayo (Exército do Povo Paraguaio), grupo insurgente marxista-leninista que atua principalmente no norte do país.}), o que motiva patrulhamento ostensivo da FTC e da Polícia Nacional. O núcleo urbano é um polo comercial vigiado e relativamente seguro.
\item[Leis Local: Município em rápido desenvolvimento e com forte apoio institucional internacional (Coreia do Sul). Livre de Restrição de Fronteira]  Fora da faixa de 50km da Lei 2532/\allowbreak 05, sem restrições para brasileiros.
\end{description}
\paragraph{Solo (SoilGrids 2.0, média ponderada 0--30
cm)}

\begin{longtable}[]{@{}ll@{}}
\toprule\noalign{}
Parâmetro & Valor \\
\midrule\noalign{}
\endhead
\bottomrule\noalign{}
\endlastfoot
pH (H$_2$O) & 5.4 \\
Carbono orgânico (SOC) & 17.47 g/\allowbreak kg \\
Argila & 24.2 \% \\
Areia & 57.97 \% \\
Silte (calc.) & 17.8 \% \\
Densidade aparente & 1.29 g/\allowbreak cm³ \\
\textbf{Aptidão agrícola} & \textbf{Média} \\
\end{longtable}

Fonte: ISRIC SoilGrids 2.0 via WCS (acesso em 2026-03-21). Coords:
-23.8833°, -56.4833°. Média ponderada camadas 0-5, 5-15, 15-30 cm.
\newpage
\secaoDiagnostico{San Vicente Pancholo}{\mapaConteudo{-22.4666}{-56.8000}}{San Vicente Pancholo (San Pedro) apresenta perfil operacional
consolidado para comparacao territorial, com leitura conjunta de
infraestrutura, risco hidroclimatico e capacidade de continuidade de
servicos essenciais.

\subsubsection{Por que sim}

\begin{itemize}
\tightlist
\item
  Base institucional de referencia disponivel e rastreavel.
\item
  Estrutura comparativa (A-E/\allowbreak GSS) consistente para decisao relativa.
\item
  Plano de mitigacao acionavel ja definido no dossie.
\end{itemize}

\subsubsection{Por que não}

\begin{itemize}
\tightlist
\item
  Serie oficial distrital granular ainda pode apresentar lacunas
  temporais.
\item
  Sensibilidade do score exige monitoramento continuo de risco e
  conectividade.
\item
  Decisao final depende de validacao de campo e janela temporal recente.
\end{itemize}

\subsubsection{Pre-condicoes Minimas de
Decisao}

\begin{itemize}
\tightlist
\item
  Atualizar indicadores criticos em ciclo trimestral.
\item
  Confirmar trafegabilidade e contingencia hidrica/\allowbreak energetica local.
\item
  Validar checklist de resiliencia domiciliar/\allowbreak logistica antes de decisao
  final.
\end{itemize}

\subsubsection{Recomendação
Técnica}

Uso recomendado como candidato em analise multicriterio, condicionado ao
cumprimento das pre-condicoes acima. Referencia atual de score: GSS 8.3
em 2026-03-05; risco predominante: baixo.

\subsubsection{}

\begin{itemize}
\tightlist
\item
  \begin{enumerate}
  \def\labelenumi{\arabic{enumi})}
  \tightlist
  \item
    Delimitacao territorial validada por georreferencia institucional
    (Fonte: acesso: 2026-03-05).
  \end{enumerate}
\item
  \begin{enumerate}
  \def\labelenumi{\arabic{enumi})}
  \setcounter{enumi}{1}
  \tightlist
  \item
    População municipal/\allowbreak distrital referenciada por base censitaria
    nacional (Fonte: acesso: 2026-03-05).
  \end{enumerate}
\item
  \begin{enumerate}
  \def\labelenumi{\arabic{enumi})}
  \setcounter{enumi}{2}
  \tightlist
  \item
    Comparabilidade intra-departamental apoiada em indicadores
    distritais oficiais (Fonte: acesso: 2026-03-05).
  \end{enumerate}
\item
  \begin{enumerate}
  \def\labelenumi{\arabic{enumi})}
  \setcounter{enumi}{3}
  \tightlist
  \item
    Estado macro de conectividade viaria ancorado em informacoes de
    rodovias (Fonte: acesso: 2026-03-05).
  \end{enumerate}
\item
  \begin{enumerate}
  \def\labelenumi{\arabic{enumi})}
  \setcounter{enumi}{4}
  \tightlist
  \item
    Exposicao hidrometeorologica monitorada por avisos oficiais (Fonte:
    acesso: 2026-03-05).
  \end{enumerate}
\item
  \begin{enumerate}
  \def\labelenumi{\arabic{enumi})}
  \setcounter{enumi}{5}
  \tightlist
  \item
    Historico de contexto meteorologico apoiado em publicacoes tecnicas
    (Fonte: acesso: 2026-03-05).
  \end{enumerate}
\item
  \begin{enumerate}
  \def\labelenumi{\arabic{enumi})}
  \setcounter{enumi}{6}
  \tightlist
  \item
    Capacidade institucional de resposta a eventos apoiada em acoes da
    SEN\footnote{Secretaría de Emergencia Nacional (Secretaria de Emergência Nacional), órgão governamental de resposta a desastres.} (Fonte: acesso: 2026-03-05).
  \end{enumerate}
\item
  \begin{enumerate}
  \def\labelenumi{\arabic{enumi})}
  \setcounter{enumi}{7}
  \tightlist
  \item
    Infraestrutura energetica analisada via operador nacional (Fonte:
    acesso: 2026-03-05).
  \end{enumerate}
\item
  \begin{enumerate}
  \def\labelenumi{\arabic{enumi})}
  \setcounter{enumi}{8}
  \tightlist
  \item
    Referencias de obras e logistica nacional verificadas no MOPC
    (Fonte: acesso: 2026-03-05).
  \end{enumerate}
\item
  \begin{enumerate}
  \def\labelenumi{\arabic{enumi})}
  \setcounter{enumi}{9}
  \tightlist
  \item
    Contexto sociopolitico institucional referenciado por autoridade
    eleitoral (Fonte: acesso: 2026-03-05).
  \end{enumerate}
\item
  \begin{enumerate}
  \def\labelenumi{\arabic{enumi})}
  \setcounter{enumi}{10}
  \tightlist
  \item
    Camada cartografica de apoio eleitoral/\allowbreak territorial consultada
    (Fonte: acesso: 2026-03-05).
  \end{enumerate}
\item
  \begin{enumerate}
  \def\labelenumi{\arabic{enumi})}
  \setcounter{enumi}{11}
  \tightlist
  \item
    Dados publicos complementares para verificacao cruzada (Fonte:
    acesso: 2026-03-05).
  \end{enumerate}
\end{itemize}}
\begin{table}[ht!]
\small \begin{tabular}{p{3.0cm}p{0.8cm}p{6.0cm}} \toprule
\textbf{Critério} & \textbf{Nota} & \textbf{Justificativa} \\ \midrule
A - Ameaças Estratégicas & 9.2 & - \\
B - Risco Social & 8.1 & - \\
C - Riscos Naturais & 8.3 & - \\
D - Autossuficiência & 7.5 & - \\
E - Institucional & 8.4 & - \\
\midrule \textbf{GSS FINAL} & \textbf{8.3} & \textbf{Seguro (bom potencial com preparacao).} \\ \bottomrule \end{tabular}
\end{table}
\subsubsection{Vulnerabilidades e
Mitigação}\label{vuln-San_Vicente_Pancholo-vulnerabilidades-e-mitigauxe7uxe3o}

\begin{itemize}
\tightlist
\item
  Exposicao hidrometeorologica controlada (\textless45/\allowbreak 100).
\item
  Conectividade viaria favoravel (\textgreater=70/\allowbreak 100).
\item
  Estoque de contingencia para 14 dias (agua/\allowbreak alimentos/\allowbreak combustivel),
  priorizado quando risco hidro=21/\allowbreak 100.
\item
  Duplicar rotas/\allowbreak logistica quando conectividade viaria=93/\allowbreak 100 for
  \textless70; manter rota secundaria mapeada e testada trimestralmente.
\item
  Plano de energia resiliente com autonomia minima de 72h
  (geracao/\allowbreak backup), revisao semestral.
\end{itemize}

Sintese aplicada: - Dados textuais consolidados com base nas fontes
oficiais acima; proximos ciclos podem aprofundar indicadores numericos
especificos por distrito.

\begin{itemize}
\tightlist
\item
  População municipal/\allowbreak distrital (consolidado operacional): 118132
  habitantes (referencia temporal: 2026-03-05; base: consolidacao de
  fontes oficiais INE\footnote{Instituto Nacional de Estadística (Instituto Nacional de Estatística), órgão oficial de dados demográficos e censitários.} listadas na secao 8).
\item
  Densidade demografica (consolidado operacional): 9 hab/\allowbreak km2 (referencia
  temporal: 2026-03-05; base: consolidacao territorial oficial).
\item
  Conectividade viaria funcional (consolidado operacional): 93/\allowbreak 100
  (referencia temporal: 2026-03-05; base: MOPC estado de rutas).
\item
  Exposicao hidrometeorologica relativa (consolidado operacional):
  21/\allowbreak 100 (referencia temporal: 2026-03-05; base: DMH/\allowbreak SEN\footnote{Secretaría de Emergencia Nacional (Secretaria de Emergência Nacional), órgão governamental de resposta a desastres.} avisos e
  ocorrencias).
\end{itemize}

NOTA METODOLOGICA: Indicadores consolidados para comparabilidade
distrital nesta fase; quando serie oficial granular não estiver publica,
registrar lacuna oficial e manter rastreabilidade da fonte
institucional. - Cenario curto prazo: interrupcoes logisticas por
eventos climaticos. - Cenario medio prazo: pressao sobre servicos
essenciais. - Mitigacao: redundancia de rota, agua, energia e estoque.

\begin{itemize}
\tightlist
\item
  Cenario otimista (B+1, C+1): GSS 8.7
\item
  Cenario conservador (B-1, C-1): GSS 8.0
\item
  Amplitude de sensibilidade: 0.7 ponto(s)
\item
  Leitura: variacao controlada para decisao comparativa entre distritos;
  priorizar monitoramento de B e C.
\end{itemize}

\subsubsection{Diagnostico Integrado}\label{vuln-San_Vicente_Pancholo-diagnostico-integrado}

San Vicente Pancholo (San Pedro) apresenta perfil operacional
consolidado para comparacao territorial, com leitura conjunta de
infraestrutura, risco hidroclimatico e capacidade de continuidade de
servicos essenciais.

\subsubsection{Por que sim}\label{vuln-San_Vicente_Pancholo-por-que-sim}

\begin{itemize}
\tightlist
\item
  Base institucional de referencia disponivel e rastreavel.
\item
  Estrutura comparativa (A-E/\allowbreak GSS) consistente para decisao relativa.
\item
  Plano de mitigacao acionavel ja definido no dossie.
\end{itemize}

\subsubsection{Por que não}\label{vuln-San_Vicente_Pancholo-por-que-nuxe3o}

\begin{itemize}
\tightlist
\item
  Serie oficial distrital granular ainda pode apresentar lacunas
  temporais.
\item
  Sensibilidade do score exige monitoramento continuo de risco e
  conectividade.
\item
  Decisao final depende de validacao de campo e janela temporal recente.
\end{itemize}

\subsubsection{Pre-condicoes Minimas de
Decisao}\label{vuln-San_Vicente_Pancholo-pre-condicoes-minimas-de-decisao}

\begin{itemize}
\tightlist
\item
  Atualizar indicadores criticos em ciclo trimestral.
\item
  Confirmar trafegabilidade e contingencia hidrica/\allowbreak energetica local.
\item
  Validar checklist de resiliencia domiciliar/\allowbreak logistica antes de decisao
  final.
\end{itemize}

\subsubsection{Recomendação
Técnica}\label{vuln-San_Vicente_Pancholo-recomendauxe7uxe3o-tuxe9cnica}

Uso recomendado como candidato em analise multicriterio, condicionado ao
cumprimento das pre-condicoes acima. Referencia atual de score: GSS 8.3
em 2026-03-05; risco predominante: baixo.

\begin{itemize}
\tightlist
\item
  \begin{enumerate}
  \def\labelenumi{\arabic{enumi})}
  \tightlist
  \item
    Delimitacao territorial validada por georreferencia institucional
    (Fonte: acesso: 2026-03-05).
  \end{enumerate}
\item
  \begin{enumerate}
  \def\labelenumi{\arabic{enumi})}
  \setcounter{enumi}{1}
  \tightlist
  \item
    População municipal/\allowbreak distrital referenciada por base censitaria
    nacional (Fonte: acesso: 2026-03-05).
  \end{enumerate}
\item
  \begin{enumerate}
  \def\labelenumi{\arabic{enumi})}
  \setcounter{enumi}{2}
  \tightlist
  \item
    Comparabilidade intra-departamental apoiada em indicadores
    distritais oficiais (Fonte: acesso: 2026-03-05).
  \end{enumerate}
\item
  \begin{enumerate}
  \def\labelenumi{\arabic{enumi})}
  \setcounter{enumi}{3}
  \tightlist
  \item
    Estado macro de conectividade viaria ancorado em informacoes de
    rodovias (Fonte: acesso: 2026-03-05).
  \end{enumerate}
\item
  \begin{enumerate}
  \def\labelenumi{\arabic{enumi})}
  \setcounter{enumi}{4}
  \tightlist
  \item
    Exposicao hidrometeorologica monitorada por avisos oficiais (Fonte:
    acesso: 2026-03-05).
  \end{enumerate}
\item
  \begin{enumerate}
  \def\labelenumi{\arabic{enumi})}
  \setcounter{enumi}{5}
  \tightlist
  \item
    Historico de contexto meteorologico apoiado em publicacoes tecnicas
    (Fonte: acesso: 2026-03-05).
  \end{enumerate}
\item
  \begin{enumerate}
  \def\labelenumi{\arabic{enumi})}
  \setcounter{enumi}{6}
  \tightlist
  \item
    Capacidade institucional de resposta a eventos apoiada em acoes da
    SEN\footnote{Secretaría de Emergencia Nacional (Secretaria de Emergência Nacional), órgão governamental de resposta a desastres.} (Fonte: acesso: 2026-03-05).
  \end{enumerate}
\item
  \begin{enumerate}
  \def\labelenumi{\arabic{enumi})}
  \setcounter{enumi}{7}
  \tightlist
  \item
    Infraestrutura energetica analisada via operador nacional (Fonte:
    acesso: 2026-03-05).
  \end{enumerate}
\item
  \begin{enumerate}
  \def\labelenumi{\arabic{enumi})}
  \setcounter{enumi}{8}
  \tightlist
  \item
    Referencias de obras e logistica nacional verificadas no MOPC
    (Fonte: acesso: 2026-03-05).
  \end{enumerate}
\item
  \begin{enumerate}
  \def\labelenumi{\arabic{enumi})}
  \setcounter{enumi}{9}
  \tightlist
  \item
    Contexto sociopolitico institucional referenciado por autoridade
    eleitoral (Fonte: acesso: 2026-03-05).
  \end{enumerate}
\item
  \begin{enumerate}
  \def\labelenumi{\arabic{enumi})}
  \setcounter{enumi}{10}
  \tightlist
  \item
    Camada cartografica de apoio eleitoral/\allowbreak territorial consultada
    (Fonte: acesso: 2026-03-05).
  \end{enumerate}
\item
  \begin{enumerate}
  \def\labelenumi{\arabic{enumi})}
  \setcounter{enumi}{11}
  \tightlist
  \item
    Dados publicos complementares para verificacao cruzada (Fonte:
    acesso: 2026-03-05).
  \end{enumerate}
\end{itemize}

\subsubsection{Combustível}\label{vuln-San_Vicente_Pancholo-combustuxedvel}

\textbf{Referência:} PETROPAR /\allowbreak  postos locais (2024) \textbf{Tipo de
localidade:} Interior

\begin{longtable}[]{@{}lll@{}}
\toprule\noalign{}
Combustível & USD/\allowbreak litro & Gs/\allowbreak litro (aprox.) \\
\midrule\noalign{}
\endhead
\bottomrule\noalign{}
\endlastfoot
Gasolina 93 oct & 1.00 & 7,400 \\
Gasolina 97 oct (premium) & 1.10 & 8,140 \\
Diesel & 0.93 & 6,882 \\
\end{longtable}

\begin{quote}
Preços podem variar ±5\% conforme posto e sazonalidade. Chaco e interior
remoto apresentam maior variação.
\end{quote}

\subsubsection{Cobertura Celular}\label{vuln-San_Vicente_Pancholo-cobertura-celular}

\textbf{Fonte:} CONATEL PY /\allowbreak  operadoras (2024)

\begin{longtable}[]{@{}ll@{}}
\toprule\noalign{}
Parâmetro & Valor \\
\midrule\noalign{}
\endhead
\bottomrule\noalign{}
\endlastfoot
Cobertura 4G população (dept.) & 80\% \\
Cobertura 4G área rural & 58\% \\
Melhor operadora & Tigo \\
Qualidade rural & moderada \\
\end{longtable}

\begin{quote}
Para áreas rurais fora do núcleo urbano, recomenda-se chip Tigo como
principal e Personal como backup.
\end{quote}

\subsubsection{Internet}\label{vuln-San_Vicente_Pancholo-internet}

\textbf{Fonte:} CONATEL /\allowbreak  Speedtest Ookla (2024)

\begin{longtable}[]{@{}ll@{}}
\toprule\noalign{}
Parâmetro & Valor \\
\midrule\noalign{}
\endhead
\bottomrule\noalign{}
\endlastfoot
Velocidade média download & 35 Mbps \\
Domicílios com internet (dept.) & 45\% \\
Tecnologia predominante & rádio \\
Opção rural & Starlink disponível (\textasciitilde USD 44/\allowbreak mês) \\
\end{longtable}

\subsubsection{Mercado Imobiliário e Terra
Rural}\label{vuln-San_Vicente_Pancholo-mercado-imobiliuxe1rio-e-terra-rural}

\textbf{Fonte:} INDERT /\allowbreak  Clasificados.com.py (2024)

\begin{longtable}[]{@{}ll@{}}
\toprule\noalign{}
Tipo & Referência \\
\midrule\noalign{}
\endhead
\bottomrule\noalign{}
\endlastfoot
Terra agrícola alta prod. (USD/\allowbreak ha) & 4,600 \\
Imóvel urbano (USD/\allowbreak m²) & 600 \\
Aluguel 2 quartos (USD/\allowbreak mês) & 250 \\
\end{longtable}

\begin{quote}
Valores de referência departamental. Localidades menores podem ter
preços 20--40\% abaixo da capital departamental. \#\#\# Saúde
\end{quote}

\textbf{Fonte:} MSPBS /\allowbreak  IPS Paraguay (2024-2026), consolidação
departamental e proxy local

\begin{longtable}[]{@{}
  >{\raggedright\arraybackslash}p{(\linewidth - 2\tabcolsep) * \real{0.3600}}
  >{\raggedright\arraybackslash}p{(\linewidth - 2\tabcolsep) * \real{0.6400}}@{}}
\toprule\noalign{}
\begin{minipage}[b]{\linewidth}\raggedright
Serviço
\end{minipage} & \begin{minipage}[b]{\linewidth}\raggedright
Disponibilidade
\end{minipage} \\
\midrule\noalign{}
\endhead
\bottomrule\noalign{}
\endlastfoot
USF /\allowbreak  Posto de Saúde & sim \\
Hospital Regional & não \\
IPS (seguro social) & não \\
Farmácia & sim \\
Distância ao hospital de referência & \textasciitilde89 km (San Pedro de
Ycuamandiyu) \\
\end{longtable}

\textbf{Principais estabelecimentos:} USF local; referencia hospitalar
em San Pedro de Ycuamandiyu

\textbf{Observação para imigrantes:} Atendimento primario local ou em
raio curto; casos de maior complexidade seguem para San Pedro de
Ycuamandiyu. Cobertura privada continua recomendada para especialidades
e urgências de maior complexidade.
\subsubsection{Dossiê de Campo}\label{dossie-San_Vicente_Pancholo-dossiuxea-de-campo}
\begin{description}[style=nextline,leftmargin=0.5cm,font=\bfseries]
\item[Coordenadas]  22°28'S 56°48'W (geocodificado via Nominatim).
\end{description}
\paragraph{Solo (SoilGrids 2.0, média ponderada 0--30
cm)}

\begin{longtable}[]{@{}ll@{}}
\toprule\noalign{}
Parâmetro & Valor \\
\midrule\noalign{}
\endhead
\bottomrule\noalign{}
\endlastfoot
pH (H$_2$O) & 5.5 \\
Carbono orgânico (SOC) & 13.52 g/\allowbreak kg \\
Argila & 18.56 \% \\
Areia & 61.08 \% \\
Silte (calc.) & 20.4 \% \\
Densidade aparente & 1.34 g/\allowbreak cm³ \\
\textbf{Aptidão agrícola} & \textbf{Média} \\
\end{longtable}

Fonte: ISRIC SoilGrids 2.0 via WCS (acesso em 2026-03-21). Coords:
-22.4667°, -56.8°. Média ponderada camadas 0-5, 5-15, 15-30 cm.
\newpage
\secaoDiagnostico{Tacuati}{\mapaConteudo{-23.4500}{-56.7333}}{Tacuati (San Pedro) apresenta perfil operacional consolidado para
comparacao territorial, com leitura conjunta de infraestrutura, risco
hidroclimatico e capacidade de continuidade de servicos essenciais.

\subsubsection{Por que sim}

\begin{itemize}
\tightlist
\item
  Base institucional de referencia disponivel e rastreavel.
\item
  Estrutura comparativa (A-E/\allowbreak GSS) consistente para decisao relativa.
\item
  Plano de mitigacao acionavel ja definido no dossie.
\end{itemize}

\subsubsection{Por que não}

\begin{itemize}
\tightlist
\item
  Serie oficial distrital granular ainda pode apresentar lacunas
  temporais.
\item
  Sensibilidade do score exige monitoramento continuo de risco e
  conectividade.
\item
  Decisao final depende de validacao de campo e janela temporal recente.
\end{itemize}

\subsubsection{Pre-condicoes Minimas de
Decisao}

\begin{itemize}
\tightlist
\item
  Atualizar indicadores criticos em ciclo trimestral.
\item
  Confirmar trafegabilidade e contingencia hidrica/\allowbreak energetica local.
\item
  Validar checklist de resiliencia domiciliar/\allowbreak logistica antes de decisao
  final.
\end{itemize}

\subsubsection{Recomendação
Técnica}

Uso recomendado como candidato em analise multicriterio, condicionado ao
cumprimento das pre-condicoes acima. Referencia atual de score: GSS 7.8
em 2026-03-05; risco predominante: moderado.

\subsubsection{}

\begin{itemize}
\tightlist
\item
  \begin{enumerate}
  \def\labelenumi{\arabic{enumi})}
  \tightlist
  \item
    Delimitacao territorial validada por georreferencia institucional
    (Fonte: acesso: 2026-03-05).
  \end{enumerate}
\item
  \begin{enumerate}
  \def\labelenumi{\arabic{enumi})}
  \setcounter{enumi}{1}
  \tightlist
  \item
    População municipal/\allowbreak distrital referenciada por base censitaria
    nacional (Fonte: acesso: 2026-03-05).
  \end{enumerate}
\item
  \begin{enumerate}
  \def\labelenumi{\arabic{enumi})}
  \setcounter{enumi}{2}
  \tightlist
  \item
    Comparabilidade intra-departamental apoiada em indicadores
    distritais oficiais (Fonte: acesso: 2026-03-05).
  \end{enumerate}
\item
  \begin{enumerate}
  \def\labelenumi{\arabic{enumi})}
  \setcounter{enumi}{3}
  \tightlist
  \item
    Estado macro de conectividade viaria ancorado em informacoes de
    rodovias (Fonte: acesso: 2026-03-05).
  \end{enumerate}
\item
  \begin{enumerate}
  \def\labelenumi{\arabic{enumi})}
  \setcounter{enumi}{4}
  \tightlist
  \item
    Exposicao hidrometeorologica monitorada por avisos oficiais (Fonte:
    acesso: 2026-03-05).
  \end{enumerate}
\item
  \begin{enumerate}
  \def\labelenumi{\arabic{enumi})}
  \setcounter{enumi}{5}
  \tightlist
  \item
    Historico de contexto meteorologico apoiado em publicacoes tecnicas
    (Fonte: acesso: 2026-03-05).
  \end{enumerate}
\item
  \begin{enumerate}
  \def\labelenumi{\arabic{enumi})}
  \setcounter{enumi}{6}
  \tightlist
  \item
    Capacidade institucional de resposta a eventos apoiada em acoes da
    SEN\footnote{Secretaría de Emergencia Nacional (Secretaria de Emergência Nacional), órgão governamental de resposta a desastres.} (Fonte: acesso: 2026-03-05).
  \end{enumerate}
\item
  \begin{enumerate}
  \def\labelenumi{\arabic{enumi})}
  \setcounter{enumi}{7}
  \tightlist
  \item
    Infraestrutura energetica analisada via operador nacional (Fonte:
    acesso: 2026-03-05).
  \end{enumerate}
\item
  \begin{enumerate}
  \def\labelenumi{\arabic{enumi})}
  \setcounter{enumi}{8}
  \tightlist
  \item
    Referencias de obras e logistica nacional verificadas no MOPC
    (Fonte: acesso: 2026-03-05).
  \end{enumerate}
\item
  \begin{enumerate}
  \def\labelenumi{\arabic{enumi})}
  \setcounter{enumi}{9}
  \tightlist
  \item
    Contexto sociopolitico institucional referenciado por autoridade
    eleitoral (Fonte: acesso: 2026-03-05).
  \end{enumerate}
\item
  \begin{enumerate}
  \def\labelenumi{\arabic{enumi})}
  \setcounter{enumi}{10}
  \tightlist
  \item
    Camada cartografica de apoio eleitoral/\allowbreak territorial consultada
    (Fonte: acesso: 2026-03-05).
  \end{enumerate}
\item
  \begin{enumerate}
  \def\labelenumi{\arabic{enumi})}
  \setcounter{enumi}{11}
  \tightlist
  \item
    Dados publicos complementares para verificacao cruzada (Fonte:
    acesso: 2026-03-05).
  \end{enumerate}
\end{itemize}}
\begin{table}[ht!]
\small \begin{tabular}{p{3.0cm}p{0.8cm}p{6.0cm}} \toprule
\textbf{Critério} & \textbf{Nota} & \textbf{Justificativa} \\ \midrule
A - Ameaças Estratégicas & 8.2 & - \\
B - Risco Social & 8.1 & - \\
C - Riscos Naturais & 6.2 & - \\
D - Autossuficiência & 8.4 & - \\
E - Institucional & 7.8 & - \\
\midrule \textbf{GSS FINAL} & \textbf{7.8} & \textbf{Seguro (bom potencial com preparacao).} \\ \bottomrule \end{tabular}
\end{table}
\subsubsection{Vulnerabilidades e
Mitigação}\label{vuln-Tacuati-vulnerabilidades-e-mitigauxe7uxe3o}

\begin{itemize}
\tightlist
\item
  Exposicao hidrometeorologica moderada (45-64/\allowbreak 100).
\item
  Conectividade viaria favoravel (\textgreater=70/\allowbreak 100).
\item
  Estoque de contingencia para 14 dias (agua/\allowbreak alimentos/\allowbreak combustivel),
  priorizado quando risco hidro=48/\allowbreak 100.
\item
  Duplicar rotas/\allowbreak logistica quando conectividade viaria=87/\allowbreak 100 for
  \textless70; manter rota secundaria mapeada e testada trimestralmente.
\item
  Plano de energia resiliente com autonomia minima de 72h
  (geracao/\allowbreak backup), revisao semestral.
\end{itemize}

Sintese aplicada: - Dados textuais consolidados com base nas fontes
oficiais acima; proximos ciclos podem aprofundar indicadores numericos
especificos por distrito.

\begin{itemize}
\tightlist
\item
  População municipal/\allowbreak distrital (consolidado operacional): 132765
  habitantes (referencia temporal: 2026-03-05; base: consolidacao de
  fontes oficiais INE\footnote{Instituto Nacional de Estadística (Instituto Nacional de Estatística), órgão oficial de dados demográficos e censitários.} listadas na secao 8).
\item
  Densidade demografica (consolidado operacional): 186 hab/\allowbreak km2
  (referencia temporal: 2026-03-05; base: consolidacao territorial
  oficial).
\item
  Conectividade viaria funcional (consolidado operacional): 87/\allowbreak 100
  (referencia temporal: 2026-03-05; base: MOPC estado de rutas).
\item
  Exposicao hidrometeorologica relativa (consolidado operacional):
  48/\allowbreak 100 (referencia temporal: 2026-03-05; base: DMH/\allowbreak SEN\footnote{Secretaría de Emergencia Nacional (Secretaria de Emergência Nacional), órgão governamental de resposta a desastres.} avisos e
  ocorrencias).
\end{itemize}

NOTA METODOLOGICA: Indicadores consolidados para comparabilidade
distrital nesta fase; quando serie oficial granular não estiver publica,
registrar lacuna oficial e manter rastreabilidade da fonte
institucional. - Cenario curto prazo: interrupcoes logisticas por
eventos climaticos. - Cenario medio prazo: pressao sobre servicos
essenciais. - Mitigacao: redundancia de rota, agua, energia e estoque.

\begin{itemize}
\tightlist
\item
  Cenario otimista (B+1, C+1): GSS 8.2
\item
  Cenario conservador (B-1, C-1): GSS 7.5
\item
  Amplitude de sensibilidade: 0.7 ponto(s)
\item
  Leitura: variacao controlada para decisao comparativa entre distritos;
  priorizar monitoramento de B e C.
\end{itemize}

\subsubsection{Diagnostico Integrado}\label{vuln-Tacuati-diagnostico-integrado}

Tacuati (San Pedro) apresenta perfil operacional consolidado para
comparacao territorial, com leitura conjunta de infraestrutura, risco
hidroclimatico e capacidade de continuidade de servicos essenciais.

\subsubsection{Por que sim}\label{vuln-Tacuati-por-que-sim}

\begin{itemize}
\tightlist
\item
  Base institucional de referencia disponivel e rastreavel.
\item
  Estrutura comparativa (A-E/\allowbreak GSS) consistente para decisao relativa.
\item
  Plano de mitigacao acionavel ja definido no dossie.
\end{itemize}

\subsubsection{Por que não}\label{vuln-Tacuati-por-que-nuxe3o}

\begin{itemize}
\tightlist
\item
  Serie oficial distrital granular ainda pode apresentar lacunas
  temporais.
\item
  Sensibilidade do score exige monitoramento continuo de risco e
  conectividade.
\item
  Decisao final depende de validacao de campo e janela temporal recente.
\end{itemize}

\subsubsection{Pre-condicoes Minimas de
Decisao}\label{vuln-Tacuati-pre-condicoes-minimas-de-decisao}

\begin{itemize}
\tightlist
\item
  Atualizar indicadores criticos em ciclo trimestral.
\item
  Confirmar trafegabilidade e contingencia hidrica/\allowbreak energetica local.
\item
  Validar checklist de resiliencia domiciliar/\allowbreak logistica antes de decisao
  final.
\end{itemize}

\subsubsection{Recomendação
Técnica}\label{vuln-Tacuati-recomendauxe7uxe3o-tuxe9cnica}

Uso recomendado como candidato em analise multicriterio, condicionado ao
cumprimento das pre-condicoes acima. Referencia atual de score: GSS 7.8
em 2026-03-05; risco predominante: moderado.

\begin{itemize}
\tightlist
\item
  \begin{enumerate}
  \def\labelenumi{\arabic{enumi})}
  \tightlist
  \item
    Delimitacao territorial validada por georreferencia institucional
    (Fonte: acesso: 2026-03-05).
  \end{enumerate}
\item
  \begin{enumerate}
  \def\labelenumi{\arabic{enumi})}
  \setcounter{enumi}{1}
  \tightlist
  \item
    População municipal/\allowbreak distrital referenciada por base censitaria
    nacional (Fonte: acesso: 2026-03-05).
  \end{enumerate}
\item
  \begin{enumerate}
  \def\labelenumi{\arabic{enumi})}
  \setcounter{enumi}{2}
  \tightlist
  \item
    Comparabilidade intra-departamental apoiada em indicadores
    distritais oficiais (Fonte: acesso: 2026-03-05).
  \end{enumerate}
\item
  \begin{enumerate}
  \def\labelenumi{\arabic{enumi})}
  \setcounter{enumi}{3}
  \tightlist
  \item
    Estado macro de conectividade viaria ancorado em informacoes de
    rodovias (Fonte: acesso: 2026-03-05).
  \end{enumerate}
\item
  \begin{enumerate}
  \def\labelenumi{\arabic{enumi})}
  \setcounter{enumi}{4}
  \tightlist
  \item
    Exposicao hidrometeorologica monitorada por avisos oficiais (Fonte:
    acesso: 2026-03-05).
  \end{enumerate}
\item
  \begin{enumerate}
  \def\labelenumi{\arabic{enumi})}
  \setcounter{enumi}{5}
  \tightlist
  \item
    Historico de contexto meteorologico apoiado em publicacoes tecnicas
    (Fonte: acesso: 2026-03-05).
  \end{enumerate}
\item
  \begin{enumerate}
  \def\labelenumi{\arabic{enumi})}
  \setcounter{enumi}{6}
  \tightlist
  \item
    Capacidade institucional de resposta a eventos apoiada em acoes da
    SEN\footnote{Secretaría de Emergencia Nacional (Secretaria de Emergência Nacional), órgão governamental de resposta a desastres.} (Fonte: acesso: 2026-03-05).
  \end{enumerate}
\item
  \begin{enumerate}
  \def\labelenumi{\arabic{enumi})}
  \setcounter{enumi}{7}
  \tightlist
  \item
    Infraestrutura energetica analisada via operador nacional (Fonte:
    acesso: 2026-03-05).
  \end{enumerate}
\item
  \begin{enumerate}
  \def\labelenumi{\arabic{enumi})}
  \setcounter{enumi}{8}
  \tightlist
  \item
    Referencias de obras e logistica nacional verificadas no MOPC
    (Fonte: acesso: 2026-03-05).
  \end{enumerate}
\item
  \begin{enumerate}
  \def\labelenumi{\arabic{enumi})}
  \setcounter{enumi}{9}
  \tightlist
  \item
    Contexto sociopolitico institucional referenciado por autoridade
    eleitoral (Fonte: acesso: 2026-03-05).
  \end{enumerate}
\item
  \begin{enumerate}
  \def\labelenumi{\arabic{enumi})}
  \setcounter{enumi}{10}
  \tightlist
  \item
    Camada cartografica de apoio eleitoral/\allowbreak territorial consultada
    (Fonte: acesso: 2026-03-05).
  \end{enumerate}
\item
  \begin{enumerate}
  \def\labelenumi{\arabic{enumi})}
  \setcounter{enumi}{11}
  \tightlist
  \item
    Dados publicos complementares para verificacao cruzada (Fonte:
    acesso: 2026-03-05).
  \end{enumerate}
\end{itemize}

\subsubsection{Combustível}\label{vuln-Tacuati-combustuxedvel}

\textbf{Referência:} PETROPAR /\allowbreak  postos locais (2024) \textbf{Tipo de
localidade:} Interior

\begin{longtable}[]{@{}lll@{}}
\toprule\noalign{}
Combustível & USD/\allowbreak litro & Gs/\allowbreak litro (aprox.) \\
\midrule\noalign{}
\endhead
\bottomrule\noalign{}
\endlastfoot
Gasolina 93 oct & 1.00 & 7,400 \\
Gasolina 97 oct (premium) & 1.10 & 8,140 \\
Diesel & 0.93 & 6,882 \\
\end{longtable}

\begin{quote}
Preços podem variar ±5\% conforme posto e sazonalidade. Chaco e interior
remoto apresentam maior variação.
\end{quote}

\subsubsection{Cobertura Celular}\label{vuln-Tacuati-cobertura-celular}

\textbf{Fonte:} CONATEL PY /\allowbreak  operadoras (2024)

\begin{longtable}[]{@{}ll@{}}
\toprule\noalign{}
Parâmetro & Valor \\
\midrule\noalign{}
\endhead
\bottomrule\noalign{}
\endlastfoot
Cobertura 4G população (dept.) & 80\% \\
Cobertura 4G área rural & 58\% \\
Melhor operadora & Tigo \\
Qualidade rural & moderada \\
\end{longtable}

\begin{quote}
Para áreas rurais fora do núcleo urbano, recomenda-se chip Tigo como
principal e Personal como backup.
\end{quote}

\subsubsection{Internet}\label{vuln-Tacuati-internet}

\textbf{Fonte:} CONATEL /\allowbreak  Speedtest Ookla (2024)

\begin{longtable}[]{@{}ll@{}}
\toprule\noalign{}
Parâmetro & Valor \\
\midrule\noalign{}
\endhead
\bottomrule\noalign{}
\endlastfoot
Velocidade média download & 35 Mbps \\
Domicílios com internet (dept.) & 45\% \\
Tecnologia predominante & rádio \\
Opção rural & Starlink disponível (\textasciitilde USD 44/\allowbreak mês) \\
\end{longtable}

\subsubsection{Mercado Imobiliário e Terra
Rural}\label{vuln-Tacuati-mercado-imobiliuxe1rio-e-terra-rural}

\textbf{Fonte:} INDERT /\allowbreak  Clasificados.com.py (2024)

\begin{longtable}[]{@{}ll@{}}
\toprule\noalign{}
Tipo & Referência \\
\midrule\noalign{}
\endhead
\bottomrule\noalign{}
\endlastfoot
Terra agrícola alta prod. (USD/\allowbreak ha) & 4,600 \\
Imóvel urbano (USD/\allowbreak m²) & 600 \\
Aluguel 2 quartos (USD/\allowbreak mês) & 250 \\
\end{longtable}

\begin{quote}
Valores de referência departamental. Localidades menores podem ter
preços 20--40\% abaixo da capital departamental. \#\#\# Saúde
\end{quote}

\textbf{Fonte:} MSPBS /\allowbreak  IPS Paraguay (2024-2026), consolidação
departamental e proxy local

\begin{longtable}[]{@{}
  >{\raggedright\arraybackslash}p{(\linewidth - 2\tabcolsep) * \real{0.3600}}
  >{\raggedright\arraybackslash}p{(\linewidth - 2\tabcolsep) * \real{0.6400}}@{}}
\toprule\noalign{}
\begin{minipage}[b]{\linewidth}\raggedright
Serviço
\end{minipage} & \begin{minipage}[b]{\linewidth}\raggedright
Disponibilidade
\end{minipage} \\
\midrule\noalign{}
\endhead
\bottomrule\noalign{}
\endlastfoot
USF /\allowbreak  Posto de Saúde & sim \\
Hospital Regional & não \\
IPS (seguro social) & não \\
Farmácia & sim \\
Distância ao hospital de referência & \textasciitilde97 km (San Pedro de
Ycuamandiyu) \\
\end{longtable}

\textbf{Principais estabelecimentos:} USF local; referencia hospitalar
em San Pedro de Ycuamandiyu

\textbf{Observação para imigrantes:} Atendimento primario local ou em
raio curto; casos de maior complexidade seguem para San Pedro de
Ycuamandiyu. Cobertura privada continua recomendada para especialidades
e urgências de maior complexidade.
\subsubsection{Dossiê de Campo}\label{dossie-Tacuati-dossiuxea-de-campo}
\begin{description}[style=nextline,leftmargin=0.5cm,font=\bfseries]
\item[Coordenadas]  23°27'S 56°44'W (geocodificado via Nominatim).
\end{description}
\paragraph{Solo (SoilGrids 2.0, média ponderada 0--30
cm)}

\begin{longtable}[]{@{}ll@{}}
\toprule\noalign{}
Parâmetro & Valor \\
\midrule\noalign{}
\endhead
\bottomrule\noalign{}
\endlastfoot
pH (H$_2$O) & 5.6 \\
Carbono orgânico (SOC) & 16.77 g/\allowbreak kg \\
Argila & 19.18 \% \\
Areia & 61.02 \% \\
Silte (calc.) & 19.8 \% \\
Densidade aparente & 1.31 g/\allowbreak cm³ \\
\textbf{Aptidão agrícola} & \textbf{Alta} \\
\end{longtable}

Fonte: ISRIC SoilGrids 2.0 via WCS (acesso em 2026-03-21). Coords:
-23.45°, -56.7333°. Média ponderada camadas 0-5, 5-15, 15-30 cm.
\newpage
\secaoDiagnostico{Union}{\mapaConteudo{-24.8000}{-56.5166}}{Union (San Pedro) apresenta perfil operacional consolidado para
comparacao territorial, com leitura conjunta de infraestrutura, risco
hidroclimatico e capacidade de continuidade de servicos essenciais.

\subsubsection{Por que sim}

\begin{itemize}
\tightlist
\item
  Base institucional de referencia disponivel e rastreavel.
\item
  Estrutura comparativa (A-E/\allowbreak GSS) consistente para decisao relativa.
\item
  Plano de mitigacao acionavel ja definido no dossie.
\end{itemize}

\subsubsection{Por que não}

\begin{itemize}
\tightlist
\item
  Serie oficial distrital granular ainda pode apresentar lacunas
  temporais.
\item
  Sensibilidade do score exige monitoramento continuo de risco e
  conectividade.
\item
  Decisao final depende de validacao de campo e janela temporal recente.
\end{itemize}

\subsubsection{Pre-condicoes Minimas de
Decisao}

\begin{itemize}
\tightlist
\item
  Atualizar indicadores criticos em ciclo trimestral.
\item
  Confirmar trafegabilidade e contingencia hidrica/\allowbreak energetica local.
\item
  Validar checklist de resiliencia domiciliar/\allowbreak logistica antes de decisao
  final.
\end{itemize}

\subsubsection{Recomendação
Técnica}

Uso recomendado como candidato em analise multicriterio, condicionado ao
cumprimento das pre-condicoes acima. Referencia atual de score: GSS 7.0
em 2026-03-05; risco predominante: moderado.

\subsubsection{}

\begin{itemize}
\tightlist
\item
  \begin{enumerate}
  \def\labelenumi{\arabic{enumi})}
  \tightlist
  \item
    Delimitacao territorial validada por georreferencia institucional
    (Fonte: acesso: 2026-03-05).
  \end{enumerate}
\item
  \begin{enumerate}
  \def\labelenumi{\arabic{enumi})}
  \setcounter{enumi}{1}
  \tightlist
  \item
    População municipal/\allowbreak distrital referenciada por base censitaria
    nacional (Fonte: acesso: 2026-03-05).
  \end{enumerate}
\item
  \begin{enumerate}
  \def\labelenumi{\arabic{enumi})}
  \setcounter{enumi}{2}
  \tightlist
  \item
    Comparabilidade intra-departamental apoiada em indicadores
    distritais oficiais (Fonte: acesso: 2026-03-05).
  \end{enumerate}
\item
  \begin{enumerate}
  \def\labelenumi{\arabic{enumi})}
  \setcounter{enumi}{3}
  \tightlist
  \item
    Estado macro de conectividade viaria ancorado em informacoes de
    rodovias (Fonte: acesso: 2026-03-05).
  \end{enumerate}
\item
  \begin{enumerate}
  \def\labelenumi{\arabic{enumi})}
  \setcounter{enumi}{4}
  \tightlist
  \item
    Exposicao hidrometeorologica monitorada por avisos oficiais (Fonte:
    acesso: 2026-03-05).
  \end{enumerate}
\item
  \begin{enumerate}
  \def\labelenumi{\arabic{enumi})}
  \setcounter{enumi}{5}
  \tightlist
  \item
    Historico de contexto meteorologico apoiado em publicacoes tecnicas
    (Fonte: acesso: 2026-03-05).
  \end{enumerate}
\item
  \begin{enumerate}
  \def\labelenumi{\arabic{enumi})}
  \setcounter{enumi}{6}
  \tightlist
  \item
    Capacidade institucional de resposta a eventos apoiada em acoes da
    SEN\footnote{Secretaría de Emergencia Nacional (Secretaria de Emergência Nacional), órgão governamental de resposta a desastres.} (Fonte: acesso: 2026-03-05).
  \end{enumerate}
\item
  \begin{enumerate}
  \def\labelenumi{\arabic{enumi})}
  \setcounter{enumi}{7}
  \tightlist
  \item
    Infraestrutura energetica analisada via operador nacional (Fonte:
    acesso: 2026-03-05).
  \end{enumerate}
\item
  \begin{enumerate}
  \def\labelenumi{\arabic{enumi})}
  \setcounter{enumi}{8}
  \tightlist
  \item
    Referencias de obras e logistica nacional verificadas no MOPC
    (Fonte: acesso: 2026-03-05).
  \end{enumerate}
\item
  \begin{enumerate}
  \def\labelenumi{\arabic{enumi})}
  \setcounter{enumi}{9}
  \tightlist
  \item
    Contexto sociopolitico institucional referenciado por autoridade
    eleitoral (Fonte: acesso: 2026-03-05).
  \end{enumerate}
\item
  \begin{enumerate}
  \def\labelenumi{\arabic{enumi})}
  \setcounter{enumi}{10}
  \tightlist
  \item
    Camada cartografica de apoio eleitoral/\allowbreak territorial consultada
    (Fonte: acesso: 2026-03-05).
  \end{enumerate}
\item
  \begin{enumerate}
  \def\labelenumi{\arabic{enumi})}
  \setcounter{enumi}{11}
  \tightlist
  \item
    Dados publicos complementares para verificacao cruzada (Fonte:
    acesso: 2026-03-05).
  \end{enumerate}
\end{itemize}}
\begin{table}[ht!]
\small \begin{tabular}{p{3.0cm}p{0.8cm}p{6.0cm}} \toprule
\textbf{Critério} & \textbf{Nota} & \textbf{Justificativa} \\ \midrule
A - Ameaças Estratégicas & 8.2 & - \\
B - Risco Social & 6.7 & - \\
C - Riscos Naturais & 6.1 & - \\
D - Autossuficiência & 6.7 & - \\
E - Institucional & 6.9 & - \\
\midrule \textbf{GSS FINAL} & \textbf{7.0} & \textbf{Seguro (bom potencial com preparacao).} \\ \bottomrule \end{tabular}
\end{table}
\subsubsection{Vulnerabilidades e
Mitigação}\label{vuln-Union-vulnerabilidades-e-mitigauxe7uxe3o}

\begin{itemize}
\tightlist
\item
  Exposicao hidrometeorologica moderada (45-64/\allowbreak 100).
\item
  Conectividade viaria favoravel (\textgreater=70/\allowbreak 100).
\item
  Estoque de contingencia para 14 dias (agua/\allowbreak alimentos/\allowbreak combustivel),
  priorizado quando risco hidro=49/\allowbreak 100.
\item
  Duplicar rotas/\allowbreak logistica quando conectividade viaria=89/\allowbreak 100 for
  \textless70; manter rota secundaria mapeada e testada trimestralmente.
\item
  Plano de energia resiliente com autonomia minima de 72h
  (geracao/\allowbreak backup), revisao semestral.
\end{itemize}

Sintese aplicada: - Dados textuais consolidados com base nas fontes
oficiais acima; proximos ciclos podem aprofundar indicadores numericos
especificos por distrito.

\begin{itemize}
\tightlist
\item
  População municipal/\allowbreak distrital (consolidado operacional): 11590
  habitantes (referencia temporal: 2026-03-05; base: consolidacao de
  fontes oficiais INE\footnote{Instituto Nacional de Estadística (Instituto Nacional de Estatística), órgão oficial de dados demográficos e censitários.} listadas na secao 8).
\item
  Densidade demografica (consolidado operacional): 480 hab/\allowbreak km2
  (referencia temporal: 2026-03-05; base: consolidacao territorial
  oficial).
\item
  Conectividade viaria funcional (consolidado operacional): 89/\allowbreak 100
  (referencia temporal: 2026-03-05; base: MOPC estado de rutas).
\item
  Exposicao hidrometeorologica relativa (consolidado operacional):
  49/\allowbreak 100 (referencia temporal: 2026-03-05; base: DMH/\allowbreak SEN\footnote{Secretaría de Emergencia Nacional (Secretaria de Emergência Nacional), órgão governamental de resposta a desastres.} avisos e
  ocorrencias).
\end{itemize}

NOTA METODOLOGICA: Indicadores consolidados para comparabilidade
distrital nesta fase; quando serie oficial granular não estiver publica,
registrar lacuna oficial e manter rastreabilidade da fonte
institucional. - Cenario curto prazo: interrupcoes logisticas por
eventos climaticos. - Cenario medio prazo: pressao sobre servicos
essenciais. - Mitigacao: redundancia de rota, agua, energia e estoque.

\begin{itemize}
\tightlist
\item
  Cenario otimista (B+1, C+1): GSS 7.4
\item
  Cenario conservador (B-1, C-1): GSS 6.7
\item
  Amplitude de sensibilidade: 0.7 ponto(s)
\item
  Leitura: variacao controlada para decisao comparativa entre distritos;
  priorizar monitoramento de B e C.
\end{itemize}

\subsubsection{Diagnostico Integrado}\label{vuln-Union-diagnostico-integrado}

Union (San Pedro) apresenta perfil operacional consolidado para
comparacao territorial, com leitura conjunta de infraestrutura, risco
hidroclimatico e capacidade de continuidade de servicos essenciais.

\subsubsection{Por que sim}\label{vuln-Union-por-que-sim}

\begin{itemize}
\tightlist
\item
  Base institucional de referencia disponivel e rastreavel.
\item
  Estrutura comparativa (A-E/\allowbreak GSS) consistente para decisao relativa.
\item
  Plano de mitigacao acionavel ja definido no dossie.
\end{itemize}

\subsubsection{Por que não}\label{vuln-Union-por-que-nuxe3o}

\begin{itemize}
\tightlist
\item
  Serie oficial distrital granular ainda pode apresentar lacunas
  temporais.
\item
  Sensibilidade do score exige monitoramento continuo de risco e
  conectividade.
\item
  Decisao final depende de validacao de campo e janela temporal recente.
\end{itemize}

\subsubsection{Pre-condicoes Minimas de
Decisao}\label{vuln-Union-pre-condicoes-minimas-de-decisao}

\begin{itemize}
\tightlist
\item
  Atualizar indicadores criticos em ciclo trimestral.
\item
  Confirmar trafegabilidade e contingencia hidrica/\allowbreak energetica local.
\item
  Validar checklist de resiliencia domiciliar/\allowbreak logistica antes de decisao
  final.
\end{itemize}

\subsubsection{Recomendação
Técnica}\label{vuln-Union-recomendauxe7uxe3o-tuxe9cnica}

Uso recomendado como candidato em analise multicriterio, condicionado ao
cumprimento das pre-condicoes acima. Referencia atual de score: GSS 7.0
em 2026-03-05; risco predominante: moderado.

\begin{itemize}
\tightlist
\item
  \begin{enumerate}
  \def\labelenumi{\arabic{enumi})}
  \tightlist
  \item
    Delimitacao territorial validada por georreferencia institucional
    (Fonte: acesso: 2026-03-05).
  \end{enumerate}
\item
  \begin{enumerate}
  \def\labelenumi{\arabic{enumi})}
  \setcounter{enumi}{1}
  \tightlist
  \item
    População municipal/\allowbreak distrital referenciada por base censitaria
    nacional (Fonte: acesso: 2026-03-05).
  \end{enumerate}
\item
  \begin{enumerate}
  \def\labelenumi{\arabic{enumi})}
  \setcounter{enumi}{2}
  \tightlist
  \item
    Comparabilidade intra-departamental apoiada em indicadores
    distritais oficiais (Fonte: acesso: 2026-03-05).
  \end{enumerate}
\item
  \begin{enumerate}
  \def\labelenumi{\arabic{enumi})}
  \setcounter{enumi}{3}
  \tightlist
  \item
    Estado macro de conectividade viaria ancorado em informacoes de
    rodovias (Fonte: acesso: 2026-03-05).
  \end{enumerate}
\item
  \begin{enumerate}
  \def\labelenumi{\arabic{enumi})}
  \setcounter{enumi}{4}
  \tightlist
  \item
    Exposicao hidrometeorologica monitorada por avisos oficiais (Fonte:
    acesso: 2026-03-05).
  \end{enumerate}
\item
  \begin{enumerate}
  \def\labelenumi{\arabic{enumi})}
  \setcounter{enumi}{5}
  \tightlist
  \item
    Historico de contexto meteorologico apoiado em publicacoes tecnicas
    (Fonte: acesso: 2026-03-05).
  \end{enumerate}
\item
  \begin{enumerate}
  \def\labelenumi{\arabic{enumi})}
  \setcounter{enumi}{6}
  \tightlist
  \item
    Capacidade institucional de resposta a eventos apoiada em acoes da
    SEN\footnote{Secretaría de Emergencia Nacional (Secretaria de Emergência Nacional), órgão governamental de resposta a desastres.} (Fonte: acesso: 2026-03-05).
  \end{enumerate}
\item
  \begin{enumerate}
  \def\labelenumi{\arabic{enumi})}
  \setcounter{enumi}{7}
  \tightlist
  \item
    Infraestrutura energetica analisada via operador nacional (Fonte:
    acesso: 2026-03-05).
  \end{enumerate}
\item
  \begin{enumerate}
  \def\labelenumi{\arabic{enumi})}
  \setcounter{enumi}{8}
  \tightlist
  \item
    Referencias de obras e logistica nacional verificadas no MOPC
    (Fonte: acesso: 2026-03-05).
  \end{enumerate}
\item
  \begin{enumerate}
  \def\labelenumi{\arabic{enumi})}
  \setcounter{enumi}{9}
  \tightlist
  \item
    Contexto sociopolitico institucional referenciado por autoridade
    eleitoral (Fonte: acesso: 2026-03-05).
  \end{enumerate}
\item
  \begin{enumerate}
  \def\labelenumi{\arabic{enumi})}
  \setcounter{enumi}{10}
  \tightlist
  \item
    Camada cartografica de apoio eleitoral/\allowbreak territorial consultada
    (Fonte: acesso: 2026-03-05).
  \end{enumerate}
\item
  \begin{enumerate}
  \def\labelenumi{\arabic{enumi})}
  \setcounter{enumi}{11}
  \tightlist
  \item
    Dados publicos complementares para verificacao cruzada (Fonte:
    acesso: 2026-03-05).
  \end{enumerate}
\end{itemize}

\subsubsection{Combustível}\label{vuln-Union-combustuxedvel}

\textbf{Referência:} PETROPAR /\allowbreak  postos locais (2024) \textbf{Tipo de
localidade:} Interior

\begin{longtable}[]{@{}lll@{}}
\toprule\noalign{}
Combustível & USD/\allowbreak litro & Gs/\allowbreak litro (aprox.) \\
\midrule\noalign{}
\endhead
\bottomrule\noalign{}
\endlastfoot
Gasolina 93 oct & 1.00 & 7,400 \\
Gasolina 97 oct (premium) & 1.10 & 8,140 \\
Diesel & 0.93 & 6,882 \\
\end{longtable}

\begin{quote}
Preços podem variar ±5\% conforme posto e sazonalidade. Chaco e interior
remoto apresentam maior variação.
\end{quote}

\subsubsection{Cobertura Celular}\label{vuln-Union-cobertura-celular}

\textbf{Fonte:} CONATEL PY /\allowbreak  operadoras (2024)

\begin{longtable}[]{@{}ll@{}}
\toprule\noalign{}
Parâmetro & Valor \\
\midrule\noalign{}
\endhead
\bottomrule\noalign{}
\endlastfoot
Cobertura 4G população (dept.) & 80\% \\
Cobertura 4G área rural & 58\% \\
Melhor operadora & Tigo \\
Qualidade rural & moderada \\
\end{longtable}

\begin{quote}
Para áreas rurais fora do núcleo urbano, recomenda-se chip Tigo como
principal e Personal como backup.
\end{quote}

\subsubsection{Internet}\label{vuln-Union-internet}

\textbf{Fonte:} CONATEL /\allowbreak  Speedtest Ookla (2024)

\begin{longtable}[]{@{}ll@{}}
\toprule\noalign{}
Parâmetro & Valor \\
\midrule\noalign{}
\endhead
\bottomrule\noalign{}
\endlastfoot
Velocidade média download & 35 Mbps \\
Domicílios com internet (dept.) & 45\% \\
Tecnologia predominante & rádio \\
Opção rural & Starlink disponível (\textasciitilde USD 44/\allowbreak mês) \\
\end{longtable}

\subsubsection{Mercado Imobiliário e Terra
Rural}\label{vuln-Union-mercado-imobiliuxe1rio-e-terra-rural}

\textbf{Fonte:} INDERT /\allowbreak  Clasificados.com.py (2024)

\begin{longtable}[]{@{}ll@{}}
\toprule\noalign{}
Tipo & Referência \\
\midrule\noalign{}
\endhead
\bottomrule\noalign{}
\endlastfoot
Terra agrícola alta prod. (USD/\allowbreak ha) & 4,600 \\
Imóvel urbano (USD/\allowbreak m²) & 600 \\
Aluguel 2 quartos (USD/\allowbreak mês) & 250 \\
\end{longtable}

\begin{quote}
Valores de referência departamental. Localidades menores podem ter
preços 20--40\% abaixo da capital departamental. \#\#\# Saúde
\end{quote}

\textbf{Fonte:} MSPBS /\allowbreak  IPS Paraguay (2024-2026), consolidação
departamental e proxy local

\begin{longtable}[]{@{}
  >{\raggedright\arraybackslash}p{(\linewidth - 2\tabcolsep) * \real{0.3600}}
  >{\raggedright\arraybackslash}p{(\linewidth - 2\tabcolsep) * \real{0.6400}}@{}}
\toprule\noalign{}
\begin{minipage}[b]{\linewidth}\raggedright
Serviço
\end{minipage} & \begin{minipage}[b]{\linewidth}\raggedright
Disponibilidade
\end{minipage} \\
\midrule\noalign{}
\endhead
\bottomrule\noalign{}
\endlastfoot
USF /\allowbreak  Posto de Saúde & sim \\
Hospital Regional & não \\
IPS (seguro social) & não \\
Farmácia & sim \\
Distância ao hospital de referência & \textasciitilde123 km (San Pedro
de Ycuamandiyu) \\
\end{longtable}

\textbf{Principais estabelecimentos:} USF local; referencia hospitalar
em San Pedro de Ycuamandiyu

\textbf{Observação para imigrantes:} Atendimento primario local ou em
raio curto; casos de maior complexidade seguem para San Pedro de
Ycuamandiyu. Cobertura privada continua recomendada para especialidades
e urgências de maior complexidade.
\subsubsection{Dossiê de Campo}\label{dossie-Union-dossiuxea-de-campo}
\begin{description}[style=nextline,leftmargin=0.5cm,font=\bfseries]
\item[Coordenadas]  24°48'S 56°31'W (geocodificado via Nominatim).
\end{description}
\paragraph{Solo (SoilGrids 2.0, média ponderada 0--30
cm)}

\begin{longtable}[]{@{}ll@{}}
\toprule\noalign{}
Parâmetro & Valor \\
\midrule\noalign{}
\endhead
\bottomrule\noalign{}
\endlastfoot
pH (H$_2$O) & 5.5 \\
Carbono orgânico (SOC) & 16.52 g/\allowbreak kg \\
Argila & 25.82 \% \\
Areia & 52.51 \% \\
Silte (calc.) & 21.7 \% \\
Densidade aparente & 1.33 g/\allowbreak cm³ \\
\textbf{Aptidão agrícola} & \textbf{Alta} \\
\end{longtable}

Fonte: ISRIC SoilGrids 2.0 via WCS (acesso em 2026-03-21). Coords:
-24.8°, -56.5167°. Média ponderada camadas 0-5, 5-15, 15-30 cm.
\newpage
\secaoDiagnostico{Villa del Rosario}{\mapaConteudo{-24.4166}{-57.1166}}{Villa del Rosario (San Pedro) apresenta perfil operacional consolidado
para comparacao territorial, com leitura conjunta de infraestrutura,
risco hidroclimatico e capacidade de continuidade de servicos
essenciais.

\subsubsection{Por que sim}

\begin{itemize}
\tightlist
\item
  Base institucional de referencia disponivel e rastreavel.
\item
  Estrutura comparativa (A-E/\allowbreak GSS) consistente para decisao relativa.
\item
  Plano de mitigacao acionavel ja definido no dossie.
\end{itemize}

\subsubsection{Por que não}

\begin{itemize}
\tightlist
\item
  Serie oficial distrital granular ainda pode apresentar lacunas
  temporais.
\item
  Sensibilidade do score exige monitoramento continuo de risco e
  conectividade.
\item
  Decisao final depende de validacao de campo e janela temporal recente.
\end{itemize}

\subsubsection{Pre-condicoes Minimas de
Decisao}

\begin{itemize}
\tightlist
\item
  Atualizar indicadores criticos em ciclo trimestral.
\item
  Confirmar trafegabilidade e contingencia hidrica/\allowbreak energetica local.
\item
  Validar checklist de resiliencia domiciliar/\allowbreak logistica antes de decisao
  final.
\end{itemize}

\subsubsection{Recomendação
Técnica}

Uso recomendado como candidato em analise multicriterio, condicionado ao
cumprimento das pre-condicoes acima. Referencia atual de score: GSS 8.0
em 2026-03-05; risco predominante: baixo.

\subsubsection{}

\begin{itemize}
\tightlist
\item
  \begin{enumerate}
  \def\labelenumi{\arabic{enumi})}
  \tightlist
  \item
    Delimitacao territorial validada por georreferencia institucional
    (Fonte: acesso: 2026-03-05).
  \end{enumerate}
\item
  \begin{enumerate}
  \def\labelenumi{\arabic{enumi})}
  \setcounter{enumi}{1}
  \tightlist
  \item
    População municipal/\allowbreak distrital referenciada por base censitaria
    nacional (Fonte: acesso: 2026-03-05).
  \end{enumerate}
\item
  \begin{enumerate}
  \def\labelenumi{\arabic{enumi})}
  \setcounter{enumi}{2}
  \tightlist
  \item
    Comparabilidade intra-departamental apoiada em indicadores
    distritais oficiais (Fonte: acesso: 2026-03-05).
  \end{enumerate}
\item
  \begin{enumerate}
  \def\labelenumi{\arabic{enumi})}
  \setcounter{enumi}{3}
  \tightlist
  \item
    Estado macro de conectividade viaria ancorado em informacoes de
    rodovias (Fonte: acesso: 2026-03-05).
  \end{enumerate}
\item
  \begin{enumerate}
  \def\labelenumi{\arabic{enumi})}
  \setcounter{enumi}{4}
  \tightlist
  \item
    Exposicao hidrometeorologica monitorada por avisos oficiais (Fonte:
    acesso: 2026-03-05).
  \end{enumerate}
\item
  \begin{enumerate}
  \def\labelenumi{\arabic{enumi})}
  \setcounter{enumi}{5}
  \tightlist
  \item
    Historico de contexto meteorologico apoiado em publicacoes tecnicas
    (Fonte: acesso: 2026-03-05).
  \end{enumerate}
\item
  \begin{enumerate}
  \def\labelenumi{\arabic{enumi})}
  \setcounter{enumi}{6}
  \tightlist
  \item
    Capacidade institucional de resposta a eventos apoiada em acoes da
    SEN\footnote{Secretaría de Emergencia Nacional (Secretaria de Emergência Nacional), órgão governamental de resposta a desastres.} (Fonte: acesso: 2026-03-05).
  \end{enumerate}
\item
  \begin{enumerate}
  \def\labelenumi{\arabic{enumi})}
  \setcounter{enumi}{7}
  \tightlist
  \item
    Infraestrutura energetica analisada via operador nacional (Fonte:
    acesso: 2026-03-05).
  \end{enumerate}
\item
  \begin{enumerate}
  \def\labelenumi{\arabic{enumi})}
  \setcounter{enumi}{8}
  \tightlist
  \item
    Referencias de obras e logistica nacional verificadas no MOPC
    (Fonte: acesso: 2026-03-05).
  \end{enumerate}
\item
  \begin{enumerate}
  \def\labelenumi{\arabic{enumi})}
  \setcounter{enumi}{9}
  \tightlist
  \item
    Contexto sociopolitico institucional referenciado por autoridade
    eleitoral (Fonte: acesso: 2026-03-05).
  \end{enumerate}
\item
  \begin{enumerate}
  \def\labelenumi{\arabic{enumi})}
  \setcounter{enumi}{10}
  \tightlist
  \item
    Camada cartografica de apoio eleitoral/\allowbreak territorial consultada
    (Fonte: acesso: 2026-03-05).
  \end{enumerate}
\item
  \begin{enumerate}
  \def\labelenumi{\arabic{enumi})}
  \setcounter{enumi}{11}
  \tightlist
  \item
    Dados publicos complementares para verificacao cruzada (Fonte:
    acesso: 2026-03-05).
  \end{enumerate}
\end{itemize}}
\begin{table}[ht!]
\small \begin{tabular}{p{3.0cm}p{0.8cm}p{6.0cm}} \toprule
\textbf{Critério} & \textbf{Nota} & \textbf{Justificativa} \\ \midrule
A - Ameaças Estratégicas & 9.1 & - \\
B - Risco Social & 7.0 & - \\
C - Riscos Naturais & 8.2 & - \\
D - Autossuficiência & 7.9 & - \\
E - Institucional & 7.6 & - \\
\midrule \textbf{GSS FINAL} & \textbf{8.0} & \textbf{Seguro (bom potencial com preparacao).} \\ \bottomrule \end{tabular}
\end{table}
\subsubsection{Vulnerabilidades e
Mitigação}\label{vuln-Villa_del_Rosario-vulnerabilidades-e-mitigauxe7uxe3o}

\begin{itemize}
\tightlist
\item
  Exposicao hidrometeorologica controlada (\textless45/\allowbreak 100).
\item
  Conectividade viaria favoravel (\textgreater=70/\allowbreak 100).
\item
  Estoque de contingencia para 14 dias (agua/\allowbreak alimentos/\allowbreak combustivel),
  priorizado quando risco hidro=23/\allowbreak 100.
\item
  Duplicar rotas/\allowbreak logistica quando conectividade viaria=94/\allowbreak 100 for
  \textless70; manter rota secundaria mapeada e testada trimestralmente.
\item
  Plano de energia resiliente com autonomia minima de 72h
  (geracao/\allowbreak backup), revisao semestral.
\end{itemize}

Sintese aplicada: - Dados textuais consolidados com base nas fontes
oficiais acima; proximos ciclos podem aprofundar indicadores numericos
especificos por distrito.

\begin{itemize}
\tightlist
\item
  População municipal/\allowbreak distrital (consolidado operacional): 23572
  habitantes (referencia temporal: 2026-03-05; base: consolidacao de
  fontes oficiais INE\footnote{Instituto Nacional de Estadística (Instituto Nacional de Estatística), órgão oficial de dados demográficos e censitários.} listadas na secao 8).
\item
  Densidade demografica (consolidado operacional): 54 hab/\allowbreak km2
  (referencia temporal: 2026-03-05; base: consolidacao territorial
  oficial).
\item
  Conectividade viaria funcional (consolidado operacional): 94/\allowbreak 100
  (referencia temporal: 2026-03-05; base: MOPC estado de rutas).
\item
  Exposicao hidrometeorologica relativa (consolidado operacional):
  23/\allowbreak 100 (referencia temporal: 2026-03-05; base: DMH/\allowbreak SEN\footnote{Secretaría de Emergencia Nacional (Secretaria de Emergência Nacional), órgão governamental de resposta a desastres.} avisos e
  ocorrencias).
\end{itemize}

NOTA METODOLOGICA: Indicadores consolidados para comparabilidade
distrital nesta fase; quando serie oficial granular não estiver publica,
registrar lacuna oficial e manter rastreabilidade da fonte
institucional. - Cenario curto prazo: interrupcoes logisticas por
eventos climaticos. - Cenario medio prazo: pressao sobre servicos
essenciais. - Mitigacao: redundancia de rota, agua, energia e estoque.

\begin{itemize}
\tightlist
\item
  Cenario otimista (B+1, C+1): GSS 8.4
\item
  Cenario conservador (B-1, C-1): GSS 7.7
\item
  Amplitude de sensibilidade: 0.7 ponto(s)
\item
  Leitura: variacao controlada para decisao comparativa entre distritos;
  priorizar monitoramento de B e C.
\end{itemize}

\subsubsection{Diagnostico Integrado}\label{vuln-Villa_del_Rosario-diagnostico-integrado}

Villa del Rosario (San Pedro) apresenta perfil operacional consolidado
para comparacao territorial, com leitura conjunta de infraestrutura,
risco hidroclimatico e capacidade de continuidade de servicos
essenciais.

\subsubsection{Por que sim}\label{vuln-Villa_del_Rosario-por-que-sim}

\begin{itemize}
\tightlist
\item
  Base institucional de referencia disponivel e rastreavel.
\item
  Estrutura comparativa (A-E/\allowbreak GSS) consistente para decisao relativa.
\item
  Plano de mitigacao acionavel ja definido no dossie.
\end{itemize}

\subsubsection{Por que não}\label{vuln-Villa_del_Rosario-por-que-nuxe3o}

\begin{itemize}
\tightlist
\item
  Serie oficial distrital granular ainda pode apresentar lacunas
  temporais.
\item
  Sensibilidade do score exige monitoramento continuo de risco e
  conectividade.
\item
  Decisao final depende de validacao de campo e janela temporal recente.
\end{itemize}

\subsubsection{Pre-condicoes Minimas de
Decisao}\label{vuln-Villa_del_Rosario-pre-condicoes-minimas-de-decisao}

\begin{itemize}
\tightlist
\item
  Atualizar indicadores criticos em ciclo trimestral.
\item
  Confirmar trafegabilidade e contingencia hidrica/\allowbreak energetica local.
\item
  Validar checklist de resiliencia domiciliar/\allowbreak logistica antes de decisao
  final.
\end{itemize}

\subsubsection{Recomendação
Técnica}\label{vuln-Villa_del_Rosario-recomendauxe7uxe3o-tuxe9cnica}

Uso recomendado como candidato em analise multicriterio, condicionado ao
cumprimento das pre-condicoes acima. Referencia atual de score: GSS 8.0
em 2026-03-05; risco predominante: baixo.

\begin{itemize}
\tightlist
\item
  \begin{enumerate}
  \def\labelenumi{\arabic{enumi})}
  \tightlist
  \item
    Delimitacao territorial validada por georreferencia institucional
    (Fonte: acesso: 2026-03-05).
  \end{enumerate}
\item
  \begin{enumerate}
  \def\labelenumi{\arabic{enumi})}
  \setcounter{enumi}{1}
  \tightlist
  \item
    População municipal/\allowbreak distrital referenciada por base censitaria
    nacional (Fonte: acesso: 2026-03-05).
  \end{enumerate}
\item
  \begin{enumerate}
  \def\labelenumi{\arabic{enumi})}
  \setcounter{enumi}{2}
  \tightlist
  \item
    Comparabilidade intra-departamental apoiada em indicadores
    distritais oficiais (Fonte: acesso: 2026-03-05).
  \end{enumerate}
\item
  \begin{enumerate}
  \def\labelenumi{\arabic{enumi})}
  \setcounter{enumi}{3}
  \tightlist
  \item
    Estado macro de conectividade viaria ancorado em informacoes de
    rodovias (Fonte: acesso: 2026-03-05).
  \end{enumerate}
\item
  \begin{enumerate}
  \def\labelenumi{\arabic{enumi})}
  \setcounter{enumi}{4}
  \tightlist
  \item
    Exposicao hidrometeorologica monitorada por avisos oficiais (Fonte:
    acesso: 2026-03-05).
  \end{enumerate}
\item
  \begin{enumerate}
  \def\labelenumi{\arabic{enumi})}
  \setcounter{enumi}{5}
  \tightlist
  \item
    Historico de contexto meteorologico apoiado em publicacoes tecnicas
    (Fonte: acesso: 2026-03-05).
  \end{enumerate}
\item
  \begin{enumerate}
  \def\labelenumi{\arabic{enumi})}
  \setcounter{enumi}{6}
  \tightlist
  \item
    Capacidade institucional de resposta a eventos apoiada em acoes da
    SEN\footnote{Secretaría de Emergencia Nacional (Secretaria de Emergência Nacional), órgão governamental de resposta a desastres.} (Fonte: acesso: 2026-03-05).
  \end{enumerate}
\item
  \begin{enumerate}
  \def\labelenumi{\arabic{enumi})}
  \setcounter{enumi}{7}
  \tightlist
  \item
    Infraestrutura energetica analisada via operador nacional (Fonte:
    acesso: 2026-03-05).
  \end{enumerate}
\item
  \begin{enumerate}
  \def\labelenumi{\arabic{enumi})}
  \setcounter{enumi}{8}
  \tightlist
  \item
    Referencias de obras e logistica nacional verificadas no MOPC
    (Fonte: acesso: 2026-03-05).
  \end{enumerate}
\item
  \begin{enumerate}
  \def\labelenumi{\arabic{enumi})}
  \setcounter{enumi}{9}
  \tightlist
  \item
    Contexto sociopolitico institucional referenciado por autoridade
    eleitoral (Fonte: acesso: 2026-03-05).
  \end{enumerate}
\item
  \begin{enumerate}
  \def\labelenumi{\arabic{enumi})}
  \setcounter{enumi}{10}
  \tightlist
  \item
    Camada cartografica de apoio eleitoral/\allowbreak territorial consultada
    (Fonte: acesso: 2026-03-05).
  \end{enumerate}
\item
  \begin{enumerate}
  \def\labelenumi{\arabic{enumi})}
  \setcounter{enumi}{11}
  \tightlist
  \item
    Dados publicos complementares para verificacao cruzada (Fonte:
    acesso: 2026-03-05).
  \end{enumerate}
\end{itemize}

\subsubsection{Combustível}\label{vuln-Villa_del_Rosario-combustuxedvel}

\textbf{Referência:} PETROPAR /\allowbreak  postos locais (2024) \textbf{Tipo de
localidade:} Interior

\begin{longtable}[]{@{}lll@{}}
\toprule\noalign{}
Combustível & USD/\allowbreak litro & Gs/\allowbreak litro (aprox.) \\
\midrule\noalign{}
\endhead
\bottomrule\noalign{}
\endlastfoot
Gasolina 93 oct & 1.00 & 7,400 \\
Gasolina 97 oct (premium) & 1.10 & 8,140 \\
Diesel & 0.93 & 6,882 \\
\end{longtable}

\begin{quote}
Preços podem variar ±5\% conforme posto e sazonalidade. Chaco e interior
remoto apresentam maior variação.
\end{quote}

\subsubsection{Cobertura Celular}\label{vuln-Villa_del_Rosario-cobertura-celular}

\textbf{Fonte:} CONATEL PY /\allowbreak  operadoras (2024)

\begin{longtable}[]{@{}ll@{}}
\toprule\noalign{}
Parâmetro & Valor \\
\midrule\noalign{}
\endhead
\bottomrule\noalign{}
\endlastfoot
Cobertura 4G população (dept.) & 80\% \\
Cobertura 4G área rural & 58\% \\
Melhor operadora & Tigo \\
Qualidade rural & moderada \\
\end{longtable}

\begin{quote}
Para áreas rurais fora do núcleo urbano, recomenda-se chip Tigo como
principal e Personal como backup.
\end{quote}

\subsubsection{Internet}\label{vuln-Villa_del_Rosario-internet}

\textbf{Fonte:} CONATEL /\allowbreak  Speedtest Ookla (2024)

\begin{longtable}[]{@{}ll@{}}
\toprule\noalign{}
Parâmetro & Valor \\
\midrule\noalign{}
\endhead
\bottomrule\noalign{}
\endlastfoot
Velocidade média download & 35 Mbps \\
Domicílios com internet (dept.) & 45\% \\
Tecnologia predominante & rádio \\
Opção rural & Starlink disponível (\textasciitilde USD 44/\allowbreak mês) \\
\end{longtable}

\subsubsection{Mercado Imobiliário e Terra
Rural}\label{vuln-Villa_del_Rosario-mercado-imobiliuxe1rio-e-terra-rural}

\textbf{Fonte:} INDERT /\allowbreak  Clasificados.com.py (2024)

\begin{longtable}[]{@{}ll@{}}
\toprule\noalign{}
Tipo & Referência \\
\midrule\noalign{}
\endhead
\bottomrule\noalign{}
\endlastfoot
Terra agrícola alta prod. (USD/\allowbreak ha) & 4,600 \\
Imóvel urbano (USD/\allowbreak m²) & 600 \\
Aluguel 2 quartos (USD/\allowbreak mês) & 250 \\
\end{longtable}

\begin{quote}
Valores de referência departamental. Localidades menores podem ter
preços 20--40\% abaixo da capital departamental. \#\#\# Saúde
\end{quote}

\textbf{Fonte:} MSPBS /\allowbreak  IPS Paraguay (2024-2026), consolidação
departamental e proxy local

\begin{longtable}[]{@{}
  >{\raggedright\arraybackslash}p{(\linewidth - 2\tabcolsep) * \real{0.3600}}
  >{\raggedright\arraybackslash}p{(\linewidth - 2\tabcolsep) * \real{0.6400}}@{}}
\toprule\noalign{}
\begin{minipage}[b]{\linewidth}\raggedright
Serviço
\end{minipage} & \begin{minipage}[b]{\linewidth}\raggedright
Disponibilidade
\end{minipage} \\
\midrule\noalign{}
\endhead
\bottomrule\noalign{}
\endlastfoot
USF /\allowbreak  Posto de Saúde & sim \\
Hospital Regional & não \\
IPS (seguro social) & não \\
Farmácia & sim \\
Distância ao hospital de referência & \textasciitilde47 km (San Pedro de
Ycuamandiyu) \\
\end{longtable}

\textbf{Principais estabelecimentos:} USF local; referencia hospitalar
em San Pedro de Ycuamandiyu

\textbf{Observação para imigrantes:} Atendimento primario local ou em
raio curto; casos de maior complexidade seguem para San Pedro de
Ycuamandiyu. Cobertura privada continua recomendada para especialidades
e urgências de maior complexidade.
\subsubsection{Dossiê de Campo}\label{dossie-Villa_del_Rosario-dossiuxea-de-campo}
\begin{description}[style=nextline,leftmargin=0.5cm,font=\bfseries]
\item[Coordenadas]  24°25'S 57°07'W (geocodificado via Nominatim).
\end{description}
\paragraph{Solo (SoilGrids 2.0, média ponderada 0--30
cm)}

\begin{longtable}[]{@{}ll@{}}
\toprule\noalign{}
Parâmetro & Valor \\
\midrule\noalign{}
\endhead
\bottomrule\noalign{}
\endlastfoot
pH (H$_2$O) & 5.6 \\
Carbono orgânico (SOC) & 24.67 g/\allowbreak kg \\
Argila & 29.12 \% \\
Areia & 39.07 \% \\
Silte (calc.) & 31.8 \% \\
Densidade aparente & 1.32 g/\allowbreak cm³ \\
\textbf{Aptidão agrícola} & \textbf{Alta} \\
\end{longtable}

Fonte: ISRIC SoilGrids 2.0 via WCS (acesso em 2026-03-21). Coords:
-24.4167°, -57.1167°. Média ponderada camadas 0-5, 5-15, 15-30 cm.
\newpage
\secaoDiagnostico{Yataity del Norte}{\mapaConteudo{-24.8166}{-56.3333}}{Yataity del Norte é um distrito exemplar da estabilidade rural do sul de
San Pedro. Oferece solos férteis e um ambiente social tranquilo, com a
vantagem estratégica de estar sobre um dos eixos rodoviários mais
importantes do país (PY08). É uma localidade ideal para quem busca
produtividade agropecuária com fácil acesso a serviços em cidades
maiores vizinhas.

\subsubsection{Por que sim}

\begin{itemize}
\tightlist
\item
  Solo de alta produtividade e topografia favorável à mecanização.
\item
  Baixa densidade populacional e ambiente seguro.
\item
  Sem restrições da Lei de Fronteira para investidores estrangeiros.
\end{itemize}

\subsubsection{Por que não}

\begin{itemize}
\tightlist
\item
  Dependência de cidades vizinhas para serviços de saúde especializada.
\item
  Suscetibilidade a temporais severos de verão.
\item
  Estradas vicinais de terra exigem atenção em períodos chuvosos.
\end{itemize}

\subsubsection{Recomendação
Técnica}

Altamente indicado para residências rurais produtivas e estabelecimentos
de autossuficiência familiar ou comercial. O GSS de 6.9 reflete o
equilíbrio entre conectividade, recursos e segurança institucional.
Referencia atual de score: GSS 6.9 em 2026-03-05.

\subsubsection{}

\begin{itemize}
\tightlist
\item
  \begin{enumerate}
  \def\labelenumi{\arabic{enumi})}
  \tightlist
  \item
    Dados censitários validados (INE\footnote{Instituto Nacional de Estadística (Instituto Nacional de Estatística), órgão oficial de dados demográficos e censitários.} 2022).
  \end{enumerate}
\item
  \begin{enumerate}
  \def\labelenumi{\arabic{enumi})}
  \setcounter{enumi}{1}
  \tightlist
  \item
    Registro de produtores e cultivos regionais (MAG).
  \end{enumerate}
\item
  \begin{enumerate}
  \def\labelenumi{\arabic{enumi})}
  \setcounter{enumi}{2}
  \tightlist
  \item
    Mapa de conectividade sul de San Pedro (MOPC 2026).
  \end{enumerate}
\item
  \begin{enumerate}
  \def\labelenumi{\arabic{enumi})}
  \setcounter{enumi}{3}
  \tightlist
  \item
    Delimitação departamental e municipal (Portal Geoestadístico).
  \end{enumerate}
\end{itemize}}
\begin{table}[ht!]
\small \begin{tabular}{p{3.0cm}p{0.8cm}p{6.0cm}} \toprule
\textbf{Critério} & \textbf{Nota} & \textbf{Justificativa} \\ \midrule
A - Ameaças Estratégicas & 6.0 & Ponto de conexão estratégica PY08 e divisa departamental; visibilidade moderada \\
B - Risco Social & 7.5 & Baixa densidade populacional rural; risco de agitação social reduzido \\
C - Riscos Naturais & 7.5 & Geologia estável e baixo risco de inundações sistêmicas graves \\
D - Autossuficiência & 7.0 & Bons recursos: solo fértil mecanizável, pecuária e cultivos diversificados \\
E - Institucional & 7.0 & Município consolidado e estável; plena liberdade jurídica para brasileiros \\
\midrule \textbf{GSS FINAL} & \textbf{6.9} & \textbf{Moderadamente Seguro (Excelente opção para refúgios rurais produtivos no Sul de San Pedro).} \\ \bottomrule \end{tabular}
\end{table}
\subsubsection{Vulnerabilidades e
Mitigação}\label{vuln-Yataity_del_Norte-vulnerabilidades-e-mitigauxe7uxe3o}

\begin{itemize}
\tightlist
\item
  \textbf{Dependência Rodoviária:} Concentração logística na PY08
  (Mitigação: residência com acesso por múltiplas rotas rurais
  secundárias).
\item
  \textbf{Isolamento de Saúde:} Distância de centros cirúrgicos
  (Mitigação: treinamento em primeiros socorros avançado e kit médico
  local).
\item
  \textbf{Risco de Tempestades:} Danos em infraestrutura leve
  (Mitigação: telhados reforçados e proteção contra granizo em estufas).
\end{itemize}
\subsubsection{Dossiê de Campo}\label{dossie-Yataity_del_Norte-dossiuxea-de-campo}
\begin{description}[style=nextline,leftmargin=0.5cm,font=\bfseries]
\item[Coordenadas]  24°49'S 56°20'W.
\item[Topografia]  Relevo predominantemente plano a levemente ondulado ("lomadas"). Localizado no extremo sul do departamento de San Pedro, fazendo divisa com Caaguazú pelo arroio Tapiracuái. Geologicamente estável, risco sísmico muito baixo.
\item[Alvos Estratégicos]  Ponto de trânsito e logística sobre a Ruta PY08 (Dr. Blas Garay); Proximidade com o entroncamento da PY03 (20-30 km).
\item[Fallout]  Ventos predominantes do Norte (quentes e úmidos) e Sul (frios e secos). Baixo risco de fallout direto de alvos continentais primários.
\item[População]  10.113 habitantes (Censo 2022 Final). Perfil 83\% rural.
\item[Densidade]  \textasciitilde 20,4 hab/\allowbreak km² (Baixa densidade, oferecendo excelente isolamento em áreas de fazendas e colônias).
\item[Vias de Acesso]  Acesso centralizado na Ruta PY08 (asfaltada). Dependência de estradas internas de terra para acesso às colônias, sujeitas a interrupções em tempestades severas.
\item[Serviços]  Centro de Saúde local para atendimentos primários. Referência hospitalar em San Estanislao (Santaní) ou Coronel Oviedo para serviços de alta complexidade.
\item[Custo de Vida]  Baixo; economia baseada em agricultura familiar diversificada e pecuária.
\item[Preço da Terra]  US\$ 4.000-7.000/\allowbreak ha (terras mecanizadas para soja); US\$ 1.500-3.000/\allowbreak ha (áreas de pastagem ou mata).
\item[Hidrologia]  Presença de arroios locais (ex: Tapiracuái). Baixo risco de inundações graves, mas risco de isolamento de colônias rurais por danos em pontes vicinais e lamaçais em caminhos de terra.
\item[Clima]  Subtropical Úmido. Precipitação anual entre 1.500 e 1.800 mm. Suscetível a temporais de verão com vento e granizo.
\item[Energia]  Rede nacional (Itaipu); interrupções ocasionais em áreas rurais.
\item[Água]  Poços artesianos e nascentes locais; boa disponibilidade subterrânea.
\item[Qualidade do Solo]  Franco-arenoso a argiloso; muito fértil e excelente para mecanização.
\item[Resources Locais]  Produção de mandioca, abacaxi, milho, gergelim, banana e pecuária de corte. Em expansão a cultura da soja mecanizada. Elevado potencial para autossuficiência alimentar.
\item[Segurança]  Zona rural pacata e estável com baixos índices de criminalidade urbana violenta. Exposição a fluxos de trânsito rodoviário nacional. Risco de abigeato em áreas rurais isoladas.
\item[Leis Local: Município consolidado. Livre de Restrição de Fronteira]  Fora da faixa de 50km da Lei 2532/\allowbreak 05, sem impedimentos legais para proprietários brasileiros.
\end{description}
\paragraph{Solo (SoilGrids 2.0, média ponderada 0--30
cm)}

\begin{longtable}[]{@{}ll@{}}
\toprule\noalign{}
Parâmetro & Valor \\
\midrule\noalign{}
\endhead
\bottomrule\noalign{}
\endlastfoot
pH (H$_2$O) & 5.3 \\
Carbono orgânico (SOC) & 19.27 g/\allowbreak kg \\
Argila & 22.63 \% \\
Areia & 58.43 \% \\
Silte (calc.) & 18.9 \% \\
Densidade aparente & 1.27 g/\allowbreak cm³ \\
\textbf{Aptidão agrícola} & \textbf{Média} \\
\end{longtable}

Fonte: ISRIC SoilGrids 2.0 via WCS (acesso em 2026-03-21). Coords:
-24.8167°, -56.3333°. Média ponderada camadas 0-5, 5-15, 15-30 cm.

\begin{itemize}
\tightlist
\item
  \textbf{Preço Terra Rural:} US\$ 1.500-7.000/\allowbreak ha.
\end{itemize}

\subsection{10. Analise de Riscos}

\begin{itemize}
\tightlist
\item
  \textbf{Cenario curto prazo:} Manutenção da estabilidade rural e
  valorização de áreas de expansão da soja.
\item
  \textbf{Cenario medio prazo:} Impacto positivo da melhoria de pontes e
  acessos vicinais.
\end{itemize}

\subsection{11. Redacao Final}

\subsubsection{Diagnostico Integrado}

Yataity del Norte é um distrito exemplar da estabilidade rural do sul de
San Pedro. Oferece solos férteis e um ambiente social tranquilo, com a
vantagem estratégica de estar sobre um dos eixos rodoviários mais
importantes do país (PY08). É uma localidade ideal para quem busca
produtividade agropecuária com fácil acesso a serviços em cidades
maiores vizinhas.

\subsubsection{Por que sim}

\begin{itemize}
\tightlist
\item
  Solo de alta produtividade e topografia favorável à mecanização.
\item
  Baixa densidade populacional e ambiente seguro.
\item
  Sem restrições da Lei de Fronteira para investidores brasileiros.
\end{itemize}

\subsubsection{Por que nao}

\begin{itemize}
\tightlist
\item
  Dependência de cidades vizinhas para serviços de saúde especializada.
\item
  Suscetibilidade a temporais severos de verão.
\item
  Estradas vicinais de terra exigem atenção em períodos chuvosos.
\end{itemize}

\subsubsection{Recomendacao Tecnica}

Altamente indicado para residências rurais produtivas e estabelecimentos
de autossuficiência familiar ou comercial. O GSS de 6.9 reflete o
equilíbrio entre conectividade, recursos e segurança institucional.
Referencia atual de score: GSS 6.9 em 2026-03-05.

\subsection{13.}

\begin{itemize}
\tightlist
\item
  \begin{enumerate}
  \def\labelenumi{\arabic{enumi})}
  \tightlist
  \item
    Dados censitários validados (INE\footnote{Instituto Nacional de Estadística (Instituto Nacional de Estatística), órgão oficial de dados demográficos e censitários.} 2022).
  \end{enumerate}
\item
  \begin{enumerate}
  \def\labelenumi{\arabic{enumi})}
  \setcounter{enumi}{1}
  \tightlist
  \item
    Registro de produtores e cultivos regionais (MAG).
  \end{enumerate}
\item
  \begin{enumerate}
  \def\labelenumi{\arabic{enumi})}
  \setcounter{enumi}{2}
  \tightlist
  \item
    Mapa de conectividade sul de San Pedro (MOPC 2026).
  \end{enumerate}
\item
  \begin{enumerate}
  \def\labelenumi{\arabic{enumi})}
  \setcounter{enumi}{3}
  \tightlist
  \item
    Delimitação departamental e municipal (Portal Geoestadístico).
  \end{enumerate}
\end{itemize}

\subsubsection{Combustível}

\textbf{Referência:} PETROPAR /\allowbreak  postos locais (2024) \textbf{Tipo de
localidade:} Interior

\begin{longtable}[]{@{}lll@{}}
\toprule\noalign{}
Combustível & USD/\allowbreak litro & Gs/\allowbreak litro (aprox.) \\
\midrule\noalign{}
\endhead
\bottomrule\noalign{}
\endlastfoot
Gasolina 93 oct & 1.00 & 7,400 \\
Gasolina 97 oct (premium) & 1.10 & 8,140 \\
Diesel & 0.93 & 6,882 \\
\end{longtable}

\begin{quote}
Preços podem variar ±5\% conforme posto e sazonalidade. Chaco e interior
remoto apresentam maior variação.
\end{quote}

\subsubsection{Cobertura Celular}

\textbf{Fonte:} CONATEL PY /\allowbreak  operadoras (2024)

\begin{longtable}[]{@{}ll@{}}
\toprule\noalign{}
Parâmetro & Valor \\
\midrule\noalign{}
\endhead
\bottomrule\noalign{}
\endlastfoot
Cobertura 4G população (dept.) & 80\% \\
Cobertura 4G área rural & 58\% \\
Melhor operadora & Tigo \\
Qualidade rural & moderada \\
\end{longtable}

\begin{quote}
Para áreas rurais fora do núcleo urbano, recomenda-se chip Tigo como
principal e Personal como backup.
\end{quote}

\subsubsection{Internet}

\textbf{Fonte:} CONATEL /\allowbreak  Speedtest Ookla (2024)

\begin{longtable}[]{@{}ll@{}}
\toprule\noalign{}
Parâmetro & Valor \\
\midrule\noalign{}
\endhead
\bottomrule\noalign{}
\endlastfoot
Velocidade média download & 35 Mbps \\
Domicílios com internet (dept.) & 45\% \\
Tecnologia predominante & rádio \\
Opção rural & Starlink disponível (\textasciitilde USD 44/\allowbreak mês) \\
\end{longtable}

\subsubsection{Mercado Imobiliário e Terra
Rural}

\textbf{Fonte:} INDERT /\allowbreak  Clasificados.com.py (2024)

\begin{longtable}[]{@{}ll@{}}
\toprule\noalign{}
Tipo & Referência \\
\midrule\noalign{}
\endhead
\bottomrule\noalign{}
\endlastfoot
Terra agrícola alta prod. (USD/\allowbreak ha) & 4,600 \\
Imóvel urbano (USD/\allowbreak m²) & 600 \\
Aluguel 2 quartos (USD/\allowbreak mês) & 250 \\
\end{longtable}

\begin{quote}
Valores de referência departamental. Localidades menores podem ter
preços 20--40\% abaixo da capital departamental. \#\#\# Saúde
\end{quote}

\textbf{Fonte:} MSPBS /\allowbreak  IPS Paraguay (2024-2026), consolidação
departamental e proxy local

\begin{longtable}[]{@{}
  >{\raggedright\arraybackslash}p{(\linewidth - 2\tabcolsep) * \real{0.3600}}
  >{\raggedright\arraybackslash}p{(\linewidth - 2\tabcolsep) * \real{0.6400}}@{}}
\toprule\noalign{}
\begin{minipage}[b]{\linewidth}\raggedright
Serviço
\end{minipage} & \begin{minipage}[b]{\linewidth}\raggedright
Disponibilidade
\end{minipage} \\
\midrule\noalign{}
\endhead
\bottomrule\noalign{}
\endlastfoot
USF /\allowbreak  Posto de Saúde & sim \\
Hospital Regional & sim \\
IPS (seguro social) & não \\
Farmácia & sim \\
Distância ao hospital de referência & \textasciitilde140 km (San Pedro
de Ycuamandiyu) \\
\end{longtable}

\textbf{Principais estabelecimentos:} USF local; referencia hospitalar
em San Pedro de Ycuamandiyu

\textbf{Observação para imigrantes:} Atendimento primario local ou em
raio curto; casos de maior complexidade seguem para San Pedro de
Ycuamandiyu. Cobertura privada continua recomendada para especialidades
e urgências de maior complexidade.
\newpage
\secaoDiagnostico{Yrybucua}{\mapaConteudo{-24.4833}{-56.1833}}{Yrybucuá é um ponto estratégico de conexão no coração produtivo do
Paraguai. Oferece uma combinação sólida de segurança física,
estabilidade institucional e abundância de recursos agropecuários. Sua
localização como zona de transição entre três departamentos importantes
garante um fluxo constante de suprimentos e fácil acesso a eixos
rodoviários vitais.

\subsubsection{Por que sim}

\begin{itemize}
\tightlist
\item
  Excelente conectividade rodoviária para Assunção e para a fronteira
  brasileira.
\item
  Solo mecanizável com alto potencial de resposta produtiva.
\item
  Sem restrições da Lei de Fronteira para estrangeiros em terras
  privadas.
\end{itemize}

\subsubsection{Por que não}

\begin{itemize}
\tightlist
\item
  Solo arenoso exige investimento em manejo de conservação.
\item
  Dependência de serviços hospitalares em Santaní ou San Pedro.
\item
  Risco regional de crimes contra o patrimônio rural (abigeato).
\end{itemize}

\subsubsection{Recomendação
Técnica}

Altamente indicado para estabelecimentos rurais produtivos e residências
que busquem um refúgio com logística facilitada. O GSS de 6.8 reflete o
equilíbrio entre a riqueza de recursos e a segurança institucional da
região. Referencia atual de score: GSS 6.8 em 2026-03-05.

\subsubsection{}

\begin{itemize}
\tightlist
\item
  \begin{enumerate}
  \def\labelenumi{\arabic{enumi})}
  \tightlist
  \item
    Dados censitários (INE\footnote{Instituto Nacional de Estadística (Instituto Nacional de Estatística), órgão oficial de dados demográficos e censitários.} 2022).
  \end{enumerate}
\item
  \begin{enumerate}
  \def\labelenumi{\arabic{enumi})}
  \setcounter{enumi}{1}
  \tightlist
  \item
    Relatórios de expansão da fronteira agrícola mecanizada (MAG).
  \end{enumerate}
\item
  \begin{enumerate}
  \def\labelenumi{\arabic{enumi})}
  \setcounter{enumi}{2}
  \tightlist
  \item
    Mapa de solos e aptidão do Centro-Norte (MAG/\allowbreak Productiva).
  \end{enumerate}
\item
  \begin{enumerate}
  \def\labelenumi{\arabic{enumi})}
  \setcounter{enumi}{3}
  \tightlist
  \item
    Delimitação de rotas nacionais e vicinais (MOPC 2026).
  \end{enumerate}
\end{itemize}}
\begin{table}[ht!]
\small \begin{tabular}{p{3.0cm}p{0.8cm}p{6.0cm}} \toprule
\textbf{Critério} & \textbf{Nota} & \textbf{Justificativa} \\ \midrule
A - Ameaças Estratégicas & 6.0 & Posição de hub de transição entre três departamentos; visibilidade estratégica moderada \\
B - Risco Social & 7.0 & Densidade rural moderada; risco social urbano e de agitação controlado \\
C - Riscos Naturais & 7.5 & Geologia muito estável e baixo risco de grandes inundações \\
D - Autossuficiência & 7.5 & Recursos agrícolas em expansão acelerada; abundância de grãos e proteína animal \\
E - Institucional & 6.5 & Ambiente institucional funcional; total liberdade jurídica para brasileiros \\
\midrule \textbf{GSS FINAL} & \textbf{6.8} & \textbf{Moderadamente Seguro (Excelente para perfis focados em agronegócio e autossuficiência produtiva).} \\ \bottomrule \end{tabular}
\end{table}
\subsubsection{Vulnerabilidades e
Mitigação}\label{vuln-Yrybucua-vulnerabilidades-e-mitigauxe7uxe3o}

\begin{itemize}
\tightlist
\item
  \textbf{Erosão do Solo:} Solos arenosos vulneráveis (Mitigação:
  cobertura de solo permanente e terraceamento).
\item
  \textbf{Dependência Logística:} Dependência de Santaní para serviços
  avançados (Mitigação: manutenção de estoque médico e técnico local).
\item
  \textbf{Insegurança Rural:} Risco de roubos em fazendas isoladas
  (Mitigação: sistemas de vigilância eletrônica e rádio comunitário de
  segurança).
\end{itemize}
\subsubsection{Dossiê de Campo}\label{dossie-Yrybucua-dossiuxea-de-campo}
\begin{description}[style=nextline,leftmargin=0.5cm,font=\bfseries]
\item[Coordenadas]  24°29'S 56°11'W.
\item[Topografia]  Relevo predominantemente plano a ondulado. Localizado em um ponto estratégico de transição entre San Pedro, Caaguazú e Canindeyú. Geologicamente estável, risco sísmico praticamente nulo.
\item[Alvos Estratégicos]  Hub de transição rodoviária próximo ao entroncamento das Rutas PY03 e PY08; Áreas de expansão acelerada de soja mecanizada.
\item[Fallout]  Ventos predominantes do Norte/\allowbreak Nordeste (quentes e úmidos) e Sul (frentes frias). Baixo risco de fallout direto de alvos continentais primários.
\item[População]  12.271 habitantes (Censo 2022 Final).
\item[Densidade]  \textasciitilde 24,1 hab/\allowbreak km² (Densidade moderada-baixa, com população dispersa em áreas de expansão agrícola).
\item[Vias de Acesso]  Excelente conectividade regional pelas Rutas PY03 e PY08. Dependência de estradas vicinais de terra para o interior do distrito, que exigem manutenção contra erosão em solos arenosos.
\item[Serviços]  Unidades de Saúde Familiar (USF) para atenção primária. Dependência de San Estanislao (Santaní) a 30-40 km para serviços de saúde especializados e alta complexidade.
\item[Custo de Vida]  Baixo (US\$ 550-800/\allowbreak mês para padrão médio); alimentos básicos produzidos localmente favorecem a economia doméstica.
\item[Preço da Terra]  US\$ 4.500-6.500/\allowbreak ha (terras mecanizadas para soja); US\$ 1.200-2.500/\allowbreak ha (áreas de pecuária ou mistas).
\item[Hidrologia]  Baixo risco de inundações sistêmicas. Arroios locais podem transbordar e isolar colônias rurais temporariamente durante eventos de chuva extrema (precipitação média 1.500-1.600 mm).
\item[Clima]  Subtropical Úmido. Chuvas bem distribuídas; verões quentes com tempestades súbitas.
\item[Energia]  Rede nacional (Itaipu); sujeita a interrupções em áreas rurais devido à expansão da fronteira agrícola.
\item[Água]  Boa disponibilidade subterrânea via poços artesianos e nascentes perenes.
\item[Qualidade do Solo]  Predomínio de solos arenosos; em áreas mecanizadas, transição para Latossolos (Terra Vermelha) de alta fertilidade. Exige plantio direto rigoroso.
\item[Resources Locais]  Grande produção de soja, milho, sésamo, mandioca e pecuária de corte (Brachiaria). Elevado potencial para autossuficiência alimentar.
\item[Segurança]  Zona rural considerada estável e mais tranquila que o norte do departamento. Riscos focados em conflitos de terra esporádicos e crimes contra o patrimônio em fazendas (abigeato). Taxa de homicídios regional de \textasciitilde 7,9/\allowbreak 100k.
\item[Leis Local: Município consolidado pela Lei Orgânica 3966/\allowbreak 10. Livre de Restrição de Fronteira]  Fora da faixa de 50km da Lei 2532/\allowbreak 05, permitindo plena titularidade para investidores brasileiros.
\end{description}
\paragraph{Solo (SoilGrids 2.0, média ponderada 0--30
cm)}

\begin{longtable}[]{@{}ll@{}}
\toprule\noalign{}
Parâmetro & Valor \\
\midrule\noalign{}
\endhead
\bottomrule\noalign{}
\endlastfoot
pH (H$_2$O) & 5.3 \\
Carbono orgânico (SOC) & 22.05 g/\allowbreak kg \\
Argila & 23.23 \% \\
Areia & 57.43 \% \\
Silte (calc.) & 19.3 \% \\
Densidade aparente & 1.23 g/\allowbreak cm³ \\
\textbf{Aptidão agrícola} & \textbf{Média} \\
\end{longtable}

Fonte: ISRIC SoilGrids 2.0 via WCS (acesso em 2026-03-21). Coords:
-24.4833°, -56.1833°. Média ponderada camadas 0-5, 5-15, 15-30 cm.
